% La presse écrite et le journal télévisé — Allemand 1ère A — Cours signé OnBuch+
% Programme officiel MINESEC. Compiler avec : tectonic ALA16-la-presse-ecrite-et-le-journal-televise.tex
\newcommand{\DOCMATIERE}{Allemand}
\newcommand{\DOCNIVEAU}{1ère A}
\newcommand{\DOCMODULE}{Chapitre 4 : Média et communication}
\newcommand{\DOCLECON}{ALA16}
\newcommand{\DOCTITRE}{La presse écrite et le journal télévisé}
% OnBuch+ — préambule « cours signés » ENRICHI (compilé avec tectonic / XeLaTeX)
% Système visuel « Pop » d'OnBuch+ : fond crème (jamais blanc), bordures encre
% épaisses, ombres « dures » décalées, titres Archivo Black en capitales.
% Version MATHÉMATIQUES : graphes (pgfplots/tikz), boîtes pédagogiques
% (définition, propriété, méthode, attention, le savais-tu, exercice résolu,
% exercice, TP, bilan), page de garde et signature OnBuch+ ancrée en bas.
%
% Macros attendues dans main.tex AVANT \input{preamble} :
%   \DOCMATIERE  \DOCNIVEAU  \DOCMODULE  \DOCLECON (numéro)  \DOCTITRE
\documentclass[11pt]{article}

\usepackage{fontspec}
\usepackage[ngerman,french]{babel} % français = langue principale ; l'allemand via \all{..} / textebox
\usepackage[a4paper,margin=2cm,top=3.4cm,bottom=2.8cm,headheight=1.4cm,footskip=1.3cm]{geometry}
\usepackage{amsmath,amssymb,amsfonts}
\usepackage{mathtools} % \xmapsto, \coloneqq... souvent utilisés par les modèles
\usepackage{cancel} % simplifications barrées (bilans de forces)
% \abs{z} (module, valeur absolue) et \norm{u} (norme), utilisés par les modèles
\providecommand{\abs}{}\let\abs\relax\DeclarePairedDelimiter{\abs}{\lvert}{\rvert}
\providecommand{\norm}{}\let\norm\relax\DeclarePairedDelimiter{\norm}{\lVert}{\rVert}
\usepackage[table]{xcolor} % [table] : fournit aussi \rowcolors (alternance de lignes)
\usepackage{pagecolor}
\usepackage{tikz}
\usetikzlibrary{arrows.meta,positioning,calc,shapes.geometric,shapes.misc,shapes.arrows,decorations.pathmorphing,decorations.pathreplacing,decorations.markings,patterns,shadows,mindmap,trees,fit,backgrounds,matrix,intersections,angles,quotes}
\usepackage{pgfplots}
\usepackage[version=4]{mhchem} % équations nucléaires \ce{^{235}_{92}U}
\usepackage[european,siunitx]{circuitikz} % schémas électriques
\pgfplotsset{compat=1.18}
\usepgfplotslibrary{fillbetween,groupplots} % aires entre courbes (calcul intégral)
\usepackage{enumitem}
\usepackage{fancyhdr}
% [most] charge toutes les bibliothèques tcolorbox (listings, minted, raster,
% documentation...) et gonfle inutilement la mémoire TeX sur un document long
% (~300 boîtes) : on ne charge que ce qui est réellement utilisé.
\usepackage[skins,breakable]{tcolorbox}
\usepackage{graphicx}
\usepackage{colortbl}
\usepackage{booktabs}
\usepackage{tabularx}
\usepackage{diagbox} % case d'en-tête diagonale des tableaux à double entrée
% Colonnes extensibles centrée / gauche / droite (C, L, R), souvent utilisées par les modèles
\newcolumntype{C}{>{\centering\arraybackslash}X}
\newcolumntype{L}{>{\raggedright\arraybackslash}X}
\newcolumntype{R}{>{\raggedleft\arraybackslash}X}
\usepackage{array}
\usepackage{multicol}
\usepackage{multirow}
\usepackage{siunitx}
% parse-numbers=false : accepte aussi \SI{2,5 \times 10^{-3}}{...}
\sisetup{locale=FR,output-decimal-marker={,},group-minimum-digits=4,parse-numbers=false}
\DeclareSIUnit\ohm{\ensuremath{\symup{\Omega}}} % U+2126 absent de la police
\DeclareSIUnit\division{div} % réglages d'oscilloscope (ms/div, V/div)
\DeclareSIUnit\tr{tr} % tours (vitesse de rotation, ex. \SI{1500}{\tr\per\minute})
\DeclareSIUnit\molar{M} % concentration molaire (ex. \SI{3}{\molar}, \SI{1}{\micro\molar})
\DeclareSIUnit\an{an} % année (le modèle confond parfois avec \year, un compteur TeX) — ex. \SI{5}{\kilo\meter\per\an}
\DeclareSIUnit\FCFA{FCFA} % franc CFA (ex. \SI{500}{\FCFA\per\kilo\gram})
\DeclareSIUnit\degreeBrix{\textdegree Brix} % teneur en sucre d'un sirop/jus (réfractométrie)
\DeclareSIUnit\month{mois} % le modèle utilise parfois \month, un compteur TeX, comme unité de durée
\usepackage{qrcode}
% \degree : commande fréquemment utilisée par les modèles pour un angle ou une
% température en dehors de \SI{}{} (non fournie par les paquets chargés).
\providecommand{\degree}{^{\circ}}
% \qed (fin de démonstration), utilisé par les modèles sans amsthm
\providecommand{\qed}{\hfill$\square$}
% \barre : double barre des valeurs interdites dans un tableau de variations (tkz-tab)
\providecommand{\barre}{\ensuremath{\|}}
% environnement proof (amsthm non chargé) utilisé par les modèles
\ifdefined\proof\else\newenvironment{proof}[1][Démonstration]{\par\noindent\textit{#1.}\ }{\qed\par}\fi
\usepackage{lastpage}
\usepackage{hyperref}
\hypersetup{hidelinks}

% ============================================================
% Palette « Pop » OnBuch+ — mêmes hex que la classe P (plus_ui.dart)
% ============================================================
\definecolor{poporange}{HTML}{F59321}
\definecolor{popdark}{HTML}{E07A0C}
\definecolor{popcream}{HTML}{FFF6E8}
\definecolor{popink}{HTML}{1C1714}
\definecolor{poppurple}{HTML}{7A5AE0}
\definecolor{popgreen}{HTML}{2E9E6A}
\definecolor{poppink}{HTML}{E0457B}
\definecolor{popblue}{HTML}{2F7DE1}
\definecolor{popgold}{HTML}{C98A12}
\definecolor{popmuted}{HTML}{8A7F76}
\definecolor{popline}{HTML}{EADFD2}
\definecolor{popgray}{HTML}{8A7F76}
\colorlet{popgrey}{popgray}
% Alias tolérants (le modèle nomme parfois les couleurs différemment)
\colorlet{popwhite}{white}
\colorlet{popblack}{popink}
\colorlet{popred}{poppink}
\colorlet{popyellow}{popgold}
% Teintes claires (fonds de tableaux/encadrés) — le modèle nomme parfois
% les couleurs avec un suffixe L (ex. popblueL) sans les avoir définies.
\colorlet{poporangeL}{poporange!20}
\colorlet{popdarkL}{popdark!20}
\colorlet{poppurpleL}{poppurple!20}
\colorlet{popgreenL}{popgreen!20}
\colorlet{poppinkL}{poppink!20}
\colorlet{popblueL}{popblue!20}
\colorlet{popgoldL}{popgold!20}
\colorlet{popredL}{popred!20}
\colorlet{popgrayL}{popgray!20}
% Teintes claires pour fonds de tableaux / zones
\colorlet{popblueL}{popblue!10!white}
\colorlet{poporangeL}{poporange!14!white}
\colorlet{popgreenL}{popgreen!12!white}
\colorlet{poppurpleL}{poppurple!12!white}
\colorlet{poppinkL}{poppink!12!white}
\colorlet{popgoldL}{popgold!14!white}

\pagecolor{popcream}

% ============================================================
% Typographie
% ============================================================
\defaultfontfeatures{Ligatures=TeX}
\setmainfont{PlusJakartaSans-Medium.ttf}[Path=../../../fonts/,BoldFont=PlusJakartaSans-Bold.ttf]
% Symboles, API et emojis absents de la police principale : repli sur DejaVu Sans (caractères actifs)
\newfontface\symfont{DejaVuSans.ttf}[Path=../../../fonts/]
\makeatletter
\newcommand\symactive[1]{\catcode#1=13 \begingroup\lccode`\~=#1 \lowercase{\endgroup\def~}{{\symfont\char#1}}}
\newcommand\symblank[1]{\catcode#1=13 \begingroup\lccode`\~=#1 \lowercase{\endgroup\def~}{}}
\newcommand\symhyph[1]{\catcode#1=13 \begingroup\lccode`\~=#1 \lowercase{\endgroup\def~}{\mbox{-}}}
\newcount\symc
\newcommand\symrange[2]{\symc=#1 \@whilenum\symc<#2 \do{\expandafter\symactive\expandafter{\the\symc}\advance\symc by 1 }}
\newcommand\symblankrange[2]{\symc=#1 \@whilenum\symc<#2 \do{\expandafter\symblank\expandafter{\the\symc}\advance\symc by 1 }}
\makeatother
\symrange{"0250}{"02FF}\symactive{"2713}\symactive{"2714}\symactive{"2717}\symactive{"2718}\symactive{"2610}\symactive{"2611}\symactive{"2612}
\symhyph{"2011}\symhyph{"2E1A}\symblankrange{"1F300}{"1FAFF}\symblank{"FE0F}
% Maths en sans-serif (Fira Math) avec chiffres et lettres droites de Plus
% Jakarta, pour que les formules se fondent dans le texte.
\usepackage[math-style=ISO,bold-style=ISO]{unicode-math}
\setmathfont{FiraMath-Regular.otf}[Path=../../../fonts/]
\setmathfont{PlusJakartaSans-Medium.ttf}[Path=../../../fonts/,range={up/num,up/{Latin,latin},\mathup}]
\usepackage{icomma} % virgule décimale sans espace en mode math
% Caractères mathématiques Unicode absents des polices de texte (Plus Jakarta,
% Archivo Black) : rendus par leur équivalent mathématique, en texte comme en
% formule — sinon ils s'affichent en carré vide (ex. « Arithmétique dans ℤ »).
\usepackage{newunicodechar}
\newunicodechar{ℝ}{\ensuremath{\mathbb{R}}}
\newunicodechar{ℤ}{\ensuremath{\mathbb{Z}}}
\newunicodechar{ℂ}{\ensuremath{\mathbb{C}}}
\newunicodechar{ℕ}{\ensuremath{\mathbb{N}}}
\newunicodechar{ℚ}{\ensuremath{\mathbb{Q}}}
\newunicodechar{𝔻}{\ensuremath{\mathbb{D}}}
\newunicodechar{α}{\ensuremath{\alpha}}
\newunicodechar{β}{\ensuremath{\beta}}
\newunicodechar{γ}{\ensuremath{\gamma}}
\newunicodechar{δ}{\ensuremath{\delta}}
\newunicodechar{ε}{\ensuremath{\varepsilon}}
\newunicodechar{θ}{\ensuremath{\theta}}
\newunicodechar{λ}{\ensuremath{\lambda}}
\newunicodechar{μ}{\ensuremath{\mu}}
\newunicodechar{σ}{\ensuremath{\sigma}}
\newunicodechar{τ}{\ensuremath{\tau}}
\newunicodechar{φ}{\ensuremath{\varphi}}
\newunicodechar{ψ}{\ensuremath{\psi}}
\newunicodechar{ω}{\ensuremath{\omega}}
\newunicodechar{Σ}{\ensuremath{\Sigma}}
% majuscules produites par \MakeUppercase dans les titres (α-aminés -> Α)
\newunicodechar{Α}{\ensuremath{\alpha}}
\newunicodechar{Β}{\ensuremath{\beta}}
\newunicodechar{Γ}{\ensuremath{\Gamma}}
\newunicodechar{Δ}{\ensuremath{\Delta}}
\newunicodechar{Ε}{\ensuremath{\varepsilon}}
\newunicodechar{Θ}{\ensuremath{\Theta}}
\newunicodechar{Λ}{\ensuremath{\Lambda}}
\newunicodechar{Μ}{\ensuremath{\mu}}
\newunicodechar{Τ}{\ensuremath{\tau}}
\newunicodechar{Φ}{\ensuremath{\Phi}}
\newunicodechar{Ψ}{\ensuremath{\Psi}}
\newunicodechar{Ω}{\ensuremath{\Omega}}
\newunicodechar{↦}{\ensuremath{\mapsto}}
\newunicodechar{∈}{\ensuremath{\in}}
\newunicodechar{∉}{\ensuremath{\notin}}
\newunicodechar{∧}{\ensuremath{\wedge}}
\newunicodechar{⇔}{\ensuremath{\Leftrightarrow}}
\newunicodechar{⟺}{\ensuremath{\Longleftrightarrow}}
\newunicodechar{⇒}{\ensuremath{\Rightarrow}}
\newunicodechar{⟹}{\ensuremath{\Longrightarrow}}
\newunicodechar{∘}{\ensuremath{\circ}}
\newunicodechar{∪}{\ensuremath{\cup}}
\newunicodechar{∩}{\ensuremath{\cap}}
\newunicodechar{≡}{\ensuremath{\equiv}}
\newunicodechar{⊂}{\ensuremath{\subset}}
\newunicodechar{⊥}{\ensuremath{\perp}}
\newunicodechar{∃}{\ensuremath{\exists}}
\newunicodechar{∀}{\ensuremath{\forall}}
\newunicodechar{∛}{\ensuremath{\sqrt[3]{\,}}}
\newunicodechar{∜}{\ensuremath{\sqrt[4]{\,}}}
\newunicodechar{⃗}{\ensuremath{{}^{\to}}}
\newunicodechar{ⁿ}{\ifmmode^{n}\else\textsuperscript{\ensuremath{n}}\fi}
\newunicodechar{ˣ}{\ifmmode^{x}\else\textsuperscript{\ensuremath{x}}\fi}
\newunicodechar{ᵇ}{\ifmmode^{b}\else\textsuperscript{\ensuremath{b}}\fi}
\newunicodechar{ʳ}{\ifmmode^{r}\else\textsuperscript{\ensuremath{r}}\fi}
\newunicodechar{ᵘ}{\ifmmode^{u}\else\textsuperscript{\ensuremath{u}}\fi}
\newunicodechar{ʸ}{\ifmmode^{y}\else\textsuperscript{\ensuremath{y}}\fi}
\newunicodechar{ᵃ}{\ifmmode^{a}\else\textsuperscript{\ensuremath{a}}\fi}
\newunicodechar{ᵉ}{\ifmmode^{e}\else\textsuperscript{\ensuremath{e}}\fi}
\newunicodechar{ᵏ}{\ifmmode^{k}\else\textsuperscript{\ensuremath{k}}\fi}
\newunicodechar{ᵖ}{\ifmmode^{p}\else\textsuperscript{\ensuremath{p}}\fi}
\newunicodechar{ᴺ}{\ifmmode^{N}\else\textsuperscript{\ensuremath{N}}\fi}
\newunicodechar{ᶜ}{\ifmmode^{c}\else\textsuperscript{\ensuremath{c}}\fi}
\newunicodechar{ʰ}{\ifmmode^{h}\else\textsuperscript{\ensuremath{h}}\fi}
\newunicodechar{ᵠ}{\ifmmode^{q}\else\textsuperscript{\ensuremath{q}}\fi}
\newunicodechar{ₙ}{\ifmmode_{n}\else\textsubscript{\ensuremath{n}}\fi}
\newunicodechar{ₐ}{\ifmmode_{a}\else\textsubscript{\ensuremath{a}}\fi}
\newunicodechar{ᵢ}{\ifmmode_{i}\else\textsubscript{\ensuremath{i}}\fi}
\newunicodechar{ᵣ}{\ifmmode_{r}\else\textsubscript{\ensuremath{r}}\fi}
\newunicodechar{ₖ}{\ifmmode_{k}\else\textsubscript{\ensuremath{k}}\fi}
\newunicodechar{ᵦ}{\ifmmode_{\beta}\else\textsubscript{\ensuremath{\beta}}\fi}
\newunicodechar{ₓ}{\ifmmode_{x}\else\textsubscript{\ensuremath{x}}\fi}
\newfontfamily\popdisplay{ArchivoBlack-Regular.ttf}[Path=../../../fonts/]
\newfontfamily\poplogo{SpaceGrotesk-Bold.ttf}[Path=../../../fonts/]
\newfontfamily\popsemi{PlusJakartaSans-SemiBold.ttf}[Path=../../../fonts/]
\newfontfamily\popxbold{PlusJakartaSans-ExtraBold.ttf}[Path=../../../fonts/]
\newfontfamily\popxboldls{PlusJakartaSans-ExtraBold.ttf}[Path=../../../fonts/,LetterSpace=3.0]

\setlist[enumerate]{leftmargin=1.7em,itemsep=5pt,topsep=4pt}
\setlist[itemize]{leftmargin=1.7em,itemsep=4pt,topsep=4pt,label={\color{poporange}\textbullet}}
\setlength{\parindent}{0pt}
\setlength{\parskip}{5pt}
\renewcommand{\arraystretch}{1.25}

% pgfplots : style Pop par défaut
\pgfplotsset{
  every axis/.append style={axis line style={popink,line width=1pt},tick style={popink},
    grid style={popline,line width=0.5pt},label style={font=\small},tick label style={font=\footnotesize}},
  popcurve/.style={poporange,line width=1.8pt,no markers,smooth},
}
% Même style disponible dans un \draw TikZ simple (hors pgfplots)
\tikzset{popcurve/.style={poporange,line width=1.8pt}}
\tikzset{popgrid/.style={popline,very thin}}
\colorlet{popcurve}{poporange} % le nom du style est parfois utilisé comme couleur

% ============================================================
% Pill / squiggle
% ============================================================
\newcommand{\poppill}[3]{%
  \tikz[baseline=-0.55ex]\node[draw=popink,line width=1.2pt,rounded corners=2.6pt,%
    fill=#2,inner xsep=6.5pt,inner ysep=3pt]{%
    \popxboldls\fontsize{7}{7}\selectfont\color{#3}\MakeUppercase{#1}};%
}
\newcommand{\popsquiggle}[1]{%
  \tikz[baseline=-0.4ex]{
    \draw[#1,line width=2.6pt,line cap=round]
      plot [smooth] coordinates {(0,0.09) (0.34,0.01) (0.68,0.21) (1.02,0.01) (1.36,0.10)};
  }%
}
% Mot-clé surligné (marqueur orange sous le texte)
\newcommand{\cle}[1]{{\popxbold\color{popdark}#1}}

% ============================================================
% En-tête / pied
% ============================================================
\pagestyle{fancy}
\fancyhf{}
\renewcommand{\headrulewidth}{0pt}
\renewcommand{\footrulewidth}{0pt}
\lhead{\raisebox{-0.15ex}{\poplogo\fontsize{13}{13}\selectfont\color{popink}OnBuch\color{poporange}+}}
\rhead{\poppill{\DOCMATIERE\ \textbullet\ \DOCNIVEAU\ \textbullet\ Leçon \DOCLECON}{white}{popink}}
\lfoot{\popxbold\fontsize{7.6}{7.6}\selectfont\color{popmuted}\MakeUppercase{Cours signé OnBuch+}\ \textbullet\ {\popsemi\DOCTITRE}}
\rfoot{\tikz[baseline=-0.6ex]\node[rounded corners=8.5pt,fill=popink,minimum height=17pt,minimum width=17pt,inner xsep=5pt,inner sep=0]{\popxbold\fontsize{8}{8}\selectfont\color{popcream}\thepage\,/\,\pageref*{LastPage}};}

% ============================================================
% Titres
% ============================================================
\newcommand{\chaptitre}[2]{%
  \begin{tcolorbox}[enhanced,colback=poporange,colframe=popink,boxrule=2.6pt,arc=11pt,
      drop shadow=popink,left=15pt,right=15pt,top=12pt,bottom=11pt,before skip=3pt,after skip=18pt]
    \poppill{#2}{popink}{popcream}\par\vspace{9pt}
    {\raggedright\hyphenpenalty=10000\exhyphenpenalty=10000\popdisplay\fontsize{22}{24}\selectfont\color{popcream}\MakeUppercase{#1}\par}
  \end{tcolorbox}
}
% Section numérotée : pastille orange avec le numéro + titre Archivo
\newcounter{popsec}
\newcommand{\coursec}[1]{%
  \stepcounter{popsec}\setcounter{popsub}{0}%
  \par\vspace{16pt}\noindent
  \begin{minipage}{\linewidth}
  \tikz[baseline=-0.7ex]\node[draw=popink,line width=1.6pt,fill=poporange,rounded corners=4pt,minimum size=22pt,inner sep=0]{\popdisplay\fontsize{12}{12}\selectfont\color{popcream}\thepopsec};%
  \hspace{8pt}{\popdisplay\fontsize{15}{17}\selectfont\color{popink}\MakeUppercase{#1}}\par
  \vspace{4pt}\noindent{\color{poporange}\rule{36pt}{3.6pt}}
  \end{minipage}\par\nopagebreak\vspace{8pt}
}
\newcounter{popsub}
\newcommand{\courssub}[1]{%
  \stepcounter{popsub}%
  \par\vspace{9pt}\noindent{\popxbold\fontsize{11.5}{13}\selectfont\color{popink}{\color{poporange}\thepopsec.\thepopsub}\hspace{6pt}#1}\par\nopagebreak\vspace{4pt}
}

% ============================================================
% Boîtes pédagogiques — même recette : fond blanc, bordure épaisse colorée,
% ombre dure encre, badge plein. Titre optionnel [#1].
% ============================================================
\tcbset{popbase/.style={breakable,enhanced,colback=white,boxrule=2.3pt,arc=9pt,
  shadow={3pt}{-3pt}{0pt}{popink},left=14pt,right=14pt,top=11pt,bottom=11pt,before skip=11pt,after skip=15pt,
  fontupper=\popsemi\fontsize{10.6}{14.5}\selectfont\color{popink},
  fontlower=\popsemi\fontsize{10.6}{14.5}\selectfont\color{popink}}}
% Coche dessinée (le glyphe ✓ n'existe pas dans les polices du document)
\newcommand{\popcheck}[1]{\tikz[baseline=-0.5ex]\draw[#1,line width=1.6pt,line cap=round,line join=round](-0.13,0.0)--(-0.04,-0.09)--(0.13,0.1);}
\providecommand{\coche}{\popcheck{popgreen}}
\providecommand{\resultat}[1]{\par\smallskip\noindent\fcolorbox{popgreen}{popgreenL}{\parbox{\dimexpr\linewidth-2\fboxsep-2\fboxrule\relax}{\textbf{Résultat.} #1}}\par\smallskip}
\newcommand{\popboxhead}[3]{\poppill{#1}{#2}{white}\ifx\relax#3\relax\else\hspace{6pt}{\popxbold\fontsize{10.6}{12}\selectfont\color{popink}#3}\fi\par\vspace{7pt}}

\newtcolorbox{aretenir}[1][]{popbase,colframe=popgreen,before upper={\popboxhead{À retenir}{popgreen}{#1}}}
% Texte en allemand (document de lecture, dialogue, extrait) : cadre
% d'étude. Un titre optionnel (source, auteur) s'affiche dans l'en-tête.
\newtcolorbox{textebox}[1][]{popbase,colframe=popblue,before upper={\popboxhead{Text}{popblue}{#1}\begin{otherlanguage}{ngerman}},after upper={\end{otherlanguage}}}
% Marques ✓ / ✗ (forme correcte / fautive) souvent utilisées par les modèles
\providecommand{\cross}{\textcolor{poppink}{\ensuremath{\times}}}
\providecommand{\xmark}{\cross}\providecommand{\crossmark}{\cross}\providecommand{\cmark}{\ensuremath{\checkmark}}
% Ligne de réponse / justification à compléter (inventée par les modèles)
\providecommand{\justification}[1]{\par\noindent\textit{Justification~:} \dotfill\par}
% \textipa (package tipa non chargé) : la police ne couvre pas l'API -> simple passage
\providecommand{\textipa}[1]{#1}
% Soulignement (ulem non chargé) : \uline, \ul
\providecommand{\uline}[1]{\underline{#1}}\providecommand{\ul}[1]{\underline{#1}}
% \glossaire{\item ...} inventée par les modèles : liste de glossaire
\providecommand{\glossaire}[1]{\begin{itemize}#1\end{itemize}}
% \alert (beamer) inventée par les modèles : texte mis en évidence
\providecommand{\alert}[1]{\textbf{\textcolor{poppink}{#1}}}
% variantes ASCII d'ordinaux écrites en macro
\providecommand{\iere}{\textsuperscript{re}}\providecommand{\ieme}{\textsuperscript{e}}\providecommand{\ere}{\textsuperscript{re}}\providecommand{\eme}{\textsuperscript{e}}\providecommand{\er}{\textsuperscript{er}}\providecommand{\ie}{\textsuperscript{e}}
% Ordinaux français écrits en macro par les modèles : 1\ière, 3\ième
\newcommand{\ière}{\textsuperscript{re}}\newcommand{\ième}{\textsuperscript{e}}
% \exemple (inventée par les modèles) : étiquette « Exemple : »
\providecommand{\exemple}{\textbf{Exemple~:}\ }
% \square utilisable aussi hors mode math (cases à cocher)
\AtBeginDocument{\let\squareMath\square\renewcommand{\square}{\ensuremath{\squareMath}}}
% Mot ou phrase en allemand à l'intérieur d'un paragraphe français
\newcommand{\all}[1]{\ifmmode\text{\textit{#1}}\else\begingroup\selectlanguage{ngerman}\itshape #1\endgroup\fi}
\let\esp\all % alias toléré
\newtcolorbox{exemplebox}[1][]{popbase,colframe=poporange,before upper={\popboxhead{Exemple}{poporange}{#1}}}
\newtcolorbox{definition}[1][]{popbase,colframe=popblue,before upper={\popboxhead{Définition}{popblue}{#1}}}
\newtcolorbox{propriete}[1][]{popbase,colframe=poppurple,before upper={\popboxhead{Propriété}{poppurple}{#1}}}
\newtcolorbox{methode}[1][]{popbase,colframe=popgold,before upper={\popboxhead{Méthode}{popgold}{#1}}}
\newtcolorbox{attention}[1][]{popbase,colframe=poppink,before upper={\popboxhead{Attention}{poppink}{#1}}}
\newtcolorbox{experience}[1][]{popbase,colframe=popblue,colback=popblueL,before upper={\popboxhead{Expérience}{popblue}{#1}}}
% « Le savais-tu ? » : carte violette pleine (contexte, culture, Cameroun)
\newtcolorbox{savaistu}[1][]{popbase,colback=poppurple,colframe=popink,
  fontupper=\popsemi\fontsize{10.4}{14.2}\selectfont\color{white},
  before upper={\poppill{Le savais-tu\,?}{white}{poppurple}\ifx\relax#1\relax\else\hspace{6pt}{\popxbold\color{white}#1}\fi\par\vspace{7pt}}}
% Exercice résolu : énoncé en haut, \tcblower puis la solution sur fond crème
\newcounter{popexo}\newcounter{popexr}
\newtcolorbox{exoresolu}[1][]{popbase,colframe=poporange,colbacklower=popcream,
  segmentation style={popink,line width=1.2pt,dashed},
  before upper={\stepcounter{popexr}\popboxhead{Exercice résolu \thepopexr}{poporange}{#1}},
  before lower={\poppill{Solution}{popink}{popcream}\par\vspace{6pt}}}
% Exercice (non résolu en place) — la correction est dans la partie Corrigés
\newtcolorbox{exercice}[1][]{popbase,colframe=popink,
  before upper={\stepcounter{popexo}\popboxhead{Exercice \thepopexo}{popink}{#1}}}
% Corrigé
\newtcolorbox{corrige}[1][]{popbase,colframe=popgreen,colback=popgreenL,
  before upper={\popboxhead{Corrigé}{popgreen}{#1}}}
% Objectifs / prérequis / bilan
\newtcolorbox{objectifs}{popbase,colframe=popink,colback=poporangeL,
  before upper={\popboxhead{Objectifs de la leçon}{popink}{}}}
\newtcolorbox{prerequis}{popbase,colframe=popink,colback=popblueL,
  before upper={\popboxhead{Prérequis}{popblue}{}}}
\newtcolorbox{bilan}{popbase,colframe=popink,colback=white,boxrule=2.8pt,
  before upper={\popboxhead{Fiche bilan}{poporange}{L'essentiel à connaître pour l'examen}}}
% Figure encadrée avec légende
\newtcolorbox{popfigure}[1][]{enhanced,breakable=false,colback=white,colframe=popline,boxrule=1.4pt,arc=8pt,
  left=10pt,right=10pt,top=10pt,bottom=8pt,before skip=10pt,after skip=12pt,halign=center,
  lower separated=false,halign lower=center,
  fontlower=\popsemi\fontsize{9}{11.5}\selectfont\color{popmuted},
  #1}
\newcommand{\legende}[1]{\par\vspace{6pt}{\popsemi\fontsize{9}{11.5}\selectfont\color{popmuted}#1}}

% ============================================================
% Signature OnBuch+
% ============================================================
\newcommand{\poplogoinline}[1]{{\poplogo\fontsize{#1}{#1}\selectfont\color{popink}OnBuch\color{poporange}+}}
\newcommand{\signatureonbuch}{%
  \begin{tcolorbox}[enhanced,colback=popink,colframe=popink,boxrule=2.2pt,arc=10pt,
      drop shadow=poporange,left=14pt,right=14pt,top=10pt,bottom=10pt,before skip=0pt,after skip=0pt]
    \tikz[baseline=-0.7ex]\node[circle,fill=popgreen,draw=popcream,line width=1.3pt,minimum size=22pt,inner sep=0]{\popcheck{white}};%
    \hspace{9pt}\begin{minipage}[c]{0.62\linewidth}
      {\popxbold\fontsize{12}{13}\selectfont\color{popcream}Signé\ \textbullet\ {\poplogo OnBuch\color{poporange}+}}\par\vspace{2pt}
      {\raggedright\popsemi\fontsize{8}{10.5}\selectfont\color{popline}\MakeUppercase{Équipe pédagogique OnBuch+}\\\MakeUppercase{Conforme au programme officiel MINESEC}\par}
    \end{minipage}\hfill
    {\poplogo\fontsize{22}{22}\selectfont\color{popcream}OnBuch\color{poporange}+}
  \end{tcolorbox}%
}

% ============================================================
% Page de garde
% #1 titre · #2 matière · #3 niveau · #4 module · #5 n° leçon · #6 durée ·
% #7 description / objectifs (texte ou itemize)
% ============================================================
\newcommand{\pagegarde}[7]{%
  \thispagestyle{empty}
  \newgeometry{a4paper,margin=1.7cm,top=1.6cm,bottom=1.4cm}%
  \begin{tikzpicture}[remember picture,overlay]
    \fill[poporange!18] ([xshift=-3.2cm,yshift=-3.6cm]current page.north east) circle (4.6cm);
    \fill[poppurple!14] ([xshift=2.4cm,yshift=7.6cm]current page.south west) circle (3.8cm);
  \end{tikzpicture}%
  {\poplogo\fontsize{30}{30}\selectfont\color{popink}OnBuch\color{poporange}+}\hfill
  \raisebox{0.6ex}{\poppill{Cours signé \textbullet\ Premium}{popink}{popcream}}\par
  \vspace{1pt}\popsquiggle{poppurple}\par
  \vspace{8pt}
  \begin{tcolorbox}[enhanced,colback=poporange,colframe=popink,boxrule=2.8pt,arc=13pt,
      drop shadow=popink,left=18pt,right=18pt,top=12pt,bottom=13pt,after skip=10pt]
    \poppill{#2 \textbullet\ #3}{popink}{popcream}\hspace{5pt}\poppill{#4}{white}{popink}\par\vspace{8pt}
    {\popdisplay\fontsize{14}{14}\selectfont\color{popink}LEÇON #5}\par\vspace{4pt}
    {\raggedright\hyphenpenalty=10000\exhyphenpenalty=10000\popdisplay\fontsize{25}{27}\selectfont\color{popcream}\MakeUppercase{#1}\par}
    \vspace{8pt}
    {\popxbold\fontsize{9.5}{9.5}\selectfont\color{popink}\MakeUppercase{Durée conseillée : #6}}
  \end{tcolorbox}
  \begin{tcolorbox}[enhanced,colback=white,colframe=popink,boxrule=2.2pt,arc=10pt,
      drop shadow=popink,left=16pt,right=16pt,top=10pt,bottom=10pt,after skip=8pt]
    \poppill{Dans cette leçon}{poppurple}{white}\par\vspace{5pt}
    \popsemi\fontsize{9.3}{12.6}\selectfont\color{popink}#7
  \end{tcolorbox}
  \noindent\begin{minipage}[t]{0.487\linewidth}
    \begin{tcolorbox}[enhanced,colback=popgreen,colframe=popink,boxrule=2.2pt,arc=10pt,
        drop shadow=popink,left=11pt,right=11pt,top=10pt,bottom=10pt,height=3.1cm,valign=center]
      \tcbox[colback=white,colframe=popink,boxrule=1.3pt,arc=5pt,boxsep=0pt,nobeforeafter,
        left=3pt,right=3pt,top=3pt,bottom=3pt,valign=center]{\qrcode[height=1.9cm]{https://play.google.com/store/apps/details?id=com.luvvix.onbuch&hl=fr}}%
      \hspace{8pt}\begin{minipage}[c]{0.5\linewidth}
        {\popdisplay\fontsize{11}{12.5}\selectfont\color{white}\MakeUppercase{Télécharge OnBuch}}\par\vspace{4pt}
        {\raggedright\popsemi\fontsize{8.4}{11}\selectfont\color{white}Gratuit \textbullet\ Google Play\\Tuteur IA, annales, concours, résultats\par}
      \end{minipage}
    \end{tcolorbox}
  \end{minipage}\hfill
  \begin{minipage}[t]{0.487\linewidth}
    \begin{tcolorbox}[enhanced,colback=poppurple,colframe=popink,boxrule=2.2pt,arc=10pt,
        drop shadow=popink,left=11pt,right=11pt,top=10pt,bottom=10pt,height=3.1cm,valign=center]
      \tikz[baseline=-0.7ex]\node[circle,fill=white,draw=popink,line width=1.3pt,minimum size=1.6cm,inner sep=0]{\popdisplay\fontsize{24}{24}\selectfont\color{poppurple}+};%
      \hspace{8pt}\begin{minipage}[c]{0.6\linewidth}
        {\popdisplay\fontsize{11}{12.5}\selectfont\color{white}\MakeUppercase{Passe à OnBuch+}}\par\vspace{4pt}
        {\raggedright\popsemi\fontsize{8.4}{11}\selectfont\color{white}Tous les cours signés, Tuteur IA illimité, fiches PDF — onglet OnBuch+ de l'app\par}
      \end{minipage}
    \end{tcolorbox}
  \end{minipage}
  \vspace{8pt}
  \popcontenu
  \vfill
  \signatureonbuch
  \restoregeometry
  \newpage
}

% Rangée « Ce cours contient » (page de garde)
\newcommand{\poptile}[3]{% #1 couleur #2 gros texte #3 légende
  \begin{tcolorbox}[enhanced,colback=#1,colframe=popink,boxrule=2pt,arc=9pt,shadow={2.5pt}{-2.5pt}{0pt}{popink},
      left=6pt,right=6pt,top=9pt,bottom=9pt,halign=center,valign=center,height=1.75cm,width=\linewidth,nobeforeafter]
    {\popdisplay\fontsize{12.5}{13}\selectfont\color{white}\MakeUppercase{#2}}\par\vspace{3pt}
    {\centering\popsemi\fontsize{7.8}{9.5}\selectfont\color{white}#3\par}
  \end{tcolorbox}}
\newcommand{\popcontenu}{%
  \par\noindent{\popxboldls\fontsize{7.5}{7.5}\selectfont\color{popmuted}CE COURS SIGNÉ CONTIENT}\par\vspace{7pt}
  \noindent\begin{minipage}[t]{0.235\linewidth}\poptile{popblue}{Cours}{complet, illustré, pas à pas}\end{minipage}\hfill
  \begin{minipage}[t]{0.235\linewidth}\poptile{poporange}{Méthodes}{et exercices résolus}\end{minipage}\hfill
  \begin{minipage}[t]{0.235\linewidth}\poptile{poppink}{Exercices}{type examen + corrigés}\end{minipage}\hfill
  \begin{minipage}[t]{0.235\linewidth}\poptile{popgreen}{Fiche bilan}{l'essentiel à retenir}\end{minipage}\par}

% Sommaire
\newtcolorbox{sommaire}{enhanced,colback=white,colframe=popink,boxrule=2.2pt,arc=10pt,
  drop shadow=popink,left=18pt,right=18pt,top=13pt,bottom=13pt,before skip=6pt,after skip=20pt,
  before upper={\poppill{Sommaire}{poppurple}{white}\par\vspace{9pt}\popsemi\fontsize{10.6}{14.5}\selectfont\color{popink}}}

% Signature de fin de cours — poussée tout en bas de la dernière page
\newcommand{\findecours}{%
  \par\vfill\nopagebreak
  \begin{center}\popsquiggle{poporange}\end{center}\vspace{4pt}
  \signatureonbuch
}
\AtBeginDocument{\def\llbracket{\mathopen{[\mkern-3.5mu[}}\def\rrbracket{\mathclose{]\mkern-3.5mu]}}} % intervalles d'entiers (glyphe ⟦ absent de la police)
% Glyphes absents de Fira Math : redéfinis à partir de glyphes disponibles
\AtBeginDocument{%
  \def\setminus{\mathbin{\backslash}}\let\smallsetminus\setminus
  \def\vdots{\mathord{\vbox{\baselineskip3.2pt\lineskiplimit0pt\kern4pt\hbox{.}\hbox{.}\hbox{.}}}}%
  \def\ddots{\mathinner{\mkern1mu\raise6.5pt\hbox{.}\mkern2mu\raise3.6pt\hbox{.}\mkern2mu\raise0.7pt\hbox{.}\mkern1mu}}%
  \def\triangle{\mathord{\scalebox{0.9}{$\symup{\Delta}$}}}%
  \def\nabla{\mathord{\raisebox{\depth}{\scalebox{1}[-1]{$\symup{\Delta}$}}}}%
  \def\checkmark{\popcheck{popdark}}%
  \def\star{\mathbin{\ast}}\def\bigstar{\mathord{\ast}}%
  \def\diamond{\mathbin{\scalebox{0.6}{$\Diamond$}}}%
  \def\bigcup{\mathop{\mathchoice{\raisebox{-0.3em}{\scalebox{1.7}{$\cup$}}}{\scalebox{1.25}{$\cup$}}{\cup}{\cup}}\displaylimits}%
  \def\bigcap{\mathop{\mathchoice{\raisebox{-0.3em}{\scalebox{1.7}{$\cap$}}}{\scalebox{1.25}{$\cap$}}{\cap}{\cap}}\displaylimits}%
  \def\lesseqqgtr{\mathrel{\lessgtr}}\def\bigoplus{\mathop{\oplus}\displaylimits}%
}
\providecommand{\ssi}{si et seulement si} % abréviation fréquente
\AtBeginDocument{\providecommand{\wideparen}{\overparen}} % arc de cercle

\begin{document}

\pagegarde{La presse écrite et le journal télévisé}{Allemand}{1ère A}{Chapitre 4 : Média et communication}{ALA16}{2 h}{Cette leçon fait découvrir le vocabulaire de la presse écrite et du journal télévisé à travers un texte authentique sur les médias en Allemagne. Elle consolide deux points grammaticaux majeurs du programme : le prétérit des verbes forts et l'ordre des mots dans la phrase allemande.
\vspace{4pt}
\begin{multicols}{2}\small
\begin{itemize}
\item À la fin de cette leçon, je sais comprendre globalement puis en détail un texte allemand sur la presse et le journal télévisé.
\item À la fin de cette leçon, je sais employer correctement le vocabulaire des médias (die Zeitung, die Nachrichten, der Reporter, die Schlagzeile, das Fernsehen, die Sendung) avec articles et pluriels.
\item À la fin de cette leçon, je sais former et conjuguer le prétérit des verbes forts les plus fréquents.
\item À la fin de cette leçon, je sais distinguer prétérit des verbes forts et prétérit des verbes faibles.
\item À la fin de cette leçon, je sais respecter l'ordre des mots : verbe conjugué en 2e position, sujet et compléments bien placés.
\item À la fin de cette leçon, je sais répondre par écrit et à l'oral à des questions de compréhension sur un texte.
\end{itemize}\end{multicols}}

\begin{sommaire}
\begin{multicols}{2}
\begin{enumerate}
\item Texte support : Die Medien in Deutschland
\item Compréhension du texte
\item Lexique thématique : les médias
\item Grammaire : le prétérit des verbes forts
\item Grammaire : l'ordre des mots
\item Méthode du commentaire et production guidée
\item Activité d'intégration
\item Méthodes et exercices résolus
\item Exercices
\item Corrigés des exercices
\item Fiche bilan
\end{enumerate}
\end{multicols}
\end{sommaire}

\begin{objectifs}
\begin{itemize}
\item À la fin de cette leçon, je sais comprendre globalement puis en détail un texte allemand sur la presse et le journal télévisé.
\item À la fin de cette leçon, je sais employer correctement le vocabulaire des médias (die Zeitung, die Nachrichten, der Reporter, die Schlagzeile, das Fernsehen, die Sendung) avec articles et pluriels.
\item À la fin de cette leçon, je sais former et conjuguer le prétérit des verbes forts les plus fréquents.
\item À la fin de cette leçon, je sais distinguer prétérit des verbes forts et prétérit des verbes faibles.
\item À la fin de cette leçon, je sais respecter l'ordre des mots : verbe conjugué en 2e position, sujet et compléments bien placés.
\item À la fin de cette leçon, je sais répondre par écrit et à l'oral à des questions de compréhension sur un texte.
\item À la fin de cette leçon, je sais résumer un texte sur les médias en quelques phrases au prétérit.
\item À la fin de cette leçon, je sais produire un court paragraphe guidé racontant ce que j'ai lu ou vu aux informations.
\end{itemize}
\end{objectifs}

\begin{prerequis}
\begin{itemize}
\item Le présent de l'indicatif des verbes réguliers et irréguliers (Seconde)
\item Le prétérit des verbes faibles et des auxiliaires haben/sein (début de 1ère)
\item L'ordre des mots de base : verbe en 2e position, inversion sujet-verbe (Seconde)
\item Les articles définis et indéfinis et les pluriels des noms (Seconde)
\item Le vocabulaire de la vie quotidienne et des loisirs (Seconde)
\end{itemize}
\end{prerequis}

\newpage

\coursec{Texte support : Die Medien in Deutschland}

Chaque matin à Yaoundé, beaucoup de Camerounais écoutent les nouvelles à la radio ou regardent le journal télévisé de la CRTV. En Allemagne aussi, des millions de personnes lisent le journal au petit-déjeuner ou regardent la \all{Tagesschau} à 20 heures. Mais comment parler de ce que l'on a lu ou vu hier ? En allemand, il faut raconter au passé avec le \cle{prétérit} (\all{Präteritum}) des verbes forts et placer les mots au bon endroit dans la phrase. Cette leçon nous plonge dans le monde des médias germanophones tout en maîtrisant ces deux outils grammaticaux essentiels.

Notre problématique est donc la suivante : \textbf{comment les Allemands s'informent-ils, et comment raconter au passé ce que l'on a lu ou vu dans les médias ?}

\courssub{Première lecture : découverte globale du texte}

\begin{methode}[Lire un texte allemand en trois étapes]
Pour comprendre un texte de presse ou de la vie quotidienne, procède toujours ainsi :
\begin{enumerate}
\item \textbf{Lecture globale, sans dictionnaire.} Lis le texte une première fois en silence. Ne t'arrête pas sur les mots inconnus. Cherche seulement le sens général : \emph{de qui parle-t-on ? où ? quand ? que se passe-t-il ?}
\item \textbf{Lecture détaillée, avec le glossaire.} Relis le texte phrase par phrase. Aide-toi du glossaire en marge. Souligne les mots du champ lexical du thème (ici : les médias).
\item \textbf{Repérage des mots-clés et des temps verbaux.} Entoure les verbes et identifie leur temps (présent, prétérit, parfait). Cela t'aide à situer les actions dans le temps et à préparer la grammaire de la leçon.
\end{enumerate}
\end{methode}

\textbf{Consigne de première lecture.} Lis le texte suivant en silence, sans dictionnaire, puis réponds mentalement à trois questions : \emph{Qui ?} (qui sont les personnages ?), \emph{Où ?} (où se passe l'action ?), \emph{Quoi ?} (que font-ils ?).

\begin{textebox}[Die Medien in Deutschland]
Gestern Morgen las Familie Müller die Zeitung am Frühstückstisch. Der Vater sah die Schlagzeile: „Neues Museum in München eröffnet“. Die Mutter las die Nachrichten auf ihrem Handy. Die Tochter Anna schrieb einen Kommentar im Internet.

Am Abend sahen alle zusammen die Tagesschau im Fernsehen. Der Sprecher berichtete über Politik, Sport und Wetter. Danach kam ein Film.

Am Nachmittag kam ein Reporter in die Stadt und fragte die Leute auf der Straße: „Lesen Sie noch Zeitung?“ Viele antworteten: „Nein, wir sehen lieber fern.“ Ein alter Mann sagte: „Ich lese jeden Morgen die Zeitung, das ist meine Tradition.“ Eine junge Frau antwortete: „Ich informiere mich nur im Internet, das ist schneller.“

Früher kaufte fast jeder eine Zeitung am Kiosk. Heute informieren sich viele Menschen im Internet oder mit dem Handy. Aber die Tagesschau am Abend bleibt für viele Familien ein Ritual: um 20 Uhr sitzen sie zusammen vor dem Fernseher und sehen die Nachrichten.
\end{textebox}

\textbf{Glossaire du texte :}

\begin{center}
\begin{tabularx}{\linewidth}{|X|X|}
\hline
\rowcolor{popblueL} \textbf{Deutsch} & \textbf{Français} \\
\hline
die Zeitung (-en) & le journal \\
\hline
der Frühstückstisch (-e) & la table du petit-déjeuner \\
\hline
die Schlagzeile (-n) & le gros titre \\
\hline
die Nachrichten (Pl.) & les informations \\
\hline
das Handy (-s) & le téléphone portable \\
\hline
der Kommentar (-e) & le commentaire \\
\hline
das Fernsehen & la télévision (le média) \\
\hline
der Fernseher (-) & le téléviseur (l'appareil) \\
\hline
der Sprecher (-) & le présentateur \\
\hline
die Sendung (-en) & l'émission \\
\hline
der Reporter (-) & le reporter \\
\hline
die Leute (Pl.) & les gens \\
\hline
der Kiosk (-e) & le kiosque \\
\hline
früher & autrefois \\
\hline
gestern & hier \\
\hline
sich informieren & s'informer \\
\hline
\end{tabularx}
\end{center}

\begin{attention}[Deux pièges de vocabulaire]
\begin{itemize}
\item \all{die Nachrichten} est presque toujours au \textbf{pluriel} en allemand pour dire « les informations ». On dit : \all{Ich sehe die Nachrichten.} « Je regarde les informations. » Ne dis jamais \all{die Nachricht} pour « les infos » : au singulier, cela signifie « une nouvelle, un message ».
\item \all{das Fernsehen} désigne le \textbf{média} (la télévision en général) ; \all{der Fernseher} désigne l'\textbf{appareil} (le téléviseur dans le salon). Compare : \all{Wir sehen fern.} « Nous regardons la télé. » / \all{Der Fernseher ist kaputt.} « Le téléviseur est en panne. »
\end{itemize}
\end{attention}

\courssub{Deuxième lecture : le champ lexical des médias}

Relis maintenant le texte et souligne tous les mots qui appartiennent au monde des médias. Tu dois trouver notamment : \all{die Zeitung}, \all{die Schlagzeile}, \all{die Nachrichten}, \all{das Fernsehen}, \all{die Tagesschau}, \all{der Sprecher}, \all{der Reporter}, \all{das Internet}, \all{der Kiosk}.

\begin{definition}[Le champ lexical]
Un \cle{champ lexical} est l'ensemble des mots qui se rapportent à un même thème. Ici, le thème est « les médias » : il comprend les supports (journal, télévision, Internet), les personnes (reporter, présentateur, journaliste) et les actions (lire, regarder, informer, rapporter).
\end{definition}

La carte mentale suivante organise ce vocabulaire pour t'aider à le mémoriser :

\begin{popfigure}
\begin{tikzpicture}[
  mindmap,
  grow cyclic,
  every node/.style={concept, font=\small},
  root concept/.append style={concept color=popblue, font=\small\bfseries},
  level 1/.append style={level distance=3.2cm, sibling angle=90},
  level 2/.append style={level distance=2.2cm, sibling angle=45, font=\footnotesize},
]
\node[root concept]{Die Medien}
  child[concept color=poporange]{ node {Presse}
    child { node {die Zeitung} }
    child { node {die Zeitschrift} }
    child { node {der Kiosk} }
  }
  child[concept color=popgreen]{ node {Fernsehen}
    child { node {die Sendung} }
    child { node {die Nachrichten} }
    child { node {die Tagesschau} }
  }
  child[concept color=poppurple]{ node {Personen}
    child { node {der Reporter} }
    child { node {der Journalist} }
    child { node {der Sprecher} }
  }
  child[concept color=poppink]{ node {Actions}
    child { node {lesen} }
    child { node {sehen / fernsehen} }
    child { node {berichten} }
    child { node {sich informieren} }
  };
\end{tikzpicture}
\legende{Carte mentale du champ lexical des médias : quatre familles de mots à retenir.}
\end{popfigure}

\begin{exemplebox}[Le vocabulaire en phrases]
\all{Mein Vater liest jeden Morgen die Zeitung.} « Mon père lit le journal chaque matin. »

\all{Die Schlagzeile ist heute sehr interessant.} « Le gros titre est très intéressant aujourd'hui. »

\all{Der Reporter fragt die Leute auf der Straße.} « Le reporter interroge les gens dans la rue. »

\all{Um 20 Uhr kommt die Tagesschau im Fernsehen.} « À 20 heures, le journal télévisé passe à la télévision. »

\all{Viele Jugendliche informieren sich im Internet.} « Beaucoup de jeunes s'informent sur Internet. »
\end{exemplebox}

\begin{savaistu}[La Tagesschau, une institution allemande]
La \all{Tagesschau} est le journal télévisé le plus regardé d'Allemagne. Elle est diffusée chaque soir à 20 heures précises sur la chaîne publique ARD depuis 1952. Pour beaucoup de familles allemandes, regarder la Tagesschau ensemble est un vrai rituel du soir, un peu comme le journal de 20 heures de la CRTV au Cameroun. Le mot se compose de \all{der Tag} (le jour) et \all{die Schau} (la présentation, le spectacle) : littéralement « la revue du jour ».
\end{savaistu}

\begin{savaistu}[Famille de mots : \textit{lesen} / \textit{sehen}]
\begin{itemize}
\item \textbf{lesen} (lire) $\rightarrow$ \all{der Leser} (le lecteur), \all{das Lesen} (la lecture), \all{lesenswert} (qui vaut la peine d'être lu).
\item \textbf{sehen} (voir) $\rightarrow$ \all{der Seher} (le spectateur), \all{das Fernsehen} (la télévision), \all{fernsehen} (regarder la télé, verbe à particule séparable).
\item \textbf{berichten} (rapporter) $\rightarrow$ \all{der Bericht} (le reportage, le rapport), \all{der Reporter} (le reporter).
\end{itemize}
Apprendre les familles de mots permet de multiplier ton vocabulaire sans effort.
\end{savaistu}

\courssub{Troisième lecture : repérage des verbes au prétérit}

Le texte raconte des événements d'\textbf{hier} (\all{gestern}). En allemand, pour raconter au passé dans un texte écrit, on utilise le \cle{prétérit} (\all{Präteritum}). Entoure dans le texte tous les verbes au passé. Tu dois trouver :

\begin{center}
\begin{tabularx}{\linewidth}{|X|X|X|}
\hline
\rowcolor{popblueL} \textbf{Infinitif} & \textbf{Prétérit dans le texte (3\textsuperscript{e} pers. sg/pl)} & \textbf{Français} \\
\hline
lesen & las / lasen & lire (il/elle lut / ils lurent) \\
\hline
sehen & sah / sahen & voir (il vit / ils virent) \\
\hline
kommen & kam / kamen & venir (il vint / ils vinrent) \\
\hline
schreiben & schrieb & écrire (elle écrivit) \\
\hline
fragen & fragte & demander (il demanda) \\
\hline
antworten & antworteten & répondre (ils répondirent) \\
\hline
berichten & berichtete & rapporter (il rapporta) \\
\hline
kaufen & kaufte & acheter (on achetait) \\
\hline
sagen & sagte & dire (il dit) \\
\hline
bleiben & blieb & rester (il resta) \\
\hline
sitzen & saßen & être assis / s'asseoir (ils étaient assis) \\
\hline
\end{tabularx}
\end{center}

\begin{propriete}[Observation grammaticale : fort ou faible ?]
Observe les deux groupes de verbes du tableau :
\begin{itemize}
\item Les \textbf{verbes forts} (\cle{starke Verben}) changent de voyelle au prétérit (alternance apophonique), \textbf{sans terminaison en \all{-te}} aux 1\iere et 3\ieme pers. sing. : \all{lesen $\rightarrow$ las}, \all{sehen $\rightarrow$ sah}, \all{kommen $\rightarrow$ kam}, \all{schreiben $\rightarrow$ schrieb}, \all{bleiben $\rightarrow$ blieb}, \all{sitzen $\rightarrow$ saßen}.
\item Les \textbf{verbes faibles} (\cle{schwache Verben}) ajoutent simplement \all{-te} au radical (comme notre passé composé régulier en français) : \all{fragen $\rightarrow$ fragte}, \all{antworten $\rightarrow$ antwortete}, \all{kaufen $\rightarrow$ kaufte}, \all{sagen $\rightarrow$ sagte}, \all{berichten $\rightarrow$ berichtete}.
\end{itemize}
Nous étudierons la conjugaison complète du prétérit des verbes forts dans la section \textbf{Grammaire : Le prétérit des verbes forts}. Pour l'instant, retiens seulement que ces formes expriment le passé dans un récit écrit.
\end{propriete}

\begin{exoresolu}[Repérage dans le texte]
Énoncé. Relevez dans le texte trois phrases contenant un verbe fort au prétérit et traduisez-les en français.
\tcblower
Solution détaillée.
\begin{enumerate}
\item \all{Gestern Morgen las Familie Müller die Zeitung.} « Hier matin, la famille Müller a lu le journal. » Verbe fort : \all{las} (infinitif \all{lesen}, voyelle \all{e $\rightarrow$ a}).
\item \all{Der Vater sah die Schlagzeile.} « Le père a vu le gros titre. » Verbe fort : \all{sah} (infinitif \all{sehen}, voyelle \all{eh $\rightarrow$ a}).
\item \all{Ein Reporter kam in die Stadt.} « Un reporter est venu en ville. » Verbe fort : \all{kam} (infinitif \all{kommen}, voyelle \all{o $\rightarrow$ a}).
\item \all{Die Tochter Anna schrieb einen Kommentar.} « La fille Anna a écrit un commentaire. » Verbe fort : \all{schrieb} (infinitif \all{schreiben}, voyelle \all{ei $\rightarrow$ ie}).
\item \all{Ein alter Mann sagte: „Ich lese jeden Morgen die Zeitung...“} Verbe faible : \all{sagte} (terminaison \all{-te}).
\end{enumerate}
Méthode : pour vérifier qu'un verbe est fort, cherche s'il change de voyelle par rapport à l'infinitif (\all{e $\rightarrow$ a}, \all{ei $\rightarrow$ ie}, \all{o $\rightarrow$ a}, \all{ie $\rightarrow$ a}) et s'il n'a \textbf{pas} de terminaison en \all{-te} aux personnes du singulier.
\end{exoresolu}

\courssub{Lecture à voix haute : prononciation et intonation}

L'enseignant lit maintenant le texte à voix haute, puis des élèves lisent à leur tour, phrase par phrase. Voici quelques points de prononciation à soigner :

\begin{itemize}
\item \all{Zeitung} se prononce « tsaï-toung » : le \all{Z} allemand se dit \all{[ts]}.
\item \all{Schlagzeile} se prononce « chlak-tsaï-lé » : \all{sch} se dit \all{[ʃ]} (comme \emph{ch}at) ; \all{Z} = \all{[ts]}.
\item \all{Fernsehen} se prononce « fèrrn-zé-enn » : le \all{s} entre deux voyelles se prononce \all{[z]} (sonore).
\item \all{Tagesschau} se prononce « ta-guess-chaou » : \all{sch} = \all{[ʃ]} ; \all{au} = \all{[aʊ]}.
\item \all{Reporter} se prononce « ré-porteur » \all{[ʁeˈpɔʁtɐ]}, avec l'accent sur la deuxième syllabe.
\item \all{Nachrichten} se prononce « na-khrich-ten » : \all{ch} = \all{[ç]} (ich-Laut) après \all{i}.
\end{itemize}

\begin{propriete}[L'intonation des phrases]
En allemand comme en français, la voix descend à la fin d'une phrase déclarative et monte à la fin d'une question fermée (sans mot interrogatif). Mais attention : dans les questions ouvertes avec mot interrogatif comme \all{Lesen Sie noch Zeitung?}, la voix \textbf{redescend} à la fin, car la question est déjà marquée par le mot interrogatif (ici l'inversion verbe-sujet joue ce rôle) et l'absence de particule interrogative type \all{ob}.
\end{propriete}

\begin{attention}[L'ordre des mots : première observation]
Observe la première phrase du texte : \all{Gestern Morgen las Familie Müller die Zeitung.}
Le verbe conjugué \all{las} est en \textbf{deuxième position}, juste après le complément de temps \all{Gestern Morgen}, et le sujet \all{Familie Müller} vient \emph{après} le verbe (inversion).
En français, on dirait « Hier matin, la famille Müller a lu le journal » : le sujet reste avant le verbe.
\cle{Règle d'or} : En allemand, le verbe conjugué occupe \textbf{toujours la deuxième place} dans la phrase déclarative principale, quoi qu'il y ait en première position (sujet, complément de temps, de lieu, d'objet...). Nous étudierons cette règle en détail dans la section \textbf{Grammaire : L'ordre des mots}.
\end{attention}

\courssub{Questions de compréhension globale}

\begin{exercice}[Vérifie ta compréhension]
Réponds aux questions suivantes par des phrases complètes, en allemand. Tu peux répondre au présent (langue de communication) ou au prétérit (langue du récit écrit).
\begin{enumerate}
\item Wer liest die Zeitung am Frühstückstisch?
\item Was sieht der Vater in der Zeitung?
\item Was sehen alle am Abend im Fernsehen?
\item Wen fragt der Reporter auf der Straße?
\item Wie informieren sich viele Menschen heute?
\end{enumerate}
\end{exercice}

\begin{corrige}[Exercice de compréhension globale]
\begin{enumerate}
\item \all{Familie Müller liest die Zeitung am Frühstückstisch.} (ou \all{las} pour le récit) « La famille Müller lit le journal à la table du petit-déjeuner. »
\item \all{Der Vater sieht die Schlagzeile: „Neues Museum in München eröffnet“.} (ou \all{sah}) « Le père voit le gros titre : "Ouverture d'un nouveau musée à Munich". »
\item \all{Am Abend sehen alle die Tagesschau im Fernsehen.} (ou \all{sahen}) « Le soir, tous regardent le journal télévisé à la télévision. »
\item \all{Der Reporter fragt die Leute auf der Straße.} (ou \all{fragte}) « Le reporter interroge les gens dans la rue. »
\item \all{Heute informieren sich viele Menschen im Internet oder mit dem Handy.} « Aujourd'hui, beaucoup de gens s'informent sur Internet ou avec leur portable. »
\end{enumerate}
\end{corrige}

\begin{aretenir}[Ce qu'il faut retenir de cette section]
\begin{itemize}
\item Le texte \all{Die Medien in Deutschland} présente une famille allemande qui lit le journal le matin et regarde la \all{Tagesschau} le soir ; un reporter interroge les passants sur leurs habitudes.
\item Le vocabulaire-clé : \all{die Zeitung}, \all{die Nachrichten}, \all{die Schlagzeile}, \all{das Fernsehen}, \all{die Sendung}, \all{der Reporter}, \all{der Sprecher}, \all{der Kiosk}.
\item Le texte est raconté au \textbf{prétérit} : verbes forts (\all{las}, \all{sah}, \all{kam}, \all{schrieb}, \all{blieb}, \all{saßen}) et verbes faibles (\all{fragte}, \all{antwortete}, \all{kaufte}, \all{sagte}, \all{berichtete}).
\item En allemand, le verbe conjugué est toujours en deuxième position dans la phrase déclarative principale.
\end{itemize}
\end{aretenir}

% ============================================================
% SECTION GRAMMAIRE 1 : PRÉTÉRIT DES VERBES FORTS
% ============================================================

\coursec{Grammaire : Le prétérit des verbes forts (Präteritum)}

\begin{propriete}[Formation du prétérit des verbes forts]
Les verbes forts (\cle{starke Verben}) forment le prétérit \textbf{sans suffixe \all{-te}}.
\begin{enumerate}
\item \textbf{Changement de voyelle (Ablaut)} : la voyelle du radical change selon des schémas fixes (ex. \all{e $\rightarrow$ a}, \all{ei $\rightarrow$ ie}, \all{o $\rightarrow$ a}, \all{ie $\rightarrow$ a}, \all{ä $\rightarrow$ o}).
\item \textbf{Terminaisons personnelles} : elles sont identiques à celles du présent pour \all{wir}, \all{ihr}, \all{sie/Sie} (\all{-en}, \all{-t}, \all{-en}), mais spécifiques aux 1\iere/3\ieme pers. sing. (\textbf{zéro terminaison}) et 2\ieme pers. sing. (\all{-st}).
\end{enumerate}
\textbf{Tableau des terminaisons (ex. \all{lesen} $\rightarrow$ \all{las-}) :}
\begin{center}
\begin{tabular}{|l|l|l|l|}
\hline
\rowcolor{popblueL} \textbf{Personne} & \textbf{Terminaison} & \textbf{Exemple : lesen} & \textbf{Français} \\
\hline
ich & -- & \all{ich las} & je lus \\
\hline
du & \all{-st} & \all{du last} & tu lus \\
\hline
er/sie/es & -- & \all{er las} & il/elle lut \\
\hline
wir & \all{-en} & \all{wir lasen} & nous lûmes \\
\hline
ihr & \all{-t} & \all{ihr last} & vous lûtes \\
\hline
sie/Sie & \all{-en} & \all{sie lasen} / \all{Sie lasen} & ils/elles lurent / Vous lûtes \\
\hline
\end{tabular}
\end{center}
\cle{Point clé} : Les formes \textbf{ich} et \textbf{er/sie/es} sont \textbf{identiques} (pas de \all{-e} à la 1\iere pers. sing. contrairement aux verbes faibles \all{ich kaufte}).
\end{propriete}

\begin{exemplebox}[Les 12 verbes forts indispensables (niveau A1/A2)]
Apprends ces formes par cœur (Infinitif -- Prétérit 3\ieme sg -- Participe passé -- Français) :
\begin{center}
\begin{tabularx}{\linewidth}{|X|X|X|X|}
\hline
\rowcolor{popblueL} \textbf{Infinitif} & \textbf{Präteritum (er/sie/es)} & \textbf{Partizip II} & \textbf{Français} \\
\hline
\all{lesen} & \all{las} & \all{gelesen} & lire \\
\hline
\all{sehen} & \all{sah} & \all{gesehen} & voir \\
\hline
\all{kommen} & \all{kam} & \all{gekommen} & venir \\
\hline
\all{schreiben} & \all{schrieb} & \all{geschrieben} & écrire \\
\hline
\all{gehen} & \all{ging} & \all{gegangen} & aller (à pied) \\
\hline
\all{stehen} & \all{stand} & \all{gestanden} & être debout / se lever \\
\hline
\all{finden} & \all{fand} & \all{gefunden} & trouver \\
\hline
\all{geben} & \all{gab} & \all{gegeben} & donner \\
\hline
\all{nehmen} & \all{nahm} & \all{genommen} & prendre \\
\hline
\all{sprechen} & \all{sprach} & \all{gesprochen} & parler \\
\hline
\all{essen} & \all{aß} & \all{gegessen} & manger \\
\hline
\all{trinken} & \all{trank} & \all{getrunken} & boire \\
\hline
\end{tabularx}
\end{center}
\cle{Astuce} : Regroupe-les par alternance vocalique pour les mémoriser :
\begin{itemize}
\item \all{e $\rightarrow$ a} : \all{geben/nahm}, \all{nehmen/nahm}, \all{sprechen/sprach}, \all{essen/aß}, \all{helfen/half}, \all{treffen/traf}.
\item \all{ei $\rightarrow$ ie} : \all{schreiben/schrieb}, \all{bleiben/blieb}, \all{scheinen/schien}, \all{beißen/biss}.
\item \all{ie $\rightarrow$ a} : \all{lesen/las}, \all{sehen/sah} (historique \all{ie}), \all{liegen/lag}, \all{bieten/bot}.
\item \all{o $\rightarrow$ a} : \all{kommen/kam}, \all{treffen? non}, \all{stoßen/stieß}. \textbf{Mieux} : \all{a $\rightarrow$ u} : \all{laufen/lief} (ä), \all{fahren/fuhr}, \all{tragen/trug}, \all{waschen/wusch}, \all{schlafen/slief}.
\end{itemize}
\end{exemplebox}

\begin{attention}[Pièges fréquents pour francophones]
\begin{itemize}
\item \textbf{Oublier le changement de voyelle} : \all{*er lesete} (faux) $\rightarrow$ \all{er las}. \all{*er kamte} (faux) $\rightarrow$ \all{er kam}.
\item \textbf{Ajouter \all{-e} à la 1\iere pers. sing.} : \all{*ich lase} (faux) $\rightarrow$ \all{ich las}. Les verbes forts n'ont \textbf{jamais} de \all{-e} à \all{ich} et \all{er/sie/es} au prétérit.
\item \textbf{Confondre \all{sehen} et \all{sehen} (faible)} : \all{sehen} (voir) est fort \rightarrow \all{sah}. Mais \all{sähen} (subjonctif II) ou \all{sehen} comme verbe faible n'existe pas. Attention à \all{lesen} (fort) vs \all{lesen} (faible, rare).
\item \textbf{\all{ß} vs \all{ss}} : \all{er aß} (il mangea, \all{ß} après voyelle longue \all{a}) ; \all{er las} (pas de \all{ß}). \all{er saß} (il s'assit, \all{ß} après \all{a} long).
\end{itemize}
\end{attention}

\begin{exercice}[Entraînement au prétérit des verbes forts]
Conjugue les verbes entre parenthèses au prétérit (toutes les personnes indiquées).
\begin{enumerate}
\item \all{ich} (lesen) die Zeitung. \quad \all{wir} (lesen) die Zeitung.
\item \all{er} (sehen) den Film. \quad \all{sie} (sehen) den Film.
\item \all{ich} (kommen) gestern an. \quad \all{ihr} (kommen) gestern an.
\item \all{Anna} (schreiben) einen Brief. \quad \all{wir} (schreiben) einen Brief.
\item \all{du} (finden) das Handy. \quad \all{Sie} (finden) das Handy.
\item \all{wir} (geben) dem Reporter Auskunft. \quad \all{er} (nehmen) die Zeitung.
\end{enumerate}
\end{exercice}

\begin{corrige}[Exercice : Prétérit des verbes forts]
\begin{enumerate}
\item \all{ich las} -- \all{wir lasen}
\item \all{er sah} -- \all{sie sahen}
\item \all{ich kam} -- \all{ihr kamt}
\item \all{Anna schrieb} -- \all{wir schrieben}
\item \all{du fandst} -- \all{Sie fanden} (\all{finden} : \all{e $\rightarrow$ a}, radical \all{fand-}, du \all{fandst})
\item \all{wir gaben} -- \all{er nahm}
\end{enumerate}
\cle{Note} : \all{du fandst} (pas \all{*fandest}) car le radical se termine par \all{d/t} $\rightarrow$ \all{-st} devient \all{-st} (assimilation). De même \all{du gabst}, \all{du namst} (rare), \all{du aßest} (mais \all{du aß} est forme courte courante, standard \all{du aßest}).
\end{corrige}

% ============================================================
% SECTION GRAMMAIRE 2 : L'ORDRE DES MOTS
% ============================================================

\coursec{Grammaire : L'ordre des mots dans la phrase déclarative (Satzglieder)}

\begin{propriete}[La règle de la \cle{deuxième position} (Verbzweitstellung)]
Dans toute phrase déclarative principale en allemand, le \textbf{verbe conjugué occupe obligatoirement la deuxième place} (Position 2). La première place (Position 1 / \all{Vorfeld}) est occupée par \textbf{un seul élément} (groupe de mots) : sujet, complément de temps, de lieu, d'objet, adverbe, etc.
\begin{center}
\begin{tabularx}{\linewidth}{|X|X|X|}
\hline
\rowcolor{popblueL} \textbf{Position 1 (Vorfeld)} & \textbf{Position 2 (Verbe)} & \textbf{Reste de la phrase (Mittelfeld + Ende)} \\
\hline
\all{Familie Müller} (Sujet) & \all{las} & \all{die Zeitung am Frühstückstisch.} \\
\hline
\all{Gestern Morgen} (Temps) & \all{las} & \all{Familie Müller die Zeitung.} \\
\hline
\all{Am Frühstückstisch} (Lieu) & \all{las} & \all{Familie Müller die Zeitung.} \\
\hline
\all{Die Zeitung} (Objet/Acc) & \all{las} & \all{Familie Müller am Frühstückstisch.} \\
\hline
\end{tabularx}
\end{center}
\cle{Conséquence} : Si un complément (temps, lieu, objet) est en Position 1, le \textbf{sujet se décale après le verbe} (inversion). C'est la différence majeure avec le français (SVO rigide).
\end{propriete}

\begin{exemplebox}[Exemples variés avec le texte de la leçon]
\begin{enumerate}
\item \textbf{Standard (Sujet en P1) :} \all{Der Reporter kam in die Stadt.} (Le reporter vint en ville.)
\item \textbf{Temps en P1 :} \all{Am Nachmittag kam der Reporter in die Stadt.} (L'après-midi, le reporter vint en ville.)
\item \textbf{Lieu en P1 :} \all{In die Stadt kam der Reporter.} (En ville vint le reporter — insistance sur le lieu.)
\item \textbf{Objet en P1 :} \all{Die Zeitung las der Vater.} (Le journal, le père le lut — insistance sur l'objet.)
\item \textbf{Adverbe en P1 :} \all{Früher kaufte fast jeder eine Zeitung.} (Autrefois, presque tout le monde achetait un journal.)
\end{enumerate}
Dans tous les cas : \textbf{Verbe = Position 2}.
\end{exemplebox}

\begin{propriete}[Ordre dans le Mittelfeld (milieu de phrase) : Tekamolo]
Quand la phrase s'allonge, l'ordre habituel des compléments entre la Position 2 (verbe) et la fin de phrase suit la règle mnémotechnique \cle{Tekamolo} :
\begin{center}
\textbf{Te} (Temporel) -- \textbf{ka} (Kausal) -- \textbf{mo} (Modal) -- \textbf{lo} (Lokal)
\end{center}
\emph{Du temps $\rightarrow$ De la cause $\rightarrow$ De la manière $\rightarrow$ Du lieu.}
Exemple : \all{Der Reporter [P1] fragte [P2] \textbf{gestern} (Te) \textbf{freundlich} (Mo) \textbf{auf der Straße} (Lo) die Leute.}
\end{propriete}

\begin{attention}[Erreurs fréquentes des francophones]
\begin{itemize}
\item \textbf{Garder l'ordre français (Sujet-Verbe) après un complément en tête de phrase :} \all{*Gestern Morgen Familie Müller las...} (FAUX) $\rightarrow$ \all{Gestern Morgen las Familie Müller...} (VERBE EN 2).
\item \textbf{Oublier l'inversion du sujet :} \all{*Am Abend die Familie sah die Tagesschau.} (FAUX) $\rightarrow$ \all{Am Abend sah die Familie die Tagesschau.}
\item \textbf{Placer le verbe à la fin comme en subordonnée :} \all{*Weil er die Zeitung las,} (correct en subordonnée) mais \all{*Er die Zeitung las.} (FAUX en principale) $\rightarrow$ \all{Er las die Zeitung.}
\end{itemize}
\end{attention}

\begin{exercice}[Mets les mots dans l'ordre correct (Phrase déclarative)]
\begin{enumerate}
\item (gestern / die Nachrichten / der Vater / hören / im Radio)
\item (die Tagesschau / um 20 Uhr / sehen / wir / immer)
\item (einen Kommentar / Anna / im Internet / schreiben / gestern)
\item (früher / eine Zeitung / fast jeder / kaufen / am Kiosk)
\item (die Leute / der Reporter / auf der Straße / fragen / am Nachmittag)
\end{enumerate}
\end{exercice}

\begin{corrige}[Exercice : Ordre des mots]
\begin{enumerate}
\item \all{Gestern hörte der Vater die Nachrichten im Radio.} (ou \all{Der Vater hörte gestern...})
\item \all{Um 20 Uhr sehen wir immer die Tagesschau.} (Verbe \all{sehen} au présent ou \all{sahen} au prétérit : \all{Um 20 Uhr sahen wir immer die Tagesschau.})
\item \all{Gestern schrieb Anna einen Kommentar im Internet.}
\item \all{Früher kaufte fast jeder eine Zeitung am Kiosk.}
\item \all{Am Nachmittag fragte der Reporter die Leute auf der Straße.}
\end{enumerate}
\cle{Vérification} : Compte les positions. Le verbe conjugué (\all{hörte}, \all{sehen/sahen}, \all{schrieb}, \all{kaufte}, \all{fragte}) est-il bien le \textbf{deuxième élément} de la phrase ?
\end{corrige}

% ============================================================
% SECTION MÉTHODE \& PRODUCTION GUIDÉE
% ============================================================

\coursec{Méthode et production guidée : Résumer et raconter au passé}

\begin{methode}[Rédiger un résumé ou un court récit au prétérit (Schriftliche Erzählung)]
Pour produire un texte cohérent au passé (compte-rendu, résumé, journal intime), suis ces étapes :
\begin{enumerate}
\item \textbf{Collecte d'idées (Stichworte).} Note les mots-clés : \emph{Wer? Was? Wann? Wo? Wie?} (Qui? Quoi? Quand? Où? Comment?). Utilise le vocabulaire de la leçon (\all{Zeitung, Nachrichten, Reporter, Fernsehen, lesen, sehen, berichten, kommen, fragen, antworten}).
\item \textbf{Choix du temps.} Pour un récit écrit (journal, article, résumé de texte) : \textbf{Prétérit}. Pour l'oral (raconter à un ami) : \textbf{Parfait} (\all{Perfekt}). Ici, on s'entraîne au \textbf{Prétérit}.
\item \textbf{Conjugaison.} Vérifie chaque verbe : Fort (voyelle change, pas de \all{-te}) ou Faible (\all{-te}) ? Applique les terminaisons (\all{ich/er/sie/es} = radical nu pour les forts).
\item \textbf{Structure des phrases.} Respecte la \textbf{Position 2 du verbe}. Varie la Position 1 (Temps, Lieu, Sujet) pour éviter la monotonie. Utilise des connecteurs : \all{zuerst} (d'abord), \all{dann} / \all{danach} (ensuite), \all{aber} (mais), \all{weil} (parce que - verbe à la fin !).
\item \textbf{Relecture.} Vérifie : Majuscules aux noms ? Umlaute (\all{ä, ö, ü}) ? \all{ß} ? Verbe en Position 2 ? Terminaisons correctes ? Accords sujet-verbe ?
\end{enumerate}
\end{methode}

\begin{experience}[Production orale guidée : Interview « Mediennutzung »]
En binôme. L'un joue le \textbf{Reporter}, l'autre un \textbf{Passant}.
\textbf{Reporter} : Pose les questions (au présent ou prétérit).
\textbf{Passant} : Réponds selon ton profil (vieux lecteur de journal / jeune connecté).
\begin{itemize}
\item \all{„Lesen Sie noch Zeitung?“} -- \all{„Ja, ich lese jeden Morgen die Zeitung.“} / \all{„Nein, ich lese keine Zeitung mehr.“}
\item \all{„Wie informieren Sie sich?“} -- \all{„Ich sehe die Tagesschau.“} / \all{„Ich informiere mich im Internet.“}
\item \all{„Was war die Schlagzeile heute?“} -- \all{„Die Schlagzeile war: ...“}
\end{itemize}
Changez de rôle. L'enseignant note les phrases au prétérit entendues (\all{war, las, sah, informierte}).
\end{experience}

\begin{exercice}[Production écrite : Mon journal médiatique d'hier]
Écris un paragraphe de \textbf{6 à 8 phrases} en allemand pour raconter ta journée d'hier (ou une journée type) sous l'angle des médias.
\textbf{Contraintes :}
\begin{itemize}
\item Utilise le \textbf{prétérit} exclusivement.
\item Intègre au moins \textbf{3 verbes forts} (\all{lesen, sehen, kommen, gehen, schreiben, finden, geben, nehmen}) et \textbf{3 verbes faibles} (\all{hören, kaufen, fragen, antworten, sagen, berichten, informieren}).
\item Utilise au moins \textbf{5 mots du vocabulaire} de la leçon (\all{Zeitung, Nachrichten, Schlagzeile, Fernsehen, Tagesschau, Reporter, Handy, Internet, Kiosk, Sender}).
\item Varie la Position 1 (au moins 2 phrases ne commencent pas par le sujet).
\end{itemize}
\textbf{Aide au démarrage :} \all{Gestern Morgen... / Am Frühstückstisch... / Dann... / Am Abend... / Um 20 Uhr...}
\end{exercice}

\begin{corrige}[Exercice : Production écrite (Proposition de correction)]
\textbf{Modèle de réponse (niveau A2) :}
\all{Gestern Morgen las ich die Zeitung am Frühstückstisch. Die Schlagzeile war interessant. Danach ging ich zur Schule. Am Nachmittag kam ein Reporter in unsere Klasse und fragte uns: „Seht ihr fern?“ Ich antwortete: „Nein, ich informiere mich im Internet.“ Am Abend sahen meine Eltern die Tagesschau im Fernsehen. Ich saß daneben und las einen Kommentar auf meinem Handy. Früher kaufte mein Großvater die Zeitung am Kiosk, heute lesen wir digital.}

\textbf{Analyse du modèle :}
\begin{itemize}
\item \textbf{Verbes forts :} \all{las} (lesen), \all{war} (sein - irrégulier), \all{ging} (gehen), \all{kam} (kommen), \all{sah} (sehen), \all{saß} (sitzen), \all{las} (lesen).
\item \textbf{Verbes faibles :} \all{fragte}, \all{antwortete}, \all{informiere} (présent dans le dialogue rapporté, mais \all{informierte} au prétérit serait \all{informierte}), \all{kaufte}.
\item \textbf{Position 1 variée :} \all{Gestern Morgen}, \all{Die Schlagzeile}, \all{Dann}, \all{Am Nachmittag}, \all{Ich}, \all{Am Abend}, \all{Ich}, \all{Früher}.
\item \textbf{Vocabulaire :} \all{Zeitung, Schlagzeile, Reporter, Tagesschau, Fernsehen, Handy, Kiosk}.
\end{itemize}
\end{corrige}

\begin{aretenir}[Bilan de la leçon ALA16]
\begin{itemize}
\item \textbf{Thème :} La presse écrite (\all{die Zeitung, die Schlagzeile, der Kiosk}) et le journal télévisé (\all{die Tagesschau, das Fernsehen, der Sprecher, die Sendung, die Nachrichten}).
\item \textbf{Grammaire 1 :} \textbf{Prétérit des verbes forts} -- Alternance vocalique, pas de \all{-te}, terminaisons \all{-, -st, -, -en, -t, -en}. Formes \all{ich/er/sie/es} identiques. Liste des 12 verbes forts de base à connaître.
\item \textbf{Grammaire 2 :} \textbf{Ordre des mots (Verbzweitstellung)} -- Verbe conjugué en Position 2. Inversion sujet/verbe si Complément en Position 1. Règle \textbf{Tekamolo} pour le milieu de phrase.
\item \textbf{Compétences :} Comprendre un texte authentique, extraire le lexique, identifier les temps du récit, conjuguer correctement, produire un court récit écrit/oral structuré.
\end{itemize}
\end{aretenir}

\coursec{Compréhension du texte}

Dans la section précédente, nous avons découvert le texte support \all{Die Medien in Deutschland} et nous avons effectué une première lecture globale : nous savons déjà que le texte décrit une famille allemande et ses habitudes face aux médias. Nous allons maintenant passer à l'étape décisive du travail sur texte : la \cle{compréhension}. C'est la compétence que l'examinateur évalue en priorité au baccalauréat : savoir lire, repérer l'information dans le texte et répondre par des phrases complètes et correctes.

\courssub{Le texte support : rappel}

Relisons attentivement le texte avant de répondre aux questions. Une bonne compréhension exige toujours au moins deux lectures.

\begin{textebox}[Die Medien in Deutschland]
Familie Berger wohnt in München. Jeden Morgen kam der Vater in die Küche und las die Zeitung. Er las zuerst die Schlagzeilen auf der ersten Seite, dann den Sportteil. Die Mutter hörte beim Frühstück das Radio. Sie hörte gern die Nachrichten um acht Uhr. Die Tochter Anna kaufte einmal in der Woche eine Zeitschrift über Musik und Mode. Am Abend sahen alle zusammen das Fernsehen. Sie sahen die Tagesschau, eine Sendung mit den wichtigsten Nachrichten des Tages. Gestern kam ein Reporter ins Fernsehen und fragte die Leute auf der Straße: „Lesen Sie noch eine Zeitung oder sehen Sie nur fern?“ Viele Leute antworteten: „Wir lesen die Nachrichten heute im Internet.“ Der Vater sagte am Abend: „Früher lasen alle die Zeitung, heute sehen fast alle nur fern oder lesen im Handy.“ Anna fand das schade. Sie sagte: „Die Zeitung ist tot? Nein, ich kaufe meine Zeitschrift weiter!“ Die Familie lachte, und der Vater las am nächsten Morgen wieder seine Zeitung.
\end{textebox}

\textbf{Glossaire}

\begin{center}
\begin{tabularx}{\linewidth}{|X|X|}\hline
\rowcolor{popblueL} \textbf{Deutsch} & \textbf{Français} \\ \hline
die Zeitung, -en & le journal \\ \hline
die Schlagzeile, -n & le gros titre \\ \hline
der Sportteil, -e & la page des sports \\ \hline
das Radio, -s & la radio \\ \hline
die Nachrichten (pl.) & les informations \\ \hline
die Zeitschrift, -en & le magazine, la revue \\ \hline
das Fernsehen & la télévision \\ \hline
die Sendung, -en & l'émission \\ \hline
der Reporter, - & le reporter, le journaliste \\ \hline
die Leute (pl.) & les gens \\ \hline
das Handy, -s & le téléphone portable \\ \hline
schade & dommage \\ \hline
\end{tabularx}
\end{center}

\courssub{La démarche de compréhension : cinq étapes}

On ne répond jamais à une question de compréhension « au feeling ». Il existe une démarche rigoureuse, en cinq étapes, que tout bon élève doit suivre.

\begin{popfigure}
\begin{tikzpicture}[
  node distance=0.55cm,
  etape/.style={rectangle, rounded corners, draw=popblue, fill=popblueL, thick, align=center, text width=5.2cm, minimum height=1cm, font=\small},
  fleche/.style={-{Stealth[length=3mm]}, thick, poporange}
]
\node[etape, align=center] (e1){\textbf{1. Lecture globale}\\ Lire tout le texte une fois, sans dictionnaire.};
\node[etape, below=of e1, align=center] (e2){\textbf{2. Hypothèses}\\ Deviner le sens des mots inconnus grâce au contexte.};
\node[etape, below=of e2, align=center] (e3){\textbf{3. Lecture détaillée}\\ Relire en soulignant les informations clés.};
\node[etape, below=of e3, align=center] (e4){\textbf{4. Vérification}\\ Retrouver dans le texte la phrase qui justifie chaque réponse.};
\node[etape, below=of e4, align=center] (e5){\textbf{5. Réponse complète}\\ Rédiger une phrase : sujet + verbe conjugué en 2\textsuperscript{e} position.};
\draw[fleche] (e1) -- (e2);
\draw[fleche] (e2) -- (e3);
\draw[fleche] (e3) -- (e4);
\draw[fleche] (e4) -- (e5);
\end{tikzpicture}
\legende{Organigramme de la démarche de compréhension d'un texte}
\end{popfigure}

\begin{methode}[Répondre à une question de compréhension]
\begin{enumerate}
\item \textbf{Lire la question} et souligner le mot interrogatif (\all{wer}, \all{was}, \all{wo}, \all{wann}, \all{warum}).
\item \textbf{Repérer dans le texte} la phrase qui contient la réponse ; la souligner au crayon.
\item \textbf{Reprendre les mots de la question} dans la réponse : si la question commence par \all{Was las der Vater?}, la réponse commence par \all{Der Vater las ...}.
\item \textbf{Formuler une phrase complète} : sujet + verbe conjugué en deuxième position + compléments.
\item \textbf{Se relire} : vérifier le temps du verbe (ici le prétérit), la place du verbe et la majuscule aux noms.
\end{enumerate}
\end{methode}

\courssub{Compréhension globale}

Avant les questions de détail, l'examinateur pose toujours des questions générales. Voici les trois questions classiques et leurs réponses.

\begin{exemplebox}[Questions de compréhension globale]
\textbf{1. Worum geht es im Text?} (De quoi parle le texte ?)\par
\all{Der Text handelt von den Medien in Deutschland und von der Familie Berger.} « Le texte parle des médias en Allemagne et de la famille Berger. »\par\smallskip
\textbf{2. Wo spielt die Szene?} (Où se passe la scène ?)\par
\all{Die Szene spielt in München, in der Familie Berger.} « La scène se passe à Munich, dans la famille Berger. »\par\smallskip
\textbf{3. Wer sind die Personen?} (Qui sont les personnages ?)\par
\all{Die Personen sind der Vater, die Mutter, die Tochter Anna und ein Reporter.} « Les personnages sont le père, la mère, la fille Anna et un reporter. »
\end{exemplebox}

Observez bien : chaque réponse est une \cle{phrase complète}, avec un sujet et un verbe conjugué. Jamais un seul mot.

\courssub{Compréhension détaillée : les questions en allemand}

Passons maintenant aux questions précises. Pour chaque question, nous citons d'abord l'élément du texte, puis nous formulons la réponse complète.

\begin{exemplebox}[Questions de compréhension détaillée]
\textbf{1. Was las der Vater?} (Que lisait le père ?)\par
Élément du texte : \all{... las der Vater ... die Zeitung.}\par
Réponse : \all{Der Vater las die Zeitung.} « Le père lisait le journal. »\par\smallskip
\textbf{2. Was machte die Mutter?} (Que faisait la mère ?)\par
Élément du texte : \all{Die Mutter hörte beim Frühstück das Radio.}\par
Réponse : \all{Die Mutter hörte beim Frühstück das Radio.} « La mère écoutait la radio au petit-déjeuner. »\par\smallskip
\textbf{3. Was sahen alle am Abend?} (Que regardaient tous le soir ?)\par
Élément du texte : \all{Am Abend sahen alle zusammen das Fernsehen.}\par
Réponse : \all{Am Abend sahen alle das Fernsehen und die Tagesschau.} « Le soir, tous regardaient la télévision et le journal télévisé. »\par\smallskip
\textbf{4. Was fragte der Reporter?} (Qu'a demandé le reporter ?)\par
Élément du texte : \all{... fragte die Leute: „Lesen Sie noch eine Zeitung ...?“}\par
Réponse : \all{Der Reporter fragte die Leute, ob sie noch eine Zeitung lesen.} « Le reporter a demandé aux gens s'ils lisent encore un journal. »
\end{exemplebox}

\begin{propriete}[La question sans mot interrogatif : le verbe en première position]
En allemand, il existe deux types de questions :
\begin{itemize}
\item \textbf{Question avec mot interrogatif} (\all{was}, \all{wer}, \all{wo}...) : le verbe conjugué est en \textbf{deuxième position}, juste après le mot interrogatif.\par
\all{Was las der Vater?} « Que lisait le père ? »
\item \textbf{Question sans mot interrogatif} (réponse oui/non) : le verbe conjugué est en \textbf{première position}.\par
\all{Las der Vater die Zeitung?} « Le père lisait-il le journal ? »\par
\all{Ja, der Vater las die Zeitung.} / \all{Nein, er las kein Buch.}
\end{itemize}
C'est exactement comme le français « Lisait-il le journal ? », mais l'allemand n'a pas besoin du mot « est-ce que » : la place du verbe suffit.
\end{propriete}

\begin{center}
\begin{tabularx}{\linewidth}{|X|X|X|}\hline
\rowcolor{popblueL} \textbf{Français} & \textbf{Deutsch} & \textbf{Remarque} \\ \hline
Que lisait le père ? & \all{Was las der Vater?} & mot interrogatif + verbe en 2\textsuperscript{e} position \\ \hline
Le père lisait-il le journal ? & \all{Las der Vater die Zeitung?} & verbe en 1\textsuperscript{re} position \\ \hline
Est-ce que la mère écoutait la radio ? & \all{Hörte die Mutter das Radio?} & pas d'équivalent de « est-ce que » \\ \hline
Oui, elle l'écoutait. & \all{Ja, sie hörte das Radio.} & réponse complète exigée \\ \hline
\end{tabularx}
\end{center}

\courssub{Vrai ou faux : justifier par le texte}

L'exercice \all{Richtig oder falsch?} est un classique du baccalauréat. La règle d'or : une réponse « vrai » ou « faux » \textbf{sans justification ne rapporte aucun point}. Il faut toujours citer la phrase du texte qui prouve la réponse.

\begin{exoresolu}[Richtig oder falsch?]
Jugez les affirmations suivantes et justifiez par une phrase du texte.\par
1. \all{Der Vater las jeden Abend die Zeitung.}\par
2. \all{Die Mutter hörte gern die Nachrichten um acht Uhr.}\par
3. \all{Anna kaufte jeden Tag eine Zeitschrift.}\par
4. \all{Viele Leute lesen die Nachrichten heute im Internet.}
\tcblower
1. \textbf{Falsch.} Le texte dit : \all{Jeden Morgen kam der Vater in die Küche und las die Zeitung.} Le père lisait le journal le \emph{matin}, pas le soir.\par
2. \textbf{Richtig.} Le texte dit : \all{Sie hörte gern die Nachrichten um acht Uhr.}\par
3. \textbf{Falsch.} Le texte dit : \all{Die Tochter Anna kaufte einmal in der Woche eine Zeitschrift.} Anna achetait un magazine \emph{une fois par semaine}, pas tous les jours.\par
4. \textbf{Richtig.} Le texte dit : \all{Wir lesen die Nachrichten heute im Internet.}
\end{exoresolu}

\begin{attention}[Ne répondez jamais par un seul mot]
Au baccalauréat, une réponse comme \all{Zeitung} ou \all{Radio} est considérée comme \textbf{nulle}, même si l'information est juste. L'examinateur exige une phrase complète : sujet + verbe conjugué + complément.\par\smallskip
Mauvaise réponse : \all{Die Zeitung.}\par
Bonne réponse : \all{Der Vater las die Zeitung.}\par\smallskip
Autre piège des francophones : recopier la question sans la transformer. À la question \all{Was las der Vater?}, on ne répond pas \all{Was las der Vater die Zeitung} : on reprend les mots de la question pour construire une \textbf{phrase déclarative}, avec le verbe en deuxième position.
\end{attention}

\courssub{Recherche d'équivalents dans le texte}

Un autre exercice fréquent consiste à retrouver dans le texte l'équivalent allemand d'un mot français. La méthode : penser au contexte, puis balayer le texte à la recherche du mot.

\begin{exercice}[Trouvez dans le texte l'équivalent allemand]
Retrouvez dans le texte les mots allemands correspondant à :
\begin{enumerate}
\item le gros titre
\item l'émission
\item les informations
\item le magazine
\item les gens
\end{enumerate}
\end{exercice}

\begin{corrige}[Recherche d'équivalents]
\begin{enumerate}
\item le gros titre : \all{die Schlagzeile} (\all{Er las zuerst die Schlagzeilen.})
\item l'émission : \all{die Sendung} (\all{... eine Sendung mit den wichtigsten Nachrichten des Tages.})
\item les informations : \all{die Nachrichten} (\all{Sie hörte gern die Nachrichten um acht Uhr.})
\item le magazine : \all{die Zeitschrift} (\all{Anna kaufte einmal in der Woche eine Zeitschrift.})
\item les gens : \all{die Leute} (\all{... fragte die Leute auf der Straße.})
\end{enumerate}
\end{corrige}

\begin{aretenir}[Les mots-clés du texte à mémoriser]
\all{die Zeitung} (le journal), \all{die Schlagzeile} (le gros titre), \all{die Nachrichten} (les informations), \all{die Sendung} (l'émission), \all{das Fernsehen} (la télévision), \all{der Reporter} (le reporter), \all{die Zeitschrift} (le magazine). Attention aux articles : on dit \all{die Zeitung} mais \all{das Fernsehen} ; le genre fait partie du mot et doit être appris avec lui.
\end{aretenir}

\courssub{Question d'ouverture : à vous de parler}

La dernière question d'une épreuve de compréhension est souvent une \cle{question d'ouverture} : elle quitte le texte et vous invite à donner votre avis. C'est une préparation directe à l'expression orale et écrite guidée.

\begin{exemplebox}[Question d'ouverture et modèles de réponses]
\textbf{Und Sie? Lesen Sie die Zeitung oder sehen Sie die Nachrichten?} (Et vous, lisez-vous le journal ou regardez-vous les informations ?)\par\smallskip
Modèles de réponses au prétérit ou au présent :\par
\all{Ich sehe jeden Abend die Nachrichten im Fernsehen.} « Je regarde les informations à la télévision tous les soirs. »\par
\all{Ich las gestern eine Zeitung in Yaoundé.} « J'ai lu un journal hier à Yaoundé. »\par
\all{Ich lese die Nachrichten im Handy, das ist schnell.} « Je lis les informations sur mon portable, c'est rapide. »\par
\all{Mein Vater hörte früher das Radio, heute sieht er fern.} « Mon père écoutait la radio autrefois, aujourd'hui il regarde la télévision. »
\end{exemplebox}

\begin{experience}[Mini-débat en classe]
Par groupes de deux, posez-vous la question \all{Liest du die Zeitung oder siehst du die Nachrichten?} et répondez en deux phrases complètes. Puis rapportez la réponse de votre camarade à la classe au prétérit : \all{Aminata las gestern eine Zeitung. Moussa sah die Nachrichten im Fernsehen.}
\end{experience}

\courssub{Correction collective : rédiger les réponses au prétérit}

Pour terminer, rédigeons ensemble la correction complète, comme on le ferait au tableau. Observez la structure de chaque réponse : reprise des mots de la question, verbe au prétérit en deuxième position, phrase complète.

\begin{corrige}[Correction collective des questions de compréhension]
\textbf{1. Worum geht es im Text?}\par
\all{Der Text handelt von den Medien und von der Familie Berger in München.}\par
\textbf{2. Was las der Vater?}\par
\all{Der Vater las jeden Morgen die Zeitung.}\par
\textbf{3. Was machte die Mutter?}\par
\all{Die Mutter hörte beim Frühstück das Radio.}\par
\textbf{4. Was sahen alle am Abend?}\par
\all{Am Abend sahen alle zusammen das Fernsehen und die Tagesschau.}\par
\textbf{5. Was fragte der Reporter?}\par
\all{Der Reporter fragte die Leute, ob sie noch eine Zeitung lesen.}\par
\textbf{6. Las der Vater die Zeitung am Abend?}\par
\all{Nein, der Vater las die Zeitung am Morgen, nicht am Abend.}
\end{corrige}

\begin{aretenir}[L'essentiel de la section]
\begin{itemize}
\item On répond toujours par une \textbf{phrase complète} : sujet + verbe conjugué en deuxième position.
\item On \textbf{reprend les mots de la question} pour construire la réponse.
\item Pour un vrai/faux, on \textbf{justifie obligatoirement} par une phrase du texte.
\item Question sans mot interrogatif : le verbe passe en \textbf{première position} (\all{Las der Vater ...?}).
\item La question d'ouverture prépare l'expression personnelle : donnez votre avis en deux ou trois phrases simples.
\end{itemize}
\end{aretenir}

La compréhension du texte est maintenant assurée. Dans la section suivante, nous étudierons en détail le \cle{lexique thématique des médias}, afin d'enrichir votre vocabulaire et de préparer vos propres productions.

\coursec{Lexique thématique : les médias}

Après avoir lu et compris le texte \all{Die Medien in Deutschland}, nous avons rencontré de nombreux mots nouveaux : \all{die Zeitung}, \all{die Nachrichten}, \all{der Reporter}, \all{die Sendung}... Pour bien parler des médias en allemand, il faut maintenant organiser ce vocabulaire, l'apprendre avec ses \cle{articles} et ses \cle{pluriels}, et surtout savoir l'employer dans des phrases. C'est l'objectif de cette section : construire, champ par champ, ton lexique des médias.

\courssub{Observer avant d'apprendre : le vocabulaire dans son contexte}

Lisons d'abord ce court reportage. Souligne mentalement tous les mots qui appartiennent au monde des médias. Tu remarqueras que le texte est écrit au \cle{prétérit} : c'est le temps du récit au passé, que nous étudierons en détail dans la section de grammaire.

\begin{textebox}[Ein Tag beim Fernsehen — Un jour à la télévision]
Gestern besuchte unsere Klasse ein Fernsehstudio in der Stadt. Zuerst sahen wir die Nachrichten im Fernsehen: eine Reporterin berichtete über ein Fest in Douala, und ein Journalist las die Schlagzeilen der Zeitungen vor. Der Reporter trug ein Mikrofon und stellte viele Fragen. Die Zuschauer zu Hause sahen die Sendung um 20 Uhr. Danach sprachen wir mit einer Journalistin. Sie sagte: „Ich schreibe jeden Tag einen Artikel für die Zeitung, und am Abend informiere ich die Zuschauer im Fernsehen.“ Ein Leser schrieb einen Brief an die Zeitung, weil er eine Schlagzeile falsch fand. Am Abend sahen wir zu Hause fern: es gab eine Sendung über Fußball in Kamerun. Mein Vater las die Zeitung, meine Mutter hörte den Rundfunk, und ich sah die Nachrichten. Es war ein interessanter Tag!
\end{textebox}

\textbf{Glossaire :} \all{das Fernsehstudio} (-s) le studio de télévision ; \all{vorlesen} (las vor, hat vorgelesen) lire à voix haute ; \all{das Mikrofon} (-e) le micro ; \all{tragen} (trug, hat getragen) porter ; \all{falsch} faux ; \all{hören} écouter.

\textbf{Questions de compréhension :}
\begin{enumerate}
\item \all{Was sahen die Schüler zuerst?} (Qu'ont vu les élèves d'abord ?)
\item \all{Worüber berichtete die Reporterin?} (Sur quoi la reporter a-t-elle fait un reportage ?)
\item \all{Was macht die Journalistin jeden Tag?} (Que fait la journaliste chaque jour ?)
\item \all{Was machte der Vater am Abend?} (Que faisait le père le soir ?)
\end{enumerate}

\begin{corrige}[Questions de compréhension]
1. \all{Zuerst sahen sie die Nachrichten im Fernsehen.} « D'abord, ils ont vu les informations à la télévision. » 2. \all{Sie berichtete über ein Fest in Douala.} « Elle a fait un reportage sur une fête à Douala. » 3. \all{Sie schreibt jeden Tag einen Artikel für die Zeitung.} « Elle écrit chaque jour un article pour le journal. » 4. \all{Er las die Zeitung.} « Il lisait le journal. »
\end{corrige}

\courssub{Champ 1 — La presse écrite (die Presse)}

\begin{definition}[Vocabulaire : la presse écrite]
\begin{itemize}
\item \all{die Zeitung} (-en) : le journal — \emph{Ich lese jeden Morgen die Zeitung.} « Je lis le journal chaque matin. »
\item \all{die Zeitschrift} (-en) : la revue, le magazine — \emph{Die Zeitschrift erscheint jeden Monat.} « La revue paraît chaque mois. »
\item \all{der Artikel} (-) : l'article — \emph{Der Artikel über Kamerun war sehr lang.} « L'article sur le Cameroun était très long. »
\item \all{die Schlagzeile} (-n) : le gros titre — \emph{Die Schlagzeile stand auf der ersten Seite.} « Le gros titre se trouvait en première page. »
\item \all{die Seite} (-n) : la page — \emph{Auf Seite 3 gibt es Sportnachrichten.} « À la page 3, il y a des nouvelles sportives. »
\end{itemize}
\end{definition}

\begin{definition}[die Schlagzeile]
\all{Die Schlagzeile} (-n) désigne le \cle{gros titre} placé en haut d'un article de journal, en gros caractères, pour attirer le lecteur. Le mot est composé de \all{schlagen} (frapper) et \all{die Zeile} (la ligne) : c'est littéralement « la ligne qui frappe » le regard. Exemple : \emph{„Fußball-Team gewinnt!“ — das war die Schlagzeile.} « "L'équipe de football gagne !" — c'était le gros titre. »
\end{definition}

\courssub{Champ 2 — La télévision et la radio}

\begin{definition}[Vocabulaire : télévision et radio]
\begin{itemize}
\item \all{das Fernsehen} (nur Sg.) : la télévision (le média) — \emph{Wir sahen das Spiel im Fernsehen.} « Nous avons vu le match à la télévision. »
\item \all{der Fernseher} (-) : le poste de télévision (l'appareil) — \emph{Der Fernseher ist neu.} « Le poste de télévision est neuf. »
\item \all{die Sendung} (-en) : l'émission — \emph{Die Sendung begann um 20 Uhr.} « L'émission a commencé à 20 heures. »
\item \all{die Nachrichten} (Pl.) : les informations, le journal télévisé — \emph{Die Nachrichten kommen um 20 Uhr.} « Les informations passent à 20 heures. »
\item \all{der Rundfunk} (nur Sg.) : la radio (le média) — \emph{Mein Großvater hört gern Rundfunk.} « Mon grand-père aime écouter la radio. »
\end{itemize}
\end{definition}

\begin{attention}[das Fernsehen ou der Fernseher ?]
Les francophones confondent souvent ces deux mots. \all{Das Fernsehen} désigne le \cle{média} (la télévision en général) : \all{im Fernsehen} = « à la télévision ». \all{Der Fernseher} désigne \cle{l'appareil}, le poste : \all{Der Fernseher war kaputt.} « Le poste était en panne. » On ne dit jamais \all{im Fernseher} !
\end{attention}

\begin{attention}[die Sendung ≠ envoyer]
\all{Die Sendung} signifie « l'émission » (de radio ou de télévision). Ne la confonds pas avec le verbe \all{senden} (envoyer), même si les deux mots appartiennent à la même famille : une \all{Sendung} est à l'origine « quelque chose que l'on envoie » (dans les foyers, par les ondes). \emph{Die Sendung über Yaoundé war spannend.} « L'émission sur Yaoundé était passionnante. »
\end{attention}

\courssub{Champ 3 — Les personnes (die Personen)}

\begin{definition}[Vocabulaire : les gens des médias]
\begin{itemize}
\item \all{der Reporter} (-) / \all{die Reporterin} (-nen) : le/la reporter — \emph{Der Reporter berichtete aus Douala.} « Le reporter a fait un reportage depuis Douala. »
\item \all{der Journalist} (-en) / \all{die Journalistin} (-nen) : le/la journaliste — \emph{Die Journalistin schrieb einen Artikel.} « La journaliste a écrit un article. »
\item \all{der Zuschauer} (-) : le téléspectateur — \emph{Die Zuschauer sahen die Sendung.} « Les téléspectateurs ont vu l'émission. »
\item \all{der Leser} (-) : le lecteur — \emph{Der Leser kaufte die Zeitung.} « Le lecteur a acheté le journal. »
\end{itemize}
\end{definition}

\begin{propriete}[Le féminin des métiers avec -in]
En allemand, le féminin des noms de métiers se forme presque toujours en ajoutant le suffixe \cle{-in} au masculin ; le pluriel est alors \cle{-nen} :
\begin{center}
\begin{tabular}{|l|l|l|}
\hline
\rowcolor{popblueL} Masculin & Féminin & Pluriel féminin \\
\hline
der Journalist & die Journalistin & die Journalistinnen \\
\hline
der Reporter & die Reporterin & die Reporterinnen \\
\hline
der Leser & die Leserin & die Leserinnen \\
\hline
der Lehrer & die Lehrerin & die Lehrerinnen \\
\hline
\end{tabular}
\end{center}
Attention : \all{der Journalist} est un nom faible, il prend \cle{-en} à l'accusatif, au datif et au génitif : \emph{Ich fragte den Journalisten.} « J'ai interrogé le journaliste. »
\end{propriete}

\courssub{Champ 4 — Les actions (die Verben)}

\begin{definition}[Vocabulaire : les verbes des médias]
\begin{itemize}
\item \all{lesen} (las, hat gelesen) : lire — \emph{Er las die Zeitung.} « Il a lu le journal. »
\item \all{sehen} (sah, hat gesehen) : voir — \emph{Wir sahen den Film.} « Nous avons vu le film. »
\item \all{fernsehen} (sah fern, hat ferngesehen) : regarder la télévision — \emph{Die Kinder sahen abends fern.} « Les enfants regardaient la télé le soir. »
\item \all{berichten} (berichtete, hat berichtet) : rapporter, faire un reportage — \emph{Sie berichtete über das Fest.} « Elle a fait un reportage sur la fête. »
\item \all{informieren} (informierte, hat informiert) : informer — \emph{Die Zeitung informierte die Leser.} « Le journal a informé les lecteurs. »
\item \all{fragen / antworten} : demander / répondre — \emph{Der Reporter fragte, der Minister antwortete.} « Le reporter a posé une question, le ministre a répondu. »
\end{itemize}
\end{definition}

\begin{attention}[fernsehen : un verbe à particule séparable]
\all{fernsehen} se sépare : \emph{Ich \textbf{sehe} abends \textbf{fern}.} « Je regarde la télé le soir. » Au prétérit : \emph{Wir \textbf{sahen} gestern \textbf{fern}.} Au Perfekt : \emph{Wir haben \textbf{ferngesehen}.} La particule \all{fern-} retourne à la fin de la phrase, comme avec tous les verbes à particule.
\end{attention}

\courssub{Expressions utiles et familles de mots}

\begin{aretenir}[Expressions à retenir par cœur]
\begin{itemize}
\item \all{eine Zeitung lesen} : lire un journal — \emph{Mein Vater las eine Zeitung.} « Mon père lisait un journal. »
\item \all{im Fernsehen} : à la télévision — \emph{Das Spiel kam im Fernsehen.} « Le match passait à la télévision. »
\item \all{die Nachrichten sehen} : regarder les informations — \emph{Wir sahen die Nachrichten.} « Nous avons regardé les informations. »
\item \all{über etwas berichten} : faire un reportage sur quelque chose (avec l'accusatif) — \emph{Sie berichtete über den Markt.} « Elle a fait un reportage sur le marché. »
\end{itemize}
\end{aretenir}

\begin{propriete}[Familles de mots : un mot en cache souvent trois]
\begin{itemize}
\item \all{die Nachricht} (-en) « la nouvelle » → \all{die Nachrichten} « les informations » (au pluriel seulement dans ce sens).
\item \all{der Bericht} (-e) « le reportage » → \all{berichten} « rapporter » → \all{der Reporter} « celui qui rapporte ».
\item \all{die Schlagzeile} = \all{schlagen} (frapper) + \all{die Zeile} (la ligne).
\item \all{das Fernsehen} = \all{fern} (loin) + \all{sehen} (voir) : « voir loin », comme le grec \emph{télé-vision} !
\end{itemize}
\end{propriete}

\begin{exemplebox}[Exemples traduits]
\all{Der Reporter berichtete über das Fest.} « Le reporter a fait un reportage sur la fête. »

\all{Wir sahen gestern die Nachrichten im Fernsehen.} « Nous avons regardé les informations à la télé hier. »

\all{Die Journalistin schrieb einen Artikel für die Zeitung.} « La journaliste a écrit un article pour le journal. »

\all{Die Schlagzeile las ich zuerst.} « J'ai lu le gros titre en premier. »
\end{exemplebox}

\courssub{Carte mentale et tableau récapitulatif}

\begin{popfigure}
\begin{tikzpicture}[mindmap, grow cyclic, every node/.style={concept, text=white, align=center, font=\small}, concept color=popblue, level 1/.append style={level distance=3.2cm, sibling angle=90}, level 2/.append style={level distance=2.4cm, sibling angle=45}]
\node {die Medien}
 child[concept color=popgreen] { node {Presse}
   child { node {die Zeitung} }
   child { node {der Artikel} }
   child { node {die Schlagzeile} }
 }
 child[concept color=poporange] { node {TV und Radio}
   child { node {das Fernsehen} }
   child { node {die Sendung} }
   child { node {der Rundfunk} }
 }
 child[concept color=poppurple] { node {Personen}
   child { node {der Reporter} }
   child { node {der Journalist} }
   child { node {der Zuschauer} }
 }
 child[concept color=poppink] { node {Verben}
   child { node {lesen} }
   child { node {berichten} }
   child { node {fernsehen} }
 };
\end{tikzpicture}
\legende{Carte mentale du champ lexical des médias : quatre familles à mémoriser.}
\end{popfigure}

\begin{center}
\begin{tabularx}{\linewidth}{|l|X|X|X|}
\hline
\rowcolor{popblueL} Champ & Mot allemand (article + pluriel) & Traduction & Exemple au prétérit \\
\hline
Presse & die Zeitung (-en) & le journal & Ich las die Zeitung. \\
\hline
Presse & die Zeitschrift (-en) & la revue & Sie kaufte eine Zeitschrift. \\
\hline
Presse & der Artikel (-) & l'article & Er schrieb einen Artikel. \\
\hline
Presse & die Schlagzeile (-n) & le gros titre & Die Schlagzeile las ich zuerst. \\
\hline
Presse & die Seite (-n) & la page & Die Seite zeigte ein Foto. \\
\hline
TV/Radio & das Fernsehen & la télévision & Das Fernsehen brachte die Nachricht. \\
\hline
TV/Radio & der Fernseher (-) & le poste de télé & Der Fernseher war kaputt. \\
\hline
TV/Radio & die Sendung (-en) & l'émission & Die Sendung begann um 20 Uhr. \\
\hline
TV/Radio & die Nachrichten (Pl.) & les informations & Wir sahen die Nachrichten. \\
\hline
TV/Radio & der Rundfunk & la radio & Der Rundfunk meldete Regen. \\
\hline
Personen & der Reporter (-) / die Reporterin (-nen) & le/la reporter & Der Reporter fragte den Bürgermeister. \\
\hline
Personen & der Journalist (-en) / die Journalistin (-nen) & le/la journaliste & Die Journalistin schrieb schnell. \\
\hline
Personen & der Zuschauer (-) & le téléspectateur & Die Zuschauer sahen das Spiel. \\
\hline
Personen & der Leser (-) & le lecteur & Der Leser schrieb einen Brief. \\
\hline
Actions & lesen (las, hat gelesen) & lire & Er las die Zeitung. \\
\hline
Actions & sehen (sah, hat gesehen) & voir & Wir sahen den Film. \\
\hline
Actions & fernsehen (sah fern, hat ferngesehen) & regarder la télé & Die Kinder sahen fern. \\
\hline
Actions & berichten (berichtete, hat berichtet) & rapporter & Sie berichtete über das Fest. \\
\hline
Actions & informieren (informierte, hat informiert) & informer & Das Radio informierte uns. \\
\hline
\end{tabularx}
\end{center}

\courssub{Exercice résolu et méthode de mémorisation}

\begin{exoresolu}[Complète avec le mot qui convient]
Complète : 1. Gestern las mein Bruder die \ldots\ (journal). 2. Die \ldots\ (gros titre) war sehr groß. 3. Eine \ldots\ (journaliste) berichtete über den Fußball. 4. Wir \ldots\ (regarder la télé, prétérit) gestern Abend. 5. Die \ldots\ (informations) kamen um 20 Uhr.
\tcblower
1. \all{Zeitung} — accusatif féminin : \all{die Zeitung} ne change pas. 2. \all{Schlagzeile} — sujet au nominatif. 3. \all{Journalistin} — féminin en \all{-in}, indiqué par \all{eine}. 4. \all{sahen fern} — verbe à particule séparable au prétérit : \all{fern} va à la fin. 5. \all{Nachrichten} — toujours au pluriel dans le sens « informations ».
\end{exoresolu}

\begin{methode}[Comment mémoriser durablement ce vocabulaire]
\begin{enumerate}
\item Apprends chaque nom \cle{avec son article et son pluriel} : pas seulement \all{Zeitung}, mais \all{die Zeitung (-en)}.
\item Regroupe les mots par famille (presse / télé / personnes / verbes), comme sur la carte mentale.
\item Construis pour chaque mot \cle{une phrase au prétérit} qui te parle : \emph{Mein Vater las die Zeitung.} « Mon père lisait le journal. »
\item Récite les expressions figées comme des blocs : \all{im Fernsehen}, \all{über etwas berichten}.
\item Teste-toi en traduisant du français vers l'allemand, puis vérifie l'article et la place du verbe.
\end{enumerate}
\end{methode}

\begin{attention}[Pièges fréquents des francophones]
\begin{itemize}
\item Ne dis pas \all{die Nachricht} pour « les informations » : il faut le pluriel \all{die Nachrichten}.
\item \all{der Artikel} est masculin ; attention aussi à \all{die Seite} (féminin) et \all{das Fernsehen} (neutre) : le genre français ne suffit jamais à deviner le genre allemand.
\item « Regarder la télé » ne se traduit jamais par \all{schauen das Fernsehen} : utilise \all{fernsehen} ou \all{die Nachrichten sehen}.
\item N'oublie pas la majuscule à tous ces noms : \all{die Zeitung}, \all{der Reporter}.
\end{itemize}
\end{attention}

\begin{savaistu}[La redevance en Allemagne]
En Allemagne, chaque foyer paie une redevance, \all{der Rundfunkbeitrag}, pour financer la télévision et la radio publiques (comme ARD et ZDF). Grâce à ce système, les chaînes publiques peuvent diffuser des informations sans dépendre uniquement de la publicité. Le mot vient de \all{der Rundfunk} (la radio) et \all{der Beitrag} (la contribution) : littéralement « la contribution pour la radio » !
\end{savaistu}

\begin{exercice}[À toi de jouer]
Traduis en allemand : 1. Hier, j'ai lu un article dans le journal. 2. La reporter a fait un reportage sur le marché de Yaoundé. 3. Nous avons regardé les informations à la télévision. 4. Le gros titre était très intéressant. 5. Les téléspectateurs ont vu l'émission à 20 heures.
\end{exercice}

\begin{corrige}[Exercice « À toi de jouer »]
1. \all{Gestern las ich einen Artikel in der Zeitung.} 2. \all{Die Reporterin berichtete über den Markt in Yaoundé.} 3. \all{Wir sahen die Nachrichten im Fernsehen.} 4. \all{Die Schlagzeile war sehr interessant.} 5. \all{Die Zuschauer sahen die Sendung um 20 Uhr.}
\end{corrige}

Tu possèdes maintenant tout le lexique des médias, organisé en quatre champs et ancré par des phrases au prétérit. Nous allons à présent étudier de plus près la grammaire qui te permettra de raconter ces événements au passé : le prétérit des verbes forts.

\coursec{Grammaire : le prétérit des verbes forts}

Dans la section précédente, nous avons enrichi notre vocabulaire des médias : \all{die Zeitung}, \all{die Nachrichten}, \all{der Reporter}, \all{die Schlagzeile}, \all{das Fernsehen}, \all{die Sendung}. Mais pour raconter ce qui s'est passé — un reportage publié hier, une émission vue la semaine dernière —, il faut un temps du passé. En relisant notre texte support \emph{Die Medien in Deutschland}, nous avons rencontré des formes comme \all{las}, \all{sah}, \all{kam}, \all{schrieb}. Ce sont des formes du \cle{prétérit} (\all{das Präteritum}), le temps du récit écrit en allemand. C'est l'objet de cette section.

\courssub{Observation : repérer le prétérit dans un texte}

\begin{textebox}[Mini-reportage : Ein Abend bei Familie Ngo Bell]
Gestern Abend \textbf{war} es bei Familie Ngo Bell in Yaoundé sehr ruhig. Der Vater \textbf{las} die Zeitung und \textbf{sah} dann die Nachrichten im Fernsehen. Die Mutter \textbf{schrieb} einen Brief an ihre Schwester in Deutschland. Um 20 Uhr \textbf{kam} der Sohn von der Schule nach Hause. Er \textbf{fragte}: „Was gibt es heute in den Nachrichten?“ Der Vater \textbf{antwortete}: „Ein Reporter \textbf{sprach} über das neue Stadion in Douala.“ Die Tochter \textbf{nahm} ihr Handy und \textbf{fand} die Schlagzeilen im Internet. Die Familie \textbf{blieb} bis 22 Uhr zusammen und \textbf{sprach} über die Nachrichten des Tages. Dann \textbf{gingen} alle ins Bett.
\end{textebox}

\textbf{Glossaire :} \all{der Brief, -e} (la lettre) ; \all{das Handy, -s} (le téléphone portable) ; \all{die Schlagzeile, -n} (le gros titre) ; \all{ins Bett gehen} (aller au lit) ; \all{zusammen} (ensemble).

\textbf{Questions de compréhension :}
\begin{enumerate}
\item \all{Was las der Vater?} (Que lisait le père ?)
\item \all{Wann kam der Sohn nach Hause?} (Quand le fils est-il rentré ?)
\item \all{Wer schrieb einen Brief?} (Qui a écrit une lettre ?)
\end{enumerate}

Observons maintenant les verbes en gras de ce texte et classons-les :

\begin{center}
\begin{tabularx}{\linewidth}{|X|X|X|}
\hline
\rowcolor{popblueL} Infinitif & Forme dans le texte & Traduction \\
\hline
lesen & las & il lisait / il a lu \\
\hline
sehen & sah & il a vu / il regardait \\
\hline
kommen & kam & il est venu / il est rentré \\
\hline
schreiben & schrieb & elle a écrit \\
\hline
nehmen & nahm & elle a pris \\
\hline
finden & fand & elle a trouvé \\
\hline
bleiben & blieb & elle est restée \\
\hline
sprechen & sprach & il a parlé \\
\hline
gehen & gingen & ils sont allés \\
\hline
sein & war & il était / il a été \\
\hline
\rowcolor{poporangeL} fragen & fragte & il a demandé \\
\hline
\rowcolor{poporangeL} antworten & antwortete & il a répondu \\
\hline
\end{tabularx}
\end{center}

\textbf{Constat :} la plupart de ces verbes changent de \cle{voyelle} au passé (\all{lesen} $\rightarrow$ \all{las}, \all{sehen} $\rightarrow$ \all{sah}) : ce sont les \cle{verbes forts}. Deux verbes, \all{fragen} $\rightarrow$ \all{fragte} et \all{antworten} $\rightarrow$ \all{antwortete}, gardent leur radical et ajoutent simplement \all{-te} : ce sont des \cle{verbes faibles}. Cette opposition est la clé de toute la section.

\courssub{Règle de formation du prétérit des verbes forts}

\begin{definition}[Verbe fort, verbe faible]
Un \cle{verbe fort} (\all{starkes Verb}) forme son passé en \textbf{modifiant la voyelle de son radical} : \all{kommen} $\rightarrow$ \all{kam}. Un \cle{verbe faible} (\all{schwaches Verb}) garde son radical et ajoute \all{-te} : \all{fragen} $\rightarrow$ \all{fragte}. En allemand, les verbes forts sont aussi appelés « verbes irréguliers » ; leurs formes s'apprennent dans un tableau à trois colonnes : infinitif -- prétérit -- participe II.
\end{definition}

\begin{propriete}[Formation du prétérit fort]
Le prétérit d'un verbe fort se forme sur le \textbf{radical modifié}, c'est-à-dire la \textbf{2\up{e} colonne} du tableau des verbes irréguliers, \textbf{sans le \all{-t-} caractéristique} des verbes faibles. On ajoute ensuite les terminaisons suivantes :

\begin{center}
\begin{tabular}{|l|l|l|}
\hline
\rowcolor{popblueL} Personne & Terminaison & Exemple : lesen (radical \all{las}) \\
\hline
ich & -- (rien) & ich las \\
\hline
du & -st & du lasest \\
\hline
er / sie / es & -- (rien) & er las \\
\hline
wir & -en & wir lasen \\
\hline
ihr & -t & ihr last \\
\hline
sie / Sie & -en & sie lasen / Sie lasen \\
\hline
\end{tabular}
\end{center}

Points essentiels : la 1\up{re} et la 3\up{e} personne du singulier prennent le \textbf{radical seul, sans aucune terminaison} (\all{ich las}, \all{er las}). C'est la grande différence avec le présent (\all{er lies\emph{t}}) et avec le prétérit faible (\all{er frag\emph{te}}).
\end{propriete}

\begin{exemplebox}[Exemples traduits]
\all{Ich las gestern die Zeitung.} « J'ai lu le journal hier. »

\all{Sie kam um 20 Uhr nach Hause.} « Elle est rentrée à 20 heures. »

\all{Wir sahen die Sendung.} « Nous avons vu l'émission. »

\all{Der Reporter schrieb einen Artikel über Yaoundé.} « Le reporter a écrit un article sur Yaoundé. »

\all{Ihr spracht mit dem Journalisten.} « Vous avez parlé avec le journaliste. »
\end{exemplebox}

\courssub{Conjugaison complète : lesen et sehen au prétérit}

\begin{popfigure}
\begin{center}
\begin{tikzpicture}
\node[draw=popblue, fill=popblueL, rounded corners, align=center, font=\small] (t1) at (0,0)
{\textbf{lesen -- las -- gelesen}\\[2pt]
ich \textbf{las} \quad (j'ai lu / je lisais)\\
du \textbf{lasest} \quad (tu as lu)\\
er/sie/es \textbf{las} \quad (il/elle a lu)\\
wir \textbf{lasen} \quad (nous avons lu)\\
ihr \textbf{last} \quad (vous avez lu)\\
sie/Sie \textbf{lasen} \quad (ils/Vous ont lu)};
\node[draw=popgreen, fill=popgreenL, rounded corners, align=center, font=\small, right=1.2cm of t1] (t2)
{\textbf{sehen -- sah -- gesehen}\\[2pt]
ich \textbf{sah} \quad (j'ai vu / je voyais)\\
du \textbf{sahst} \quad (tu as vu)\\
er/sie/es \textbf{sah} \quad (il/elle a vu)\\
wir \textbf{sahen} \quad (nous avons vu)\\
ihr \textbf{saht} \quad (vous avez vu)\\
sie/Sie \textbf{sahen} \quad (ils/Vous ont vu)};
\end{tikzpicture}
\end{center}
\legende{Conjugaison de \all{lesen} et \all{sehen} au prétérit, avec traduction française.}
\end{popfigure}

Remarque : à la 2\up{e} personne du singulier, si le radical se termine déjà par un \all{-s} ou un son proche, on obtient \all{du lasest} (on garde un \all{-e-} de liaison) ; pour \all{sehen}, la forme régulière est \all{du sahst}. À l'oral, on emploie d'ailleurs rarement \all{du lasest} ; on préfère le parfait \all{du hast gelesen}.

\courssub{Les verbes forts essentiels de la leçon}

\begin{aretenir}[Tableau des verbes forts à connaître]
\begin{center}
\begin{tabularx}{\linewidth}{|X|X|X|X|}
\hline
\rowcolor{popblueL} Infinitif & Präteritum & Partizip II & Français \\
\hline
lesen & las & gelesen & lire \\
\hline
sehen & sah & gesehen & voir \\
\hline
kommen & kam & gekommen & venir \\
\hline
schreiben & schrieb & geschrieben & écrire \\
\hline
gehen & ging & gegangen & aller \\
\hline
geben & gab & gegeben & donner \\
\hline
nehmen & nahm & genommen & prendre \\
\hline
finden & fand & gefunden & trouver \\
\hline
bleiben & blieb & geblieben & rester \\
\hline
sprechen & sprach & gesprochen & parler \\
\hline
\end{tabularx}
\end{center}
\end{aretenir}

\begin{methode}[Comment apprendre les verbes forts]
\begin{enumerate}
\item Apprends toujours les \textbf{trois colonnes ensemble} : infinitif -- prétérit -- participe II (\all{gehen -- ging -- gegangen}), jamais une forme isolée.
\item Travaille par \textbf{petites séries de 5 verbes} par jour ; récite-les à voix haute comme une comptine.
\item Regroupe les verbes qui changent de voyelle de la même façon : \all{bleiben -- blieb} et \all{schreiben -- schrieb} (ei $\rightarrow$ ie) ; \all{kommen -- kam} et \all{nehmen -- nahm} ; \all{finden -- fand} et \all{sprechen -- sprach}.
\item Teste-toi en fabriquant une phrase au prétérit avec chaque verbe : \all{Gestern ging ich ins Kino.} « Hier, je suis allé au cinéma. »
\item Révise les séries précédentes avant d'en commencer une nouvelle (répétition espacée).
\end{enumerate}
\end{methode}

\courssub{Contraste : prétérit fort et prétérit faible}

\begin{center}
\begin{tabularx}{\linewidth}{|X|X|X|}
\hline
\rowcolor{popblueL} & Verbe faible : fragen & Verbe fort : lesen \\
\hline
Formation & radical + \textbf{te} & radical à voyelle changée \\
\hline
ich & frag-te & las \\
\hline
du & frag-test & las-est \\
\hline
er/sie/es & frag-te & las \\
\hline
wir & frag-ten & las-en \\
\hline
ihr & frag-tet & las-t \\
\hline
sie/Sie & frag-ten & las-en \\
\hline
\end{tabularx}
\end{center}

Le verbe faible est régulier et prévisible ; le verbe fort doit être mémorisé. Mais attention : les deux partagent la même logique de terminaisons (\all{-st}, \all{-en}, \all{-t}, \all{-en}), sauf que le fort n'a \textbf{jamais} de \all{-te} après le radical.

\courssub{Emploi du prétérit : le temps du récit écrit}

\begin{propriete}[Emploi du prétérit]
Le prétérit est le temps du \textbf{récit écrit} : articles de presse, comptes rendus, résumés, biographies, récits. C'est donc le temps naturel de notre chapitre sur les médias : \all{Der Reporter schrieb einen Artikel.} « Le reporter a écrit un article. »

Complément utile pour le Bac : \textbf{à l'oral}, les Allemands préfèrent souvent le \textbf{parfait} (\all{Perfekt}) : \all{Ich habe die Zeitung gelesen.} « J'ai lu le journal. » À l'écrit journalistique, on dira plutôt : \all{Ich las die Zeitung.} Exception importante : \all{sein} (\all{war}), \all{haben} (\all{hatte}) et les verbes modaux s'emploient au prétérit même à l'oral.
\end{propriete}

\begin{popfigure}
\begin{center}
\begin{tikzpicture}[font=\small]
\draw[-{Stealth}, thick, popdark] (-5,0) -- (5,0);
\fill[popgreen] (-3,0) circle (2.5pt);
\node[above=4pt, align=center, popgreen] at (-3,0) {\textbf{Vergangenheit}\\Präteritum\\\all{ich las}};
\fill[popblue] (0,0) circle (2.5pt);
\node[above=4pt, align=center, popblue] at (0,0) {\textbf{Gegenwart}\\Präsens\\\all{ich lese}};
\fill[poporange] (3,0) circle (2.5pt);
\node[above=4pt, align=center, poporange] at (3,0) {\textbf{Zukunft}\\Futur\\\all{ich werde lesen}};
\node[below=6pt, popmuted] at (-3,0) {passé};
\node[below=6pt, popmuted] at (0,0) {présent};
\node[below=6pt, popmuted] at (3,0) {futur};
\end{tikzpicture}
\end{center}
\legende{Frise des temps : le prétérit situe l'action dans le passé (\all{Vergangenheit}).}
\end{popfigure}

\courssub{Comparaison avec le français}

\begin{attention}[Un temps allemand, trois traductions françaises possibles]
Le prétérit allemand recouvre à la fois le \textbf{passé composé}, l'\textbf{imparfait} et le \textbf{passé simple} français. C'est le contexte qui décide :
\begin{itemize}
\item \all{Er las die Zeitung.} = « Il \textbf{a lu} le journal » (passé composé), « Il \textbf{lut} le journal » (passé simple), « Il \textbf{lisait} le journal » (imparfait).
\end{itemize}
L'allemand ne fait pas la distinction aspectuelle du français (action ponctuelle / action habituelle) avec deux temps différents : un seul temps suffit. En traduction (version), choisis le temps français le plus naturel selon le sens de la phrase.
\end{attention}

\begin{attention}[Erreur fréquente : ajouter un -t au radical fort]
Ne dis jamais ✗ \all{er laste} ni ✗ \all{er kamte}. Les verbes forts \textbf{ne prennent pas de \all{-te}} : la 3\up{e} personne du singulier est le radical seul : \all{er las}, \all{er kam}, \all{er schrieb}, \all{er ging}. Le \all{-t} de la 3\up{e} personne n'existe qu'au présent (\all{er liest}) et au prétérit faible (\all{er fragte}).
\end{attention}

\begin{attention}[Ne pas confondre sah et war]
\all{sah} est le prétérit de \all{sehen} (voir) : \all{Er sah die Sendung.} « Il a vu l'émission. » \all{war} est le prétérit de \all{sein} (être) : \all{Er war müde.} « Il était fatigué. » Les deux formes sont courtes et fréquentes dans les textes de presse ; vérifie toujours l'infinitif avant de traduire.
\end{attention}

\courssub{Exercice résolu et entraînement}

\begin{exoresolu}[Conjugaison au prétérit]
Énoncé : mets les verbes entre parenthèses au prétérit.

1. Der Journalist (schreiben) einen Artikel. \quad 2. Wir (sehen) die Nachrichten. \quad 3. Ihr (kommen) spät nach Hause. \quad 4. Ich (lesen) die Schlagzeilen. \quad 5. Die Reporter (fragen) den Minister.
\tcblower
Solution détaillée :

1. \all{Der Journalist \textbf{schrieb} einen Artikel.} Verbe fort : 2\up{e} colonne \all{schrieb}, 3\up{e} personne du singulier = radical seul, sans terminaison. « Le journaliste a écrit un article. »

2. \all{Wir \textbf{sahen} die Nachrichten.} Verbe fort : radical \all{sah} + terminaison \all{-en} à \all{wir}. « Nous avons vu les informations. »

3. \all{Ihr \textbf{kamt} spät nach Hause.} Verbe fort : radical \all{kam} + \all{-t} à \all{ihr}. « Vous êtes rentrés tard à la maison. »

4. \all{Ich \textbf{las} die Schlagzeilen.} 1\up{re} personne du singulier = radical seul. « J'ai lu les gros titres. »

5. \all{Die Reporter \textbf{fragten} den Minister.} Verbe \textbf{faible} : radical \all{frag-} + \all{-te} + \all{-n} (3\up{e} personne du pluriel). « Les reporters ont interrogé le ministre. »
\end{exoresolu}

\begin{exercice}[À toi de jouer]
1. Conjugue au prétérit : \all{geben} (er), \all{nehmen} (sie, pluriel), \all{finden} (ich), \all{bleiben} (wir), \all{sprechen} (du).

2. Traduis en allemand : a) « Hier, j'ai lu le journal. » b) « Elle a vu l'émission à 20 heures. » c) « Nous sommes allés au marché. »

3. Version : traduis en français : \all{Der Reporter kam nach Douala, sah das neue Stadion und schrieb einen langen Artikel.}
\end{exercice}

\begin{corrige}[Exercice « À toi de jouer »]
1. \all{er gab} ; \all{sie nahmen} ; \all{ich fand} ; \all{wir blieben} ; \all{du sprachst}.

2. a) \all{Gestern las ich die Zeitung.} b) \all{Sie sah die Sendung um 20 Uhr.} c) \all{Wir gingen zum Markt.}

3. « Le reporter est venu à Douala, a vu le nouveau stade et a écrit un long article. »
\end{corrige}

\begin{savaistu}[Pourquoi « fort » et « faible » ?]
La terminologie vient des grammairiens allemands du XIX\up{e} siècle (les frères Grimm, aussi célèbres pour leurs contes !). Les verbes « forts » sont les plus anciens de la langue : ils changent de voyelle selon un système hérité du vieux germanique, appelé \emph{Ablaut} (alternance vocalique). L'anglais, cousin de l'allemand, a gardé le même phénomène : \emph{sing -- sang -- sung} correspond exactement à \all{singen -- sang -- gesungen}. Les verbes « faibles », eux, sont souvent plus récents et suivent le modèle régulier en \all{-te}, comme les verbes français en \emph{-er}.
\end{savaistu}

\coursec{Grammaire : l'ordre des mots}

Dans la section précédente, nous avons appris à conjuguer les verbes forts au prétérit : \all{lesen -- las}, \all{sehen -- sah}, \all{kommen -- kam}. Mais savoir conjuguer ne suffit pas pour construire une phrase allemande correcte. Il faut encore savoir \emph{où placer} chaque élément. Reprenons deux phrases de notre texte support :

\all{Der Vater las die Zeitung.} « Le père lisait le journal. »

\all{Gestern las der Vater die Zeitung.} « Hier, le père lisait le journal. »

Observez bien : dans la deuxième phrase, le sujet \all{der Vater} n'est plus avant le verbe, mais \emph{après}. Ce n'est pas une fantaisie : c'est la conséquence de la règle la plus importante de la syntaxe allemande, celle que nous allons étudier maintenant.

\courssub{La règle d'or : le verbe conjugué en 2e position}

\begin{propriete}[La règle fondamentale V2]
Dans la phrase déclarative allemande, le \cle{verbe conjugué} occupe \textbf{toujours la 2e position}, quoi que l'on mette en 1re position (sujet, complément de temps, complément de lieu, etc.). La 1re position ne peut contenir qu'\textbf{un seul élément} (un groupe de mots qui forme un bloc : \all{der Vater}, \all{gestern Abend}, \all{am Morgen}).
\end{propriete}

\begin{popfigure}
\begin{tikzpicture}[
  box/.style={draw=popblue, fill=popblueL, rounded corners, align=center, minimum height=1.2cm, font=\small},
  vbox/.style={draw=poporange, fill=poporangeL, rounded corners, align=center, minimum height=1.2cm, font=\small\bfseries},
  fleche/.style={-{Stealth[length=3mm]}, thick, popdark}
]
\node[box, align=center] (p1) at (0,0){\textbf{Position 1}\\ sujet OU complément\\ \all{Der Vater} / \all{Gestern}};
\node[vbox, align=center] (p2) at (4.6,0){\textbf{Position 2}\\ VERBE conjugué\\ \all{las}};
\node[box, align=center] (p3) at (9.2,0){\textbf{Milieu}\\ sujet / compléments\\ \all{der Vater} / \all{die Zeitung}};
\node[box, align=center] (p4) at (13.4,0){\textbf{Fin}\\ participe, infinitif\\ ou \all{nicht}};
\draw[fleche] (p1) -- (p2);
\draw[fleche] (p2) -- (p3);
\draw[fleche] (p3) -- (p4);
\end{tikzpicture}
\legende{Les positions dans la phrase déclarative allemande : le verbe conjugué est un pivot fixe en 2e position.}
\end{popfigure}

\courssub{Cas 1 : le sujet en 1re position (ordre de base)}

C'est le cas le plus simple, identique au français : sujet -- verbe -- complément.

\begin{exemplebox}[Sujet -- verbe -- COD]
\all{Der Vater las die Zeitung.} « Le père lisait le journal. »

\all{Die Mutter sah die Nachrichten.} « La mère regardait les informations. »

\all{Der Reporter schrieb einen Artikel.} « Le journaliste écrivait un article. »
\end{exemplebox}

Ici, la position 1 est occupée par le sujet, la position 2 par le verbe conjugué : tout est en ordre.

\courssub{Cas 2 : l'inversion — un complément en 1re position}

Si l'on veut insister sur un complément (de temps, de lieu), on le place en 1re position. Mais alors, le verbe doit \emph{rester} en 2e position : le sujet passe donc \emph{après} le verbe. C'est ce qu'on appelle l'\cle{inversion}.

\begin{exemplebox}[Complément -- verbe -- sujet]
\all{Die Mutter las die Nachrichten.} « La mère lisait les informations. »

→ \all{Am Morgen las die Mutter die Nachrichten.} « Le matin, la mère lisait les informations. »

\all{Gestern las der Vater die Zeitung.} « Hier, le père lisait le journal. »

\all{Früher kaufte die Familie die Zeitung am Kiosk.} « Autrefois, la famille achetait le journal au kiosque. »

\all{In Yaoundé sahen viele Menschen abends die Nachrichten.} « À Yaoundé, beaucoup de gens regardaient les informations le soir. »
\end{exemplebox}

\begin{attention}[L'erreur n° 1 des francophones]
En français, l'ordre sujet--verbe ne change jamais : « Hier, le père lisait le journal. » Le francophone a donc tendance à écrire : ✗ \all{Gestern der Vater las die Zeitung.} C'est \textbf{faux} : il y a alors deux éléments avant le verbe. La seule phrase correcte est : ✓ \all{Gestern las der Vater die Zeitung.} Retenez : ce qui est en 1re position « pousse » le sujet derrière le verbe.
\end{attention}

\courssub{Cas 3 : l'ordre des compléments (TeKaMoLo simplifié)}

Quand plusieurs compléments se suivent au milieu de la phrase, l'allemand impose un ordre précis, résumé par le mot mnémotechnique \cle{TeKaMoLo} :

\begin{center}
\begin{tabularx}{\linewidth}{|l|X|X|}
\hline
\rowcolor{popblueL} \textbf{Ordre} & \textbf{Type de complément} & \textbf{Exemple} \\
\hline
\textbf{Te} -- Temporal & temps (quand ?) & \all{gestern Abend} \\
\hline
\textbf{Ka} -- Kausal & cause (pourquoi ?) & \all{wegen des Wetters} \\
\hline
\textbf{Mo} -- Modal & manière (comment ?) & \all{mit Interesse} \\
\hline
\textbf{Lo} -- Lokal & lieu (où ?) & \all{zu Hause} \\
\hline
\end{tabularx}
\end{center}

Au niveau de la Première, retenez surtout : \textbf{le temps avant le lieu}.

\begin{exemplebox}[Temps avant lieu]
\all{Wir sahen gestern Abend zu Hause die Nachrichten.} « Nous avons regardé les informations hier soir à la maison. »

\all{Die Schüler lasen am Morgen in der Schule die Zeitung.} « Les élèves lisaient le journal le matin à l'école. »
\end{exemplebox}

\courssub{Cas 4 : la place de la négation \all{nicht}}

Deux cas à distinguer :

\begin{itemize}
\item Si l'on nie \textbf{toute la phrase}, \all{nicht} se place en fin de phrase (ou avant le dernier bloc) : \all{Er las die Zeitung nicht.} « Il ne lisait pas le journal. »
\item Si l'on nie \textbf{un seul élément}, \all{nicht} se place juste \emph{avant} cet élément : \all{Er las nicht die Zeitung, sondern die Zeitschrift.} « Il ne lisait pas le journal, mais le magazine. » / \all{Er las die Zeitung nicht am Morgen.} « Il ne lisait pas le journal le matin. »
\end{itemize}

\courssub{Cas 5 : la phrase interrogative}

\begin{propriete}[Verbe en 1re position ou après le mot interrogatif]
\begin{itemize}
\item \textbf{Question fermée} (réponse oui/non) : le verbe conjugué est en \textbf{1re position} : \all{Las der Vater die Zeitung?} « Le père lisait-il le journal ? »
\item \textbf{Question ouverte} : le verbe est en \textbf{2e position}, juste après le mot interrogatif : \all{Was las der Vater?} « Que lisait le père ? » / \all{Wann sahen wir die Nachrichten?} « Quand regardions-nous les informations ? »
\end{itemize}
\end{propriete}

\courssub{Comparaison français / allemand}

\begin{center}
\begin{tabularx}{\linewidth}{|X|X|X|}
\hline
\rowcolor{popblueL} \textbf{Français} & \textbf{Deutsch} & \textbf{Remarque} \\
\hline
Le père lisait le journal. & \all{Der Vater las die Zeitung.} & ordre identique \\
\hline
Hier, le père lisait le journal. & \all{Gestern las der Vater die Zeitung.} & inversion obligatoire en allemand \\
\hline
Que lisait le père ? & \all{Was las der Vater?} & verbe après le mot interrogatif \\
\hline
Il ne lisait pas le journal. & \all{Er las die Zeitung nicht.} & \all{nicht} en fin de phrase \\
\hline
\end{tabularx}
\end{center}

\begin{savaistu}[La règle V2, une marque germanique]
La règle de la 2e position verbale (dite « V2 ») est une caractéristique typique des langues germaniques : on la retrouve en néerlandais, en suédois, en norvégien. Elle explique la fameuse « inversion » qui surprend tant les francophones. Le français, lui, a perdu cette règle au cours de son histoire : c'est pourquoi votre cerveau de francophone résiste... mais avec de l'entraînement, l'inversion devient un réflexe !
\end{savaistu}

\begin{methode}[Vérifier l'ordre des mots en 3 étapes]
\begin{enumerate}
\item \textbf{Repérer le verbe conjugué} dans la phrase.
\item \textbf{Compter les éléments avant le verbe} : il ne doit y en avoir qu'\textbf{un seul} (un bloc). Si deux éléments précèdent le verbe, la phrase est fausse.
\item \textbf{Vérifier le milieu et la fin} : temps avant lieu (TeKaMoLo), \all{nicht} avant l'élément nié ou en fin de phrase.
\end{enumerate}
Exemple de vérification : ✗ \all{Gestern der Vater las die Zeitung} → deux éléments avant \all{las} (\all{Gestern} + \all{der Vater}) → faux. Correction : \all{Gestern las der Vater die Zeitung.}
\end{methode}

\begin{textebox}[Ein Abend bei Familie Nkoulou]
Gestern Abend saß die Familie Nkoulou im Wohnzimmer. Der Vater las die Zeitung, und die Mutter sah die Nachrichten im Fernsehen. Um 20 Uhr begann die Tagesschau. Der Reporter zeigte Bilder aus Douala: dort regnete es den ganzen Tag. „Interessante Nachrichten!“, sagte der Vater. Die Tochter Amina las nicht die Zeitung, sie schrieb einen Brief an ihre Freundin in Deutschland. Früher hatte die Familie keinen Fernseher; damals hörten alle das Radio. Am Abend erzählte der Großvater Geschichten. Heute informiert sich die Familie mit dem Handy, dem Fernsehen und der Zeitung. Am Morgen kaufte der Vater die Zeitung am Kiosk neben der Schule. Der Verkäufer lächelte und sagte: „Die Schlagzeile heute ist sehr gut!“ Zu Hause las der Vater zuerst die Sportseite, denn er liebt Fußball. Am Wochenende spielt der berühmte Club aus Yaoundé gegen ein Team aus Deutschland. Die ganze Familie wird das Spiel im Fernsehen sehen.

\textbf{Glossaire} : das Wohnzimmer, - (le salon) -- die Tagesschau (le journal télévisé allemand) -- es regnete (il pleuvait) -- der Brief, -e (la lettre) -- damals (à cette époque) -- das Radio, -s (la radio) -- sich informieren (s'informer) -- das Handy, -s (le téléphone portable) -- der Kiosk, -e (le kiosque) -- der Verkäufer, - (le vendeur) -- die Schlagzeile, -n (le gros titre) -- zuerst (d'abord) -- die Sportseite, -n (la page des sports) -- das Spiel, -e (le match).

\textbf{Questions} : 1. Wer las die Zeitung? 2. Was sah die Mutter im Fernsehen? 3. Was machte Amina? 4. Wie informierte sich die Familie früher?
\end{textebox}

\begin{exoresolu}[Transformation : commencer par un complément]
\textbf{Énoncé} : Recopiez chaque phrase en commençant par le complément entre parenthèses. Attention à l'inversion !

a) Der Vater las die Zeitung. (\all{Gestern}) \quad b) Wir sahen die Nachrichten. (\all{Am Abend}) \quad c) Die Familie hörte das Radio. (\all{Früher})
\tcblower
\textbf{Solution détaillée} : le complément passe en position 1, le verbe reste en position 2, le sujet passe après le verbe.

a) \all{Gestern las der Vater die Zeitung.} « Hier, le père lisait le journal. »

b) \all{Am Abend sahen wir die Nachrichten.} « Le soir, nous regardions les informations. »

c) \all{Früher hörte die Familie das Radio.} « Autrefois, la famille écoutait la radio. »
\end{exoresolu}

\begin{exercice}[À vous de jouer]
1. Remettez les phrases dans l'ordre correct : a) \all{las / am Morgen / die Mutter / die Zeitung} ; b) \all{nicht / der Schüler / las / den Artikel} ; c) \all{sahen / gestern Abend / wir / zu Hause / die Nachrichten}.

2. Traduisez en allemand : a) « Hier, le reporter écrivit un article. » b) « Le matin, la famille regardait les informations. » c) « Le père ne lisait pas le journal. »

3. Transformez : \all{Die Schüler lasen die Schlagzeile.} → commencez par \all{In der Schule}.
\end{exercice}

\begin{corrige}[Exercice « À vous de jouer »]
1. a) \all{Am Morgen las die Mutter die Zeitung.} b) \all{Der Schüler las den Artikel nicht.} c) \all{Wir sahen gestern Abend zu Hause die Nachrichten.} (temps avant lieu)

2. a) \all{Gestern schrieb der Reporter einen Artikel.} b) \all{Am Morgen sah die Familie die Nachrichten.} c) \all{Der Vater las die Zeitung nicht.}

3. \all{In der Schule lasen die Schüler die Schlagzeile.} (inversion : le sujet \all{die Schüler} passe après le verbe \all{lasen})
\end{corrige}

\begin{aretenir}[L'essentiel de l'ordre des mots]
\begin{itemize}
\item Phrase déclarative : le verbe conjugué est \textbf{toujours en 2e position}.
\item Un complément en 1re position provoque l'\textbf{inversion} : le sujet passe après le verbe.
\item Ordre des compléments : \textbf{temps avant lieu} (TeKaMoLo).
\item \all{nicht} : en fin de phrase (négation globale) ou devant l'élément nié.
\item Question fermée : verbe en 1re position ; question ouverte : verbe après le mot interrogatif.
\end{itemize}
\end{aretenir}

\coursec{Méthode du commentaire et production guidée}

Dans les sections précédentes, nous avons étudié le vocabulaire des médias, le prétérit des verbes forts et l'ordre des mots. Ces outils grammaticaux ne sont pas une fin en soi : ils servent maintenant à \emph{parler} et à \emph{écrire} sur le texte \all{Die Medien in Deutschland}. Au baccalauréat comme en classe, on vous demandera de \cle{commenter} un texte, de le \cle{résumer} et de produire un court paragraphe personnel. Cette section vous donne une méthode simple, en quatre étapes, des modèles de phrases prêts à l'emploi et une synthèse des points de grammaire indispensables.

\begin{definition}[Vocabulaire clé de la leçon : Les médias]
Voici les six mots de base du programme, avec article, pluriel et traduction. Apprenez-les par cœur avec leur genre.
\begin{center}
\begin{tabular}{|l|l|l|}
\hline
\rowcolor{popblueL} \textbf{Allemand (Singulier)} & \textbf{Pluriel} & \textbf{Français} \\
\hline
die Zeitung & die Zeitungen & le journal (quotidien) \\
\hline
die Nachrichten & (Pl. only) & les informations, les nouvelles \\
\hline
der Reporter & die Reporter & le reporter, le journaliste \\
\hline
die Schlagzeile & die Schlagzeilen & le titre (de une), la manchette \\
\hline
das Fernsehen & — & la télévision (médium) \\
\hline
die Sendung & die Sendungen & l'émission, le journal télévisé \\
\hline
\end{tabular}
\end{center}
\emph{Composés utiles :} \all{die Tagesschau} (le journal télévisé de 20h, ARD), \all{das Handy, -s} (le portable), \all{das Internet} (Internet), \all{der Artikel, -e} (l'article), \all{der Bericht, -e} (le reportage).
\end{definition}

\begin{propriete}[Rappel : Le prétérit des verbes forts (Präteritum)]
Les verbes forts changent leur voyelle radicale au prétérit (souvent sans \emph{-te}). Pas de terminaison aux 1\textsuperscript{re} et 3\textsuperscript{e} pers. sing. Voici les formes rencontrées dans cette leçon :
\begin{center}
\begin{tabular}{|l|l|l|l|l|l|l|}
\hline
\rowcolor{popblueL} \textbf{Infinitif} & \textbf{ich} & \textbf{du} & \textbf{er/sie/es} & \textbf{wir} & \textbf{ihr} & \textbf{sie/Sie} \\
\hline
lesen (lire) & las & last & las & lasen & last & lasen \\
\hline
sehen (voir) & sah & sahst & sah & sahen & saht & sahen \\
\hline
gehen (aller) & ging & gingst & ging & gingen & gingt & gingen \\
\hline
kommen (venir) & kam & kamst & kam & kamen & kamt & kamen \\
\hline
finden (trouver) & fand & fandst & fand & fanden & fandet & fanden \\
\hline
geben (donner) & gab & gabst & gab & gaben & gabt & gaben \\
\hline
nehmen (prendre) & nahm & nahmst & nahm & nahmen & nahmt & nahmen \\
\hline
sprechen (parler) & sprach & sprachst & sprach & sprachen & spracht & sprachen \\
\hline
schreiben (écrire) & schrieb & schriebst & schrieb & schrieben & schriebt & schrieben \\
\hline
trinken (boire) & trank & trankst & trank & tranken & trankt & tranken \\
\hline
\end{tabular}
\end{center}
\emph{Verbes irréguliers (modaux / auxiliaires) :} \all{sein $\to$ war/ware}, \all{haben $\to$ hatte}, \all{werden $\to$ wurde}, \all{können $\to$ konnte}, \all{müssen $\to$ musste}, \all{wollen $\to$ wollte}.
\end{propriete}

\begin{propriete}[Synthèse : L'ordre des mots dans la phrase]
Trois positions de base à maîtriser pour le bac :
\begin{center}
\begin{tabularx}{\linewidth}{|X|X|X|}
\hline
\rowcolor{popblueL} \textbf{Type de phrase} & \textbf{Position du verbe conjugué} & \textbf{Exemple} \\
\hline
Phrase principale (déclarative) & \cle{2\textsuperscript{e} position} & \all{Gestern \textbf{sah} ich die Tagesschau.} \\
\hline
Phrase principale (question fermée) & \cle{1\textsuperscript{re} position} & \all{\textbf{Sah} ich die Tagesschau?} \\
\hline
Subordonnée (\all{dass, weil, wenn, als...}) & \cle{Fin de phrase} & \all{..., \textbf{dass} ich die Tagesschau \textbf{sah}.} \\
\hline
\end{tabularx}
\end{center}
\emph{Règle d'or :} En phrase principale, \textbf{un seul élément} (sujet, complément de temps, \all{Meiner Meinung nach}, \all{Zuerst}...) peut précéder le verbe (position 1). Le verbe reste \emph{toujours} en position 2.
\end{propriete}

\courssub{La méthode du commentaire en quatre étapes}

\begin{methode}[Comment rédiger un commentaire de texte]
Un commentaire de texte se construit toujours dans le même ordre :
\begin{enumerate}
\item \textbf{Présenter le texte} : indiquer le titre, le thème, la source (journal, dialogue, article) et les personnages. Une ou deux phrases suffisent.
\item \textbf{Dégager l'idée principale} : dire en une phrase de quoi parle vraiment le texte (ici : l'évolution des habitudes médiatiques).
\item \textbf{Analyser le contenu et le vocabulaire} : appuyer son analyse sur des \cle{citations courtes} du texte, entre guillemets allemands „…“, et relever les mots importants du champ lexical des médias.
\item \textbf{Donner son avis personnel} : terminer par une opinion justifiée, introduite par \all{Meiner Meinung nach...} ou \all{Ich finde...}.
\end{enumerate}
Règle d'or : jamais d'avis sans justification, jamais d'analyse sans citation du texte.
\end{methode}

\begin{popfigure}
\begin{center}
\begin{tikzpicture}[
node distance=0.55cm,
etape/.style={rectangle, rounded corners=4pt, draw=popblue, fill=popblueL, text width=5.6cm, align=center, font=\small, inner sep=6pt},
etapeLast/.style={rectangle, rounded corners=4pt, draw=poporange, fill=poporangeL, text width=5.6cm, align=center, font=\small, inner sep=6pt},
fleche/.style={-{Stealth[length=3mm]}, thick, popdark}]
\node[etape, align=center] (e1){\textbf{1. Présentation}\\ Titre, thème, source, personnages};
\node[etape, below=of e1, align=center] (e2){\textbf{2. Idée principale}\\ De quoi parle le texte ?};
\node[etape, below=of e2, align=center] (e3){\textbf{3. Analyse}\\ Contenu + citations courtes „...“\\ + vocabulaire des médias};
\node[etapeLast, below=of e3, align=center] (e4){\textbf{4. Avis personnel}\\ \all{Meiner Meinung nach...} + justification};
\draw[fleche] (e1) -- (e2);
\draw[fleche] (e2) -- (e3);
\draw[fleche] (e3) -- (e4);
\end{tikzpicture}
\end{center}
\legende{Organigramme du commentaire de texte : quatre étapes dans un ordre fixe.}
\end{popfigure}

\courssub{Les phrases-types du commentaire}

Pour chaque étape, il existe des formules toutes faites. Apprenez-les par cœur : elles garantissent un allemand correct et une structure claire.

\begin{exemplebox}[Phrases-types pour le commentaire]
\begin{itemize}
\item \all{Der Text handelt von einer Familie, die Zeitung las und fernsah.} « Le texte traite d'une famille qui lisait le journal et regardait la télévision. »
\item \all{Der Text stammt aus einer Zeitung.} « Le texte est tiré d'un journal. »
\item \all{Der Autor zeigt, dass früher fast jeder eine Zeitung kaufte.} « L'auteur montre qu'autrefois presque tout le monde achetait un journal. »
\item \all{Ein Beispiel dafür ist: „die Mutter las die Nachrichten auf dem Handy“.} « Un exemple de cela est : "la mère lisait les informations sur le portable". »
\item \all{Meiner Meinung nach sind die Nachrichten im Fernsehen wichtig.} « À mon avis, les informations télévisées sont importantes. »
\item \all{Ich finde die Zeitung interessanter, weil sie mehr Details gibt.} « Je trouve le journal plus intéressant parce qu'il donne plus de détails. »
\end{itemize}
\end{exemplebox}

\begin{attention}[Attention à l'ordre des mots dans les formules]
Après \all{Meiner Meinung nach} (« à mon avis »), le verbe conjugué reste en \cle{deuxième position} : \all{Meiner Meinung nach \textbf{ist} das Fernsehen praktisch.} et non \emph{Meiner Meinung nach das Fernsehen ist...}. Après \all{dass} et \all{weil}, en revanche, le verbe va à la \cle{fin} : \all{Der Autor zeigt, dass viele das Internet \textbf{benutzen}.}
\end{attention}

\begin{center}
\begin{tabularx}{\linewidth}{|X|X|X|}
\hline
\rowcolor{popblueL} Étape & Formule allemande & Français \\
\hline
Présentation & Der Text handelt von... & Le texte traite de... \\
\hline
Source & Der Text stammt aus... & Le texte est tiré de... \\
\hline
Idée principale & Der Autor zeigt, dass... & L'auteur montre que... \\
\hline
Analyse & Ein Beispiel ist: „...“ & Un exemple est : « ... » \\
\hline
Avis personnel & Meiner Meinung nach... & À mon avis... \\
\hline
Justification & ..., weil... & ..., parce que... \\
\hline
\end{tabularx}
\end{center}

\courssub{Commentaire modèle rédigé}

Voici un commentaire complet, rédigé selon la méthode en quatre étapes. Observez d'abord le texte, puis repérez les quatre étapes et les formules utilisées.

\begin{textebox}[Commentaire modèle : Die Medien in Deutschland]
Der Text „Die Medien in Deutschland“ handelt von Familie Müller und ihren Medien. Der Vater las die Zeitung, die Mutter las die Nachrichten auf dem Handy, und am Abend sahen alle die Tagesschau. Der Text zeigt, dass früher fast jeder eine Zeitung kaufte, heute aber viele das Internet benutzen. Meiner Meinung nach ist das Fernsehen praktisch, aber die Zeitung gibt mehr Details.

\medskip
\emph{Glossaire :} die Medien (Pl.) « les médias » ; die Nachrichten (Pl.) « les informations » ; das Handy, -s « le portable » ; die Tagesschau « le journal télévisé (allemand) » ; früher « autrefois » ; benutzen « utiliser » ; praktisch « pratique » ; das Detail, -s « le détail ».
\end{textebox}

Analysons ce modèle. La première phrase \emph{présente} le texte (titre, thème, personnages) avec la formule \all{Der Text handelt von...}. Les deux phrases suivantes \emph{analysent} le contenu : elles citent les faits du texte (le père lisait le journal, la mère lisait les informations sur son portable, toute la famille regardait la Tagesschau le soir). La phrase avec \all{Der Text zeigt, dass...} dégage \emph{l'idée principale} : l'évolution des habitudes médiatiques, du journal papier vers Internet. Enfin, la dernière phrase donne un \emph{avis personnel justifié} : la télévision est pratique, mais le journal donne plus de détails. Remarquez que tout le récit des faits est au \cle{prétérit} (\all{las}, \all{sahen}, \all{kaufte}) et que le verbe est toujours à la bonne place.

\begin{aretenir}[Les quatre étapes en quatre formules]
\begin{enumerate}
\item Présentation : \all{Der Text handelt von...}
\item Idée principale : \all{Der Autor zeigt, dass...}
\item Analyse : \all{Ein Beispiel ist: „...“}
\item Avis personnel : \all{Meiner Meinung nach...}
\end{enumerate}
\end{aretenir}

\courssub{Production guidée 1 : le résumé au prétérit}

Résumer, c'est redire l'essentiel \emph{avec ses propres mots}, en quatre ou cinq phrases. Pour enchaîner les idées, on utilise des \cle{mots-charnières} : \all{zuerst} « d'abord », \all{dann} « puis », \all{danach} « ensuite », \all{schließlich} « finalement ». Attention : ces mots occupent la première position, donc le verbe suit immédiatement (inversion) : \all{Zuerst \textbf{las} der Vater die Zeitung.}

\begin{exemplebox}[Modèle de résumé]
\all{Zuerst las der Vater die Zeitung. Dann las die Mutter die Nachrichten auf dem Handy. Danach sahen alle am Abend die Tagesschau. Schließlich zeigt der Text: Heute benutzen viele Leute das Internet.}

« D'abord le père lisait le journal. Puis la mère lisait les informations sur le portable. Ensuite tous regardaient la Tagesschau le soir. Finalement, le texte montre : aujourd'hui, beaucoup de gens utilisent Internet. »
\end{exemplebox}

\begin{attention}[Ne recopiez pas le texte !]
Dans un résumé, il est interdit de recopier des phrases entières du texte support. Il faut \cle{reformuler} avec ses propres mots, au prétérit, en gardant seulement l'information essentielle. Par exemple, au lieu de recopier toute la description de la famille Müller, on écrit simplement : \all{Die Familie benutzte verschiedene Medien.} « La famille utilisait différents médias. »
\end{attention}

\courssub{Production guidée 2 : le paragraphe personnel}

On vous demandera souvent de raconter ce que \emph{vous} avez lu ou vu aux informations. Trois amorces suffisent :

\begin{exemplebox}[Modèle de paragraphe personnel]
\all{Gestern sah ich die Nachrichten im Fernsehen. Ich sah einen Bericht über Fußball in Kamerun. Dann las ich eine Zeitung. Ich las einen Artikel über das Wetter in Yaoundé. Die Nachrichten waren interessant.}

« Hier, j'ai regardé les informations à la télévision. J'ai vu un reportage sur le football au Cameroun. Puis j'ai lu un journal. J'ai lu un article sur la météo à Yaoundé. Les informations étaient intéressantes. »
\end{exemplebox}

Notez les formes de prétérit : \all{ich sah} (sehen), \all{ich las} (lesen), \all{waren} (sein). Le verbe reste en deuxième position, même après \all{Gestern} et \all{Dann}.

\courssub{Expression orale guidée : le mini-dialogue en binôme}

\begin{experience}[Hast du gestern die Nachrichten gesehen?]
Par binômes, réalisez le dialogue suivant, puis inversez les rôles. Remplacez les éléments soulignés par vos propres informations (un match de football, un article sur Douala, la météo, un reportage sur le cacao...).

\medskip
\all{— Hast du gestern die Nachrichten gesehen?}\\
\all{— Ja, ich sah einen Bericht über Fußball. Und du?}\\
\all{— Nein, aber ich las einen Artikel in der Zeitung. Er war sehr interessant.}

« — As-tu regardé les informations hier ? — Oui, j'ai vu un reportage sur le football. Et toi ? — Non, mais j'ai lu un article dans le journal. Il était très intéressant. »

\medskip
\emph{Note pédagogique :} À l'oral authentique, les Allemands utilisent le \emph{Perfekt} (\all{ich habe gesehen}). Ici, nous nous entraînons au \emph{Prétérit} (\all{ich sah}) car c'est la forme écrite exigée au programme de Première et au Baccalauréat.
\end{experience}

\begin{exoresolu}[Rédiger un mini-commentaire]
Énoncé : Rédigez un commentaire de 4 phrases sur le texte \all{Die Medien in Deutschland} en utilisant les quatre formules : \all{Der Text handelt von...}, \all{Der Autor zeigt, dass...}, \all{Ein Beispiel ist...}, \all{Meiner Meinung nach...}.
\tcblower
Solution détaillée :
\begin{enumerate}
\item Présentation : \all{Der Text handelt von Familie Müller und ihren Medien.} (Le texte traite de la famille Müller et de ses médias.)
\item Idée principale : \all{Der Autor zeigt, dass die Familie früher die Zeitung las und heute das Handy benutzt.} (L'auteur montre que la famille lisait autrefois le journal et utilise aujourd'hui le portable.) Attention : après \all{dass}, le verbe \all{benutzt} va à la fin.
\item Analyse avec citation : \all{Ein Beispiel ist: „am Abend sahen alle die Tagesschau“.} (Un exemple est : « le soir, tous regardaient la Tagesschau ».)
\item Avis personnel : \all{Meiner Meinung nach ist das Handy praktisch, weil es schnell ist.} (À mon avis, le portable est pratique parce qu'il est rapide.) Après \all{weil}, le verbe \all{ist} va à la fin.
\end{enumerate}
Vérification : quatre phrases, quatre formules, verbe en deuxième position (sauf subordonnées), prétérit correct (\all{las}, \all{sahen}).
\end{exoresolu}

\courssub{Grille d'auto-évaluation}

Avant de rendre votre production (résumé, paragraphe ou commentaire), vérifiez chaque point de cette grille :

\begin{center}
\begin{tabularx}{\linewidth}{|X|X|}
\hline
\rowcolor{popblueL} Critère & Question à se poser \\
\hline
Vocabulaire des médias & Ai-je utilisé au moins 2 mots du thème (\all{die Zeitung}, \all{die Nachrichten}, \all{das Fernsehen}, \all{die Sendung}) ? \\
\hline
Prétérit correct & Mes verbes forts sont-ils bien conjugués (\all{las}, \all{sah}, \all{ging}, \all{kam}) ? \\
\hline
Verbe en 2\textsuperscript{e} position & Après \all{Gestern}, \all{Dann}, \all{Meiner Meinung nach}, le verbe suit-il immédiatement ? \\
\hline
Phrase complète & Chaque phrase a-t-elle un sujet, un verbe et un point final ? \\
\hline
Reformulation & Ai-je évité de recopier des phrases entières du texte ? \\
\hline
\end{tabularx}
\end{center}

\begin{savaistu}[Le commentaire au baccalauréat]
Au baccalauréat camerounais, le commentaire et le résumé sont notés sur \emph{deux} critères à la fois : la \cle{pertinence du contenu} (avez-vous compris le texte ? vos exemples sont-ils exacts ?) \emph{et} la \cle{correction grammaticale} (prétérit, ordre des mots, accords). Un commentaire plein de bonnes idées mais avec des verbes mal placés perd des points ; inversement, un allemand parfait qui ne répond pas à la question aussi. Travaillez donc les deux : la méthode en quatre étapes pour le contenu, et la grammaire des sections précédentes pour la forme.
\end{savaistu}

\begin{exercice}[À vous de produire]
\begin{enumerate}
\item Résumez le texte \all{Die Medien in Deutschland} en 4 phrases avec \all{zuerst}, \all{dann}, \all{danach}, \all{schließlich}.
\item Rédigez un paragraphe personnel de 4 phrases : \all{Gestern sah ich... / Ich las... / Meiner Meinung nach...}
\item En binôme, jouez le mini-dialogue sur les informations d'hier et présentez-le à la classe.
\end{enumerate}
\end{exercice}

\begin{corrige}[Exercice : propositions de réponses]
\begin{enumerate}
\item \all{Zuerst las der Vater die Zeitung. Dann las die Mutter die Nachrichten auf dem Handy. Danach sahen alle die Tagesschau. Schließlich benutzen heute viele das Internet.} (Réponses variables : vérifier le prétérit et l'inversion après les mots-charnières.)
\item Exemple : \all{Gestern sah ich die Nachrichten im Fernsehen. Ich las auch einen Artikel über Douala. Dann sprach ich mit meinem Bruder über die Nachrichten. Meiner Meinung nach sind die Nachrichten wichtig.}
\item Dialogue libre : exiger au moins deux répliques par élève, le prétérit (\all{sah}, \all{las}) et une question correcte avec \all{Hast du...?}.
\end{enumerate}
\end{corrige}

\newpage

\coursec{Activité d'intégration}

\coursec{Activité d’intégration : Notre mini-journal télévisé}

\courssub{Objectifs de l'activité}

À la fin de cette activité d'intégration, tu seras capable de :

\begin{itemize}
\item mobiliser le \cle{vocabulaire des médias} (die Zeitung, die Nachrichten, der Reporter, die Schlagzeile, das Fernsehen, die Sendung, berichten) dans une situation de communication authentique ;
\item rédiger de courtes nouvelles au \cle{prétérit des verbes forts} (kam, gab, sagte, berichtete, interviewte) ;
\item respecter l'\cle{ordre des mots} allemand : verbe conjugué en deuxième position, sujet et compléments bien placés ;
\item présenter oralement un mini-journal télévisé avec une prononciation claire ;
\item comprendre la production des autres groupes et répondre à des questions de compréhension.
\end{itemize}

\courssub{Situation}

Le club d'allemand de ton lycée à Yaoundé prépare la „Woche der deutschen Sprache“ (la semaine de la langue allemande). Pour clôturer la semaine, chaque classe de Première doit présenter un mini-journal télévisé en allemand devant les autres classes. Ta classe a été choisie pour ouvrir la manifestation ! Le professeur d'allemand vous demande de préparer, par groupes de quatre, une „Sendung“ de trois nouvelles : une nouvelle de l'école, une nouvelle de la ville et une nouvelle du Cameroun. La meilleure émission sera affichée au tableau du cours d'allemand et enregistrée pour le site du club.

Voici le modèle d'annonce que le club a publié :

\begin{textebox}[Ankündigung des Deutsch-Clubs]
Liebe Schülerinnen und Schüler!

Am Freitag ist unsere „Woche der deutschen Sprache“ zu Ende. Jede Klasse präsentiert eine kleine Nachrichtensendung auf Deutsch. Ein Nachrichtensprecher meldet drei Nachrichten: eine Nachricht aus der Schule, eine aus der Stadt und eine aus Kamerun. Ein Reporter interviewt einen Zeugen. Jede Nachricht hat drei oder vier Sätze im Präteritum. Vergesst nicht: das Verb steht auf Position zwei!

Die beste Sendung kommt an die Wandzeitung des Deutschkurses. Viel Erfolg!

Euer Deutsch-Club
\end{textebox}

\begin{definition}[Glossaire de l'annonce]
\begin{itemize}
\item die Ankündigung, -en : l'annonce
\item der Nachrichtensprecher, - : le présentateur du journal télévisé
\item der Zeuge, -n / die Zeugin, -nen : le témoin / la témoin
\item die Wandzeitung, -en : le journal mural (affiché au tableau)
\item vergessen (vergaß, hat vergessen) : oublier
\end{itemize}
\end{definition}

\courssub{Consignes}

Travaillez par groupes de quatre. Répartissez les rôles : un \all{Nachrichtensprecher} (présentateur), un \all{Reporter}, un \all{Zeuge} (témoin) et un \all{Redakteur} (rédacteur en chef, qui vérifie la grammaire de tous les textes).

\begin{enumerate}
\item \textbf{Choix des nouvelles.} En groupe, choisissez trois nouvelles réelles ou imaginaires de la semaine : une nouvelle de l'école (par exemple un match de football, une visite, un concours), une nouvelle de la ville (Yaoundé ou Douala : un marché, la circulation, un événement culturel) et une nouvelle du Cameroun (sport, cacao, musique). Notez pour chaque nouvelle une \all{Schlagzeile} (un titre accrocheur).

\item \textbf{Rédaction au prétérit.} Rédigez chaque nouvelle en trois ou quatre phrases. Utilisez au moins trois verbes forts au prétérit (par exemple : kam, gab, war, hatte, sagte, begann, gewann, fand ... statt) et au moins trois mots du vocabulaire de la leçon (die Nachrichten, berichten, die Schlagzeile, der Reporter, die Sendung, das Fernsehen). Vérifiez que le verbe conjugué est toujours en deuxième position.

\item \textbf{Préparation de l'interview.} Le reporter prépare deux questions au prétérit (par exemple : \all{Was sahen Sie?} — Qu'avez-vous vu ?) et le témoin prépare deux réponses courtes, également au prétérit.

\item \textbf{Relecture collective.} Le rédacteur en chef relit tous les textes : prétérit correct ? verbe en deuxième position ? majuscules aux noms ? vocabulaire de la leçon ? Corrigez ensemble les erreurs.

\item \textbf{Présentation orale.} Chaque groupe présente sa „Sendung“ devant la classe (environ trois minutes). Le présentateur annonce les nouvelles, le reporter interviewe le témoin. Parlez lentement et clairement.

\item \textbf{Compréhension orale.} Après chaque émission, les autres groupes répondent à deux questions de compréhension posées par le professeur ou par le groupe (par exemple : \all{Was passierte in der Schule? Wer interviewte den Zeugen?}).

\item \textbf{Vote et affichage.} La classe vote pour la meilleure „Sendung“ (critères : prétérit correct, ordre des mots, vocabulaire, prononciation). Le texte gagnant est affiché au tableau du cours d'allemand.
\end{enumerate}

\begin{attention}[Pièges à éviter dans votre production]
\begin{itemize}
\item Ne dites pas \all{Gestern die Mannschaft gewann das Spiel} : le verbe doit être en deuxième position : \all{Gestern gewann die Mannschaft das Spiel.}
\item Attention au prétérit des verbes forts : on écrit \all{gewann} (et non „gewinnte“), \all{kam} (et non „kamte“), \all{gab} (et non „gabte“).
\item N'oubliez pas la majuscule à tous les noms : \all{die Nachrichten, das Spiel, der Reporter}.
\item Faux ami : \all{die Nachricht} signifie „la nouvelle / l'information“, pas „la nouvelle“ littéraire (qui se dit \all{die Novelle} ou \all{die Erzählung}).
\end{itemize}
\end{attention}

\courssub{Ressources à mobiliser}

\begin{aretenir}[Boîte à outils de la Sendung]
\textbf{Vocabulaire :} die Nachrichten (les informations), die Schlagzeile, -n (le gros titre), der Reporter, - / die Reporterin, -nen, die Sendung, -en (l'émission), das Fernsehen (la télévision), die Zeitung, -en (le journal), berichten (rapporter), interviewen (interviewer), passieren (se passer).

\textbf{Prétérit des verbes forts utiles :} sein $\rightarrow$ war ; haben $\rightarrow$ hatte ; kommen $\rightarrow$ kam ; geben $\rightarrow$ gab ; gewinnen $\rightarrow$ gewann ; beginnen $\rightarrow$ begann ; sehen $\rightarrow$ sah ; sprechen $\rightarrow$ sprach ; stattfinden $\rightarrow$ fand ... statt (\all{das Spiel fand um 15 Uhr statt} : le match a eu lieu à 15 heures).

\textbf{Ordre des mots :} phrase déclarative = verbe conjugué en position 2. Si un complément de temps ouvre la phrase, le sujet passe après le verbe : \all{Am Montag kam ein Reporter in unsere Schule.}

\textbf{Structures de présentation :} \all{Guten Tag, liebe Zuschauer! Hier sind die Nachrichten.} (Bonjour chers téléspectateurs ! Voici les informations.) ; \all{Unser Reporter berichtet von...} (Notre reporter fait un reportage sur...) ; \all{Das war unsere Sendung. Auf Wiedersehen!} (C'était notre émission. Au revoir !)
\end{aretenir}

\courssub{Proposition de production et corrigé commenté}

\begin{exoresolu}[Modèle de mini-journal télévisé : „Klasse Erste A meldet sich“]
Rédige avec ton groupe une Sendung de trois nouvelles (école, ville, Cameroun) avec une interview, en respectant les consignes de l'activité.

\tcblower

\textbf{Production modèle (texte du présentateur) :}

\all{Guten Tag, liebe Zuschauer! Hier sind die Nachrichten vom Freitag.}

\all{Schlagzeile eins: Fußballfest an unserer Schule. Am Mittwoch gewann die Fußballmannschaft unseres Gymnasiums das Spiel gegen das Lyzeum in Nkolndongo. Das Spiel begann um 15 Uhr, und viele Schüler kamen auf den Sportplatz. Unser Kapitän hatte zwei Tore und war der beste Spieler.}

\all{Schlagzeile zwei: Neuer Markt in Yaoundé. In der Stadt gab es am Donnerstag ein Fest: Der neue Markt öffnete seine Türen. Viele Händler kamen mit Obst, Gemüse und Stoffen. Unser Reporter berichtet jetzt von dort.}

\textbf{Interview (reporter et témoin) :}

\all{Reporter: Guten Tag! Was sahen Sie am Donnerstag auf dem Markt?}

\all{Zeuge: Ich sah viele Menschen und viele Stände. Eine Frau verkaufte frische Mangos.}

\all{Reporter: Wie fanden Sie das Fest?}

\all{Zeuge: Es war sehr schön. Die Musik begann um 10 Uhr, und alle waren froh.}

\all{Sprecher: Schlagzeile drei: Kamerun im Fernsehen. Am Freitag sprach das Fernsehen über den Kakao aus Kamerun. Viele Bauern aus der Zentralregion kamen nach Yaoundé und zeigten ihre Produkte.}

\all{Das war unsere Sendung. Auf Wiedersehen!}

\textbf{Traduction :} « Bonjour chers téléspectateurs ! Voici les informations de vendredi. Titre un : fête du football dans notre école. Mercredi, l'équipe de football de notre lycée a gagné le match contre le lycée de Nkolndongo. Le match a commencé à 15 heures et beaucoup d'élèves sont venus sur le terrain de sport. Notre capitaine a marqué deux buts et a été le meilleur joueur. Titre deux : nouveau marché à Yaoundé. En ville, il y a eu jeudi une fête : le nouveau marché a ouvert ses portes. Beaucoup de commerçants sont venus avec des fruits, des légumes et des tissus. Notre reporter fait maintenant un reportage de là-bas. — Reporter : Bonjour ! Qu'avez-vous vu jeudi au marché ? — Témoin : J'ai vu beaucoup de gens et beaucoup d'étals. Une femme vendait des mangues fraîches. — Reporter : Comment avez-vous trouvé la fête ? — Témoin : C'était très beau. La musique a commencé à 10 heures et tout le monde était content. — Titre trois : le Cameroun à la télévision. Vendredi, la télévision a parlé du cacao du Cameroun. Beaucoup de paysans de la région du Centre sont venus à Yaoundé et ont montré leurs produits. C'était notre émission. Au revoir ! »

\textbf{Commentaire des choix de langue :}

\begin{itemize}
\item \textbf{Prétérit des verbes forts :} \all{gewann} (gewinnen), \all{begann} (beginnen), \all{kamen} (kommen), \all{gab} (geben, ici dans la structure impersonnelle \all{es gab} : „il y a eu“), \all{sah} (sehen), \all{sprach} (sprechen), \all{war} (sein), \all{hatte} (haben), \all{fand ... statt} (stattfinden). Toutes les formes sont correctes et sans terminaison ajoutée au radical fort.
\item \textbf{Ordre des mots :} quand la phrase commence par un complément de temps, le sujet se place après le verbe : \all{Am Mittwoch gewann die Fußballmannschaft...} ; \all{In der Stadt gab es...} ; \all{Am Freitag sprach das Fernsehen...}. Le verbe reste toujours en deuxième position.
\item \textbf{Vocabulaire de la leçon :} \all{die Nachrichten, die Schlagzeile, der Reporter, berichten, das Fernsehen, die Sendung} sont tous employés en contexte.
\item \textbf{Questions au prétérit :} \all{Was sahen Sie?} (verbe en première position, sujet après le verbe) et \all{Wie fanden Sie das Fest?} (mot interrogatif + verbe + sujet).
\item \textbf{Registre adapté :} le présentateur salue le public (\all{liebe Zuschauer}) et conclut avec une formule de clôture, comme dans un vrai journal télévisé.
\end{itemize}
\end{exoresolu}

\courssub{Bilan de l'activité}

Cette activité t'a permis de réinvestir l'ensemble des acquis de la leçon dans une tâche finale concrète et motivante :

\begin{itemize}
\item \textbf{Compréhension :} tu as lu une annonce authentique du club d'allemand et écouté les émissions des autres groupes pour répondre à des questions.
\item \textbf{Lexique :} tu as employé le vocabulaire des médias dans un contexte réel de communication (présenter, rapporter, interviewer).
\item \textbf{Grammaire et conjugaison :} tu as rédigé des phrases au prétérit des verbes forts et vérifié la place du verbe en deuxième position, y compris après un complément de temps en tête de phrase.
\item \textbf{Production orale :} tu as pris la parole devant la classe dans un rôle défini, avec une prononciation soignée.
\end{itemize}

\begin{methode}[Grille d'auto-évaluation]
Avant de présenter ta Sendung, vérifie chaque point :
\begin{enumerate}
\item Ai-je au moins trois verbes forts au prétérit, correctement conjugués ?
\item Le verbe conjugué est-il en deuxième position dans chaque phrase déclarative ?
\item Ai-je utilisé au moins trois mots du vocabulaire des médias ?
\item Tous les noms ont-ils une majuscule ?
\item Mon texte a-t-il une ouverture (\all{Hier sind die Nachrichten}) et une clôture (\all{Auf Wiedersehen}) ?
\end{enumerate}
Si tu peux répondre „oui“ aux cinq questions, ta Sendung est prête pour le tableau du cours d'allemand !
\end{methode}

\newpage

\coursec{Méthodes et exercices résolus}

\courssub{Fiche méthode 1 : Comprendre un texte sur les médias (global puis détaillé)}

\begin{methode}[Lire et exploiter un texte de presse]
\begin{enumerate}
\item \textbf{Première lecture (globale).} Lis le titre et le premier paragraphe : ils donnent le thème. Demande-toi : \emph{Worum geht es ?} (De quoi s'agit-il ?). Repère les mots déjà connus : \all{die Zeitung}, \all{das Fernsehen}, \all{die Nachrichten}.
\item \textbf{Repérage des mots-clés.} Souligne les noms (ils ont une majuscule !), les chiffres et les verbes au prétérit (\all{las}, \all{sah}, \all{berichtete}) : ils portent l'information.
\item \textbf{Deuxième lecture (détaillée).} Pour chaque question, localise la phrase du texte qui contient la réponse. Ne traduis pas tout le texte : cherche l'information demandée.
\item \textbf{Deviner les mots inconnus.} Utilise le contexte et les familles de mots : \all{die Schlagzeile} contient \all{die Zeile} (la ligne) ; \all{der Reporter} ressemble au français « reporter ».
\item \textbf{Rédiger la réponse.} Réponds par une phrase complète en allemand, en reprenant les mots de la question, avec le verbe conjugué en deuxième position.
\end{enumerate}
\end{methode}

\begin{exoresolu}[Compréhension de l'écrit — type Bac]
\textbf{Texte :} \all{Gestern Abend sahen viele Deutsche die Nachrichten im Fernsehen. Ein Reporter berichtete über ein Fußballspiel in München. Die Zeitungen brachten am nächsten Morgen große Schlagzeilen: „Bayern gewinnt 3:1“. Viele Leser kauften die Zeitung am Kiosk.}

\textbf{Glossaire :} \all{der Leser, -} = le lecteur ; \all{der Kiosk, -e} = le kiosque ; \all{bringen} (prétérit \all{brachte}) = apporter, publier.

\textbf{Questions :} 1. \all{Was sahen viele Deutsche gestern Abend?} 2. \all{Worüber berichtete der Reporter?} 3. \all{Wo kauften die Leser die Zeitung?}
\tcblower
\textbf{Raisonnement :} La question 1 porte sur le contenu de la phrase 1 (\all{sahen} = prétérit de \all{sehen}). La question 2 renvoie à la phrase 2 (\all{berichtete über} + accusatif : « faire un reportage sur »). La question 3 se trouve dans la dernière phrase (\all{am Kiosk} = \all{an dem Kiosk}, datif de lieu).

\textbf{Réponses rédigées :}
\begin{enumerate}
\item \all{Viele Deutsche sahen gestern Abend die Nachrichten im Fernsehen.} « Beaucoup d'Allemands ont regardé les informations à la télévision hier soir. »
\item \all{Der Reporter berichtete über ein Fußballspiel in München.} « Le reporter a fait un reportage sur un match de football à Munich. »
\item \all{Die Leser kauften die Zeitung am Kiosk.} « Les lecteurs ont acheté le journal au kiosque. »
\end{enumerate}
\textbf{Conclusion :} chaque réponse reprend le sujet et le verbe de la question, conjugués au prétérit, avec le verbe en deuxième position.
\end{exoresolu}

\courssub{Fiche méthode 2 : Employer le vocabulaire des médias}

\begin{methode}[Bien utiliser le lexique de la presse et de la télévision]
\begin{enumerate}
\item \textbf{Apprends chaque nom avec son article et son pluriel :} \all{die Zeitung, -en} ; \all{der Reporter, -} ; \all{die Schlagzeile, -n} ; \all{die Sendung, -en} ; \all{das Fernsehen} (sans pluriel) ; \all{die Nachricht, -en} (souvent au pluriel : \all{die Nachrichten} = les informations télévisées).
\item \textbf{Associe le nom à son verbe :} \all{eine Zeitung lesen} (lire un journal), \all{Nachrichten sehen} (regarder les infos), \all{über etwas berichten} (faire un reportage sur quelque chose), \all{eine Sendung anschauen} (regarder une émission).
\item \textbf{Construis des familles de mots :} \all{die Nachricht} $\to$ \all{der Nachrichtensprecher, -} (le présentateur du journal télévisé) ; \all{die Zeit} $\to$ \all{die Zeitung} ; \all{schreiben} $\to$ \all{der Schreiber} $\to$ \all{der Journalist, -en}.
\item \textbf{Vérifie le cas :} après \all{lesen}, \all{sehen}, \all{kaufen}, le nom est à l'accusatif : \all{Ich lese die Zeitung.} (et non un datif).
\end{enumerate}
\end{methode}

\begin{center}
\begin{tabularx}{\linewidth}{|X|X|X|}\hline
\rowcolor{popblueL} Français & Deutsch & Remarque \\ \hline
le journal & die Zeitung, -en & féminin \\ \hline
les informations (TV) & die Nachrichten (Pl.) & pluriel de \all{die Nachricht} \\ \hline
le reporter, le journaliste & der Reporter, - & masculin \\ \hline
le gros titre & die Schlagzeile, -n & \all{schlagen} + \all{die Zeile} \\ \hline
la télévision (le média) & das Fernsehen & sans pluriel \\ \hline
l'émission & die Sendung, -en & de \all{senden} (émettre) \\ \hline
le téléspectateur & der Zuschauer, - & de \all{zuschauen} \\ \hline
le kiosque & der Kiosk, -e & \all{am Kiosk} = au kiosque \\ \hline
\end{tabularx}
\end{center}

\begin{exoresolu}[Vocabulaire — phrases à compléter]
\textbf{Complète avec le mot qui convient :} \all{Schlagzeile, Reporter, Sendung, Nachrichten, Zeitung}

1. \all{Der \_\_\_\_\_\_\_\_ interviewt den Fußballtrainer.} 2. \all{Mein Vater liest jeden Morgen die \_\_\_\_\_\_\_\_.} 3. \all{Um 20 Uhr kommen die \_\_\_\_\_\_\_\_ im Fernsehen.} 4. \all{Die \_\_\_\_\_\_\_\_ „Sieg für Kamerun!“ steht auf der ersten Seite.} 5. \all{Diese \_\_\_\_\_\_\_\_ über Musik sehe ich gern.}
\tcblower
\textbf{Raisonnement :} on identifie la fonction de chaque mot. 1. Celui qui interviewe, c'est le journaliste : \all{Reporter}. 2. Ce qu'on lit le matin : \all{Zeitung} (accusatif féminin, l'article \all{die} ne change pas). 3. À 20 h à la télévision : les informations, \all{Nachrichten} (pluriel, d'où \all{kommen}). 4. Le gros titre en première page : \all{Schlagzeile}. 5. Une émission sur la musique : \all{Sendung}.

\textbf{Solutions :} 1. \all{Der Reporter interviewt den Fußballtrainer.} 2. \all{Mein Vater liest jeden Morgen die Zeitung.} 3. \all{Um 20 Uhr kommen die Nachrichten im Fernsehen.} 4. \all{Die Schlagzeile „Sieg für Kamerun!“ steht auf der ersten Seite.} 5. \all{Diese Sendung über Musik sehe ich gern.}

\textbf{Conclusion :} bien vérifier l'article (der/die/das), le nombre (singulier ou pluriel) et le sens du contexte avant de choisir le mot.
\end{exoresolu}

\courssub{Fiche méthode 3 : Conjuguer le prétérit des verbes forts}

\begin{propriete}[Formation du Präteritum des verbes forts]
Le prétérit des verbes forts se forme sur un \cle{radical modifié} (2\textsuperscript{e} colonne de la liste des verbes irréguliers), auquel on ajoute les terminaisons :

\begin{center}
\begin{tabular}{|l|l|l|}\hline
\rowcolor{popblueL} Personne & Terminaison & Exemple : \all{lesen} (\all{las}) \\ \hline
ich & -- & ich las \\ \hline
du & -st & du lasest \\ \hline
er/sie/es & -- & er las \\ \hline
wir & -en & wir lasen \\ \hline
ihr & -t & ihr last \\ \hline
sie/Sie & -en & sie lasen \\ \hline
\end{tabular}
\end{center}

Radicaux fréquents à connaître : \all{sehen $\to$ sah}, \all{kommen $\to$ kam}, \all{geben $\to$ gab}, \all{schreiben $\to$ schrieb}, \all{gehen $\to$ ging}, \all{finden $\to$ fand}, \all{bleiben $\to$ blieb}. Auxiliaires irréguliers : \all{sein $\to$ war}, \all{haben $\to$ hatte}, \all{werden $\to$ wurde}. Verbes mixtes (radical changé + terminaisons faibles) : \all{bringen $\to$ brachte}, \all{denken $\to$ dachte}, \all{kennen $\to$ kannte}.
\end{propriete}

\begin{methode}[Employer le Präteritum pas à pas]
\begin{enumerate}
\item \textbf{Repère le radical du prétérit} dans la liste des verbes irréguliers (2\textsuperscript{e} colonne) : \all{lesen $\to$ las}, \all{sehen $\to$ sah}, \all{schreiben $\to$ schrieb}.
\item \textbf{Ajoute les terminaisons} selon le tableau ci-dessus.
\item \textbf{Attention :} pas de terminaison à la 1\textsuperscript{re} et à la 3\textsuperscript{e} personne du singulier : \all{er kam} (jamais \emph{er kamte}).
\item \textbf{Emploi :} le prétérit s'emploie surtout à l'écrit (récits, articles de presse) : \all{Der Reporter schrieb einen Artikel.} À l'oral, on préfère souvent le Perfekt : \all{Der Reporter hat einen Artikel geschrieben.}
\item \textbf{Les auxiliaires} \all{war}, \all{hatte}, \all{wurde} sont les formes les plus fréquentes : apprends-les par cœur.
\end{enumerate}
\end{methode}

\begin{exoresolu}[Conjugaison — mettre au prétérit]
\textbf{Mets les verbes entre parenthèses au prétérit :}

1. \all{Die Zeitung (bringen) große Fotos.} 2. \all{Wir (sehen) die Sendung am Abend.} 3. \all{Der Reporter (schreiben) einen Artikel über Douala.} 4. \all{Ihr (lesen) die Schlagzeilen.} 5. \all{Es (sein) eine interessante Nachricht.}
\tcblower
\textbf{Raisonnement :} tous ces verbes sont forts, mixtes ou irréguliers : on prend le radical de la 2\textsuperscript{e} colonne, puis les terminaisons du prétérit. \all{bringen} est un verbe mixte : \all{brachte} (radical changé + terminaison faible -te).

\textbf{Solutions :}
\begin{enumerate}
\item \all{Die Zeitung brachte große Fotos.} « Le journal a publié de grandes photos. » (verbe mixte : \all{brachte})
\item \all{Wir sahen die Sendung am Abend.} « Nous avons vu l'émission le soir. » (\all{sehen $\to$ sah} + -en)
\item \all{Der Reporter schrieb einen Artikel über Douala.} « Le reporter a écrit un article sur Douala. » (\all{schreiben $\to$ schrieb}, 3\textsuperscript{e} pers. sing. sans terminaison)
\item \all{Ihr last die Schlagzeilen.} « Vous avez lu les gros titres. » (\all{lesen $\to$ las} + -t)
\item \all{Es war eine interessante Nachricht.} « C'était une information intéressante. » (\all{sein $\to$ war})
\end{enumerate}
\textbf{Conclusion :} la seule difficulté est de connaître le radical irrégulier ; les terminaisons sont régulières.
\end{exoresolu}

\courssub{Fiche méthode 4 : Respecter l'ordre des mots}

\begin{propriete}[La place du verbe conjugué]
\begin{itemize}
\item \textbf{Phrase déclarative :} verbe conjugué en \cle{2\textsuperscript{e} position}. Le sujet n'est pas obligé d'être premier : \all{Gestern las ich die Zeitung.} (complément en tête $\to$ sujet après le verbe : c'est l'\cle{inversion}).
\item \textbf{Question sans mot interrogatif :} verbe en \cle{1\textsuperscript{re} position} : \all{Siehst du die Nachrichten?}
\item \textbf{Question avec mot interrogatif :} verbe en \cle{2\textsuperscript{e} position} : \all{Was liest du?}
\item \textbf{Subordonnée} (introduite par \all{dass}, \all{weil}, \all{obwohl}...) : verbe conjugué \cle{à la fin} : \all{Ich weiß, dass er die Zeitung las.}
\item \textbf{Au Perfekt :} auxiliaire en 2\textsuperscript{e} position, participe passé à la fin : \all{Er hat die Sendung gesehen.}
\end{itemize}
\end{propriete}

\begin{methode}[Construire une phrase correcte]
\begin{enumerate}
\item \textbf{Choisis ce que tu mets en tête de phrase} (sujet ou complément de temps/lieu).
\item \textbf{Place le verbe conjugué juste après}, en 2\textsuperscript{e} position, quoi qu'il arrive.
\item \textbf{Place le sujet} : avant le verbe s'il est en tête, juste après le verbe sinon.
\item \textbf{Termine par les autres éléments} (compléments, particule séparable, participe passé ou infinitif à la fin).
\item \textbf{Vérifie la subordonnée :} après \all{weil} ou \all{dass}, le verbe conjugué va à la toute fin.
\end{enumerate}
\end{methode}

\begin{exoresolu}[Grammaire — remettre la phrase en ordre]
\textbf{Remets les mots dans l'ordre correct :}

1. \all{gestern / las / die Zeitung / mein Bruder} 2. \all{die Nachrichten / sah / am Abend / sie} 3. \all{weil / interessant / der Artikel / war / ich / las / ihn}
\tcblower
\textbf{Raisonnement :} 1. Si \all{gestern} ouvre la phrase, le verbe reste en 2\textsuperscript{e} position et le sujet passe après (inversion). 2. Même principe : \all{am Abend} peut ouvrir la phrase, le sujet \all{sie} suit alors le verbe. 3. \all{weil} introduit une subordonnée : le verbe conjugué \all{war} va à la fin de la subordonnée.

\textbf{Solutions :}
\begin{enumerate}
\item \all{Gestern las mein Bruder die Zeitung.} « Hier, mon frère a lu le journal. » (ou : \all{Mein Bruder las gestern die Zeitung.})
\item \all{Am Abend sah sie die Nachrichten.} « Le soir, elle a regardé les informations. » (ou : \all{Sie sah am Abend die Nachrichten.})
\item \all{Ich las ihn, weil der Artikel interessant war.} « Je l'ai lu parce que l'article était intéressant. »
\end{enumerate}
\textbf{Conclusion :} dans tous les cas, le verbe conjugué est en 2\textsuperscript{e} position dans la principale, à la fin dans la subordonnée.
\end{exoresolu}

\courssub{Fiche méthode 5 : Rédiger un court résumé ou paragraphe}

\begin{methode}[Produire un court texte sur les médias]
\begin{enumerate}
\item \textbf{Releve 4 à 5 informations essentielles} du texte (qui ? quoi ? où ? quand ?).
\item \textbf{Rédige des phrases simples} au prétérit ou au présent, une idée par phrase, verbe en 2\textsuperscript{e} position.
\item \textbf{Relie les phrases} avec \all{und}, \all{aber}, \all{dann}, \all{weil}.
\item \textbf{Utilise le vocabulaire du thème :} \all{die Zeitung berichtet}, \all{der Reporter schreibt}, \all{die Zuschauer sehen die Sendung}.
\item \textbf{Relis-toi :} majuscules aux noms, terminaisons des verbes, place du verbe.
\end{enumerate}
\end{methode}

\begin{exoresolu}[Expression écrite — résumé guidé]
\textbf{Écris un court résumé (4-5 phrases) du texte suivant :} \all{Am Sonntag las Familie Mbarga die Zeitung. Der Vater las die Sportseite, die Mutter las die Nachrichten aus Kamerun. Am Abend sahen alle die Sendung „Afrika heute“ im Fernsehen. Ein Reporter berichtete über den Markt in Yaoundé. Die Familie fand die Sendung sehr interessant.}
\tcblower
\textbf{Raisonnement :} on garde les informations essentielles (la famille, le journal, la télévision, le reportage), on reformule avec des connecteurs, on conserve le prétérit.

\textbf{Réponse rédigée :}

\all{Am Sonntag las Familie Mbarga die Zeitung. Der Vater interessierte sich für Sport, und die Mutter las die Nachrichten aus Kamerun. Am Abend sahen alle zusammen eine Sendung im Fernsehen. Ein Reporter berichtete über den Markt in Yaoundé. Die Familie fand die Sendung sehr interessant.}

« Dimanche, la famille Mbarga a lu le journal. Le père s'intéressait au sport et la mère a lu les nouvelles du Cameroun. Le soir, tous ont regardé ensemble une émission à la télévision. Un reporter a fait un reportage sur le marché de Yaoundé. La famille a trouvé l'émission très intéressante. »

\textbf{Conclusion :} 5 phrases, verbe en 2\textsuperscript{e} position, prétérit correct (\all{las}, \all{sahen}, \all{berichtete}, \all{fand}), vocabulaire du thème employé.
\end{exoresolu}

\begin{attention}[Les 5 erreurs qui coûtent des points au Bac]
\begin{enumerate}
\item \textbf{Oublier la majuscule aux noms :} \all{die zeitung} est faux ; écris toujours \all{die Zeitung}, \all{das Fernsehen}, \all{die Schlagzeile}.
\item \textbf{Mauvaise place du verbe :} après un complément en tête de phrase, le sujet vient après le verbe : \all{Gestern sah ich die Nachrichten.} (jamais \emph{Gestern ich sah...}).
\item \textbf{Terminaison fautive au prétérit des verbes forts :} \all{er sah} et non \emph{er sahte} ; la 3\textsuperscript{e} personne du singulier n'a pas de terminaison.
\item \textbf{Faux amis et confusions lexicales :} \all{die Nachricht} signifie « la nouvelle, l'information » ; \all{der Reporter} est le journaliste, alors que \all{der Bericht} est le compte rendu, le reportage écrit ; \all{die Sendung} est l'émission, pas « l'envoi » au sens de colis dans ce contexte.
\item \textbf{Calque du français :} « regarder la télévision » se dit \all{fernsehen} (verbe à particule : \all{Ich sehe abends fern.}) ou \all{das Fernsehen schauen}, jamais \emph{regardieren} ; « lire le journal » : \all{die Zeitung lesen}, avec l'accusatif.
\end{enumerate}
\end{attention}

\newpage

\coursec{Exercices}

\courssub{Je vérifie mes connaissances}

\begin{exercice}[QCM : les médias]
Entoure la bonne réponse.
\begin{enumerate}
\item \all{Die Zeitung} signifie : a) la télévision \quad b) le journal \quad c) la radio
\item \all{Der Reporter} est : a) celui qui écrit et enquête \quad b) celui qui regarde la télévision \quad c) celui qui vend des journaux
\item \all{Die Schlagzeile} est : a) le titre en gros caractères \quad b) la dernière page \quad c) la publicité
\item \all{Die Nachrichten} sont : a) les jeux télévisés \quad b) les informations \quad c) les films
\item Le pluriel de \all{die Zeitung} est : a) die Zeitungen \quad b) die Zeitunge \quad c) die Zeitunger
\end{enumerate}
\end{exercice}

\begin{exercice}[Vrai ou faux ?]
Réponds par \all{richtig} ou \all{falsch} et justifie ta réponse par une phrase en allemand.
\begin{enumerate}
\item \all{Das Fernsehen} bedeutet „die Zeitung“.
\item Eine \all{Sendung} sieht man im Fernsehen oder hört man im Radio.
\item \all{Der Reporter schreibt Artikel für die Zeitung.}
\item \all{Die Nachrichten} kommen nur im Radio.
\end{enumerate}
\end{exercice}

\begin{exercice}[Vocabulaire : les familles de mots]
Complète le tableau avec les mots de la leçon.
\begin{center}
\begin{tabular}{|l|l|l|}
\hline
\rowcolor{popblueL} Verbe & Nom & Personne \\
\hline
lesen & das Lesen & der \dots\dots \\
\hline
schreiben & der Artikel & der \dots\dots \\
\hline
berichten & der Bericht & der \dots\dots \\
\hline
senden & die \dots\dots & der Zuschauer \\
\hline
\end{tabular}
\end{center}
\end{exercice}

\begin{exercice}[Conjugaison : le prétérit des verbes forts]
Complète avec le prétérit du verbe entre parenthèses.
\begin{enumerate}
\item Der Reporter \dots\dots\ den Artikel. (schreiben)
\item Wir \dots\dots\ gestern die Nachrichten. (sehen)
\item Mein Vater \dots\dots\ am Morgen die Zeitung. (lesen)
\item Die Journalistin \dots\dots\ nach Deutschland. (kommen)
\item Er \dots\dots\ mit dem Redakteur. (sprechen)
\item Die Zeitung \dots\dots\ viele Informationen. (geben)
\end{enumerate}
\end{exercice}

\courssub{Je m'entraîne}

\begin{exercice}[Thème : traduction en allemand]
Traduis les phrases suivantes en allemand.
\begin{enumerate}
\item Hier, le reporter a écrit un article.
\item Nous avons regardé les informations à la télévision.
\item Le matin, mon père lisait le journal.
\item La journaliste est venue à Douala et a parlé avec les habitants.
\item Le journal donnait chaque jour des nouvelles du pays.
\end{enumerate}
\end{exercice}

\begin{exercice}[Ordre des mots : phrases mêlées]
Remets les mots dans l'ordre pour former une phrase correcte, puis transforme chaque phrase en commençant par le complément de temps.
\begin{enumerate}
\item las / gestern / der Reporter / den Artikel
\item die Nachrichten / wir / am Abend / sahen
\item kam / nach Yaoundé / letzte Woche / die Journalistin
\item mein Bruder / das Fernsehen / jeden Tag / sah
\end{enumerate}
Exemple de transformation : \all{Der Reporter las gestern den Artikel.} $\rightarrow$ \all{Gestern las der Reporter den Artikel.}
\end{exercice}

\begin{exercice}[Texte à trous : les médias]
Complète le texte avec les mots suivants (chaque mot n'est utilisé qu'une seule fois) : \all{Schlagzeilen -- Zeitung -- Nachrichten -- Reporter -- Sendung -- Artikel}.

Mein Vater kauft jeden Morgen die \dots\dots\ am Kiosk. Er liest zuerst die \dots\dots\ und dann die Sportseite. Mein Bruder ist \dots\dots\ bei einer Zeitung in Douala : er schreibt jeden Tag einen \dots\dots. Am Abend sehen wir zusammen die \dots\dots\ im Fernsehen. Danach sehen wir eine \dots\dots\ über Fußball.
\end{exercice}

\begin{exercice}[Appariements : mots et définitions]
Relie chaque mot allemand à sa définition française.
\begin{center}
\begin{tabular}{|l|l|}
\hline
\rowcolor{popblueL} Deutsch & Français \\
\hline
1. die Zeitung & a) celui qui regarde la télévision \\
\hline
2. der Reporter & b) l'émission \\
\hline
3. die Schlagzeile & c) le journal \\
\hline
4. die Sendung & d) le titre en gros caractères \\
\hline
5. der Zuschauer & e) le journaliste qui enquête \\
\hline
\end{tabular}
\end{center}
\end{exercice}

\courssub{Je me prépare au Bac}

\begin{textebox}[Ein Kameruner in Hamburg]
Jean Mbarga kommt aus Yaoundé. Vor zwei Jahren kam er nach Hamburg. Dort arbeitete er bei einer großen Zeitung. Jeden Morgen las er die Schlagzeilen und schrieb dann einen Artikel über Kamerun. Er sprach oft mit deutschen Studenten und gab ihnen Informationen über sein Land. Am Abend sah er mit seinen Kollegen die Nachrichten im Fernsehen. „Die Medien sind sehr wichtig“, sagte er. „Sie bringen die Menschen zusammen.“ Nach einem Jahr fuhr Jean nach Kamerun zurück. Heute arbeitet er bei einer Zeitung in Douala und schreibt über Deutschland.
\end{textebox}

Glossaire : \all{vor zwei Jahren} = il y a deux ans ; \all{der Kollege, -n} = le collègue ; \all{zusammenbringen} = rassembler ; \all{zurückfahren} = retourner.

\begin{exercice}[Compréhension écrite : un journaliste camerounais en Allemagne]
Lis le texte \all{Ein Kameruner in Hamburg}, puis réponds aux questions.

\textbf{A. Fragen zum Text} — Réponds en allemand par des phrases complètes.
\begin{enumerate}
\item Woher kommt Jean Mbarga ?
\item Wann kam er nach Hamburg ?
\item Was las er jeden Morgen ?
\item Mit wem sprach er oft ?
\item Wo arbeitet Jean heute ?
\end{enumerate}

\textbf{B. Richtig oder falsch ?} — Réponds par \all{richtig} ou \all{falsch} et justifie avec une phrase du texte.
\begin{enumerate}
\item Jean arbeitete bei einem Fernsehsender in Hamburg.
\item Er gab den Studenten Informationen über Kamerun.
\item Jean findet die Medien nicht wichtig.
\end{enumerate}

\textbf{C. Fragen de langue.}
\begin{enumerate}
\item Relève dans le texte cinq verbes au prétérit et donne leur infinitif.
\item Trouve dans le texte deux mots de la famille de \all{die Zeitung} et traduis-les.
\end{enumerate}
\end{exercice}

\begin{exercice}[Thème et version]
\textbf{A. Thème.} Traduis en allemand les phrases suivantes.
\begin{enumerate}
\item Hier, le reporter a écrit un article sur le football.
\item Nous avons regardé les informations à la télévision hier soir.
\item Le matin, mon père lisait le journal et buvait son café.
\item La journaliste est venue au Cameroun et a vu les marchés de Yaoundé.
\item Les journaux donnaient chaque semaine des nouvelles du monde.
\end{enumerate}
\end{exercice}

\begin{textebox}[Eine Familie in der Schweiz]
Familie Berger wohnt in Bern. Früher las der Vater jeden Morgen die Zeitung am Frühstückstisch. Die Mutter hörte im Auto die Nachrichten im Radio. Die Kinder sahen am Abend eine Sendung im Fernsehen. Am Wochenende kam die ganze Familie zusammen und sprach über die Schlagzeilen der Woche. Heute lesen alle die Nachrichten auf dem Handy, aber am Sonntag kaufen sie immer noch eine Zeitung.
\end{textebox}

Glossaire : \all{früher} = autrefois ; \all{der Frühstückstisch, -e} = la table du petit-déjeuner ; \all{das Handy, -s} = le téléphone portable ; \all{immer noch} = encore, toujours.

\begin{exercice}[Version]
Traduis en français le texte \all{Eine Familie in der Schweiz}.
\end{exercice}

\begin{exercice}[Production écrite guidée]
\textbf{A. Zusammenfassung.} Résume en cinq phrases, au prétérit, le texte \all{Ein Kameruner in Hamburg} de l'exercice de compréhension écrite.

\textbf{B. Meine Medien.} Rédige un paragraphe d'environ huit lignes (60 à 80 mots) sur tes habitudes avec les médias : quels journaux lis-tu ? Quelles émissions regardes-tu ? Quand et avec qui ? Utilise au prétérit au moins trois verbes forts de la leçon (\all{lesen, sehen, schreiben, kommen, geben, sprechen}) et les mots-charnières suivants dans l'ordre : \all{zuerst -- dann -- danach -- schließlich}.
\end{exercice}



\newpage

\coursec{Corrigés des exercices}

\begin{corrige}[Exercice 1 — QCM : les médias]
\begin{enumerate}
\item \textbf{b)} — \all{die Zeitung, -en} = le journal. (\all{das Fernsehen} = la télévision ; \all{das Radio} = la radio.)
\item \textbf{a)} — \all{der Reporter, -} = le journaliste qui écrit et enquête sur le terrain. (Pluriel identique : \all{die Reporter}.)
\item \textbf{a)} — \all{die Schlagzeile, -n} = le titre en gros caractères, à la une.
\item \textbf{b)} — \all{die Nachrichten} = les informations ; dans ce sens, le mot ne s'emploie qu'au pluriel : \all{Ich sehe die Nachrichten.} « Je regarde les informations. »
\item \textbf{a)} — \all{die Zeitung, die Zeitungen} : la plupart des noms féminins forment leur pluriel en \all{-(e)n} : \all{die Sendung, die Sendungen} ; \all{die Schlagzeile, die Schlagzeilen}.
\end{enumerate}
\textbf{Point de vigilance :} ne jamais écrire \all{die Zeitunger}. Attention au sens : \all{die Nachricht, -en} au singulier signifie « la nouvelle, le message » (\all{Ich bekomme eine Nachricht.} « Je reçois un message. »), mais « les informations » (presse, radio, télévision) se dit toujours \all{die Nachrichten} au pluriel.
\end{corrige}

\begin{corrige}[Exercice 2 — Vrai ou faux ?]
\begin{enumerate}
\item \all{Falsch. Das Fernsehen bedeutet „la télévision“, die Zeitung bedeutet „le journal“.}
\item \all{Richtig. Eine Sendung sieht man im Fernsehen oder hört man im Radio.} (Une émission se regarde à la télévision ou s'écoute à la radio.)
\item \all{Richtig. Der Reporter schreibt Artikel für die Zeitung.}
\item \all{Falsch. Die Nachrichten kommen im Radio, im Fernsehen und in der Zeitung.}
\end{enumerate}
\textbf{Point de vigilance :} dans la justification, le verbe conjugué reste à la deuxième place ; attention au genre : \all{das Fernsehen} (nom issu d'un infinitif, donc neutre, comme tous les noms formés sur un infinitif : \all{das Lesen, das Schreiben}).
\end{corrige}

\begin{corrige}[Exercice 3 — Vocabulaire : les familles de mots]
\begin{center}
\begin{tabular}{|l|l|l|}
\hline
\rowcolor{popblueL} Verbe & Nom & Personne \\
\hline
lesen & das Lesen & der \textbf{Leser} \\
\hline
schreiben & das Schreiben / der Artikel & der \textbf{Schriftsteller} / der \textbf{Journalist} \\
\hline
berichten & der Bericht & der \textbf{Reporter} \\
\hline
senden & die \textbf{Sendung} & der \textbf{Sender} \\
\hline
\end{tabular}
\end{center}
\textbf{Justification :} le suffixe \all{-er} forme le nom de la personne (ou de l'appareil) qui fait l'action : \all{lesen} $\rightarrow$ \all{der Leser} (le lecteur), \all{berichten} $\rightarrow$ \all{der Reporter} (le reporter), \all{senden} $\rightarrow$ \all{der Sender} (la chaîne, la station émettrice). Le suffixe \all{-ung} forme un nom féminin d'action : \all{senden} $\rightarrow$ \all{die Sendung} (l'émission). L'infinitif substantivé est toujours neutre : \all{das Lesen, das Schreiben}. \textbf{Point de vigilance :} \all{der Zuschauer} = le téléspectateur (celui qui regarde), à ne pas confondre avec \all{der Zuhörer} = l'auditeur (celui qui écoute la radio) ni avec \all{der Leser} = le lecteur (celui qui lit).
\end{corrige}

\begin{corrige}[Exercice 4 — Conjugaison : le prétérit des verbes forts]
\begin{enumerate}
\item Der Reporter \all{schrieb} den Artikel. (\all{schreiben -- schrieb -- hat geschrieben})
\item Wir \all{sahen} gestern die Nachrichten. (\all{sehen -- sah -- hat gesehen})
\item Mein Vater \all{las} am Morgen die Zeitung. (\all{lesen -- las -- hat gelesen})
\item Die Journalistin \all{kam} nach Deutschland. (\all{kommen -- kam -- ist gekommen})
\item Er \all{sprach} mit dem Redakteur. (\all{sprechen -- sprach -- hat gesprochen})
\item Die Zeitung \all{gab} viele Informationen. (\all{geben -- gab -- hat gegeben})
\end{enumerate}
\textbf{Justification grammaticale :} au prétérit, les verbes forts changent de voyelle radicale (\all{ei} $\rightarrow$ \all{ie} pour \all{schreiben} ; \all{e} $\rightarrow$ \all{a} pour \all{sehen, lesen, geben, sprechen}) et ne prennent \textbf{pas} de \all{-t-} de conjugaison : \all{ich las, du lasest, er las, wir lasen, ihr last, sie lasen}. La 1\iere{} et la 3\ieme{} personne du singulier sont identiques, sans terminaison : \all{ich kam / er kam}. \textbf{Point de vigilance :} au parfait, \all{kommen} se conjugue avec l'auxiliaire \all{sein} (verbe de déplacement) : \all{sie ist gekommen} ; les autres verbes de la liste prennent \all{haben}.
\end{corrige}

\begin{corrige}[Exercice 5 — Thème : traduction en allemand]
\begin{enumerate}
\item \all{Gestern schrieb der Reporter einen Artikel.} (ou au parfait : \all{Gestern hat der Reporter einen Artikel geschrieben.})
\item \all{Wir sahen die Nachrichten im Fernsehen.} (ou \all{Wir haben die Nachrichten im Fernsehen gesehen.})
\item \all{Am Morgen las mein Vater die Zeitung.}
\item \all{Die Journalistin kam nach Douala und sprach mit den Einwohnern.}
\item \all{Die Zeitung gab jeden Tag Nachrichten aus dem Land.}
\end{enumerate}
\textbf{Points de vigilance :}
\begin{itemize}
\item \textbf{Inversion :} quand la phrase commence par un complément de temps (\all{gestern}, \all{am Morgen}), le verbe conjugué reste en deuxième position et le sujet passe après : \all{Gestern schrieb der Reporter\dots} et jamais \all{Gestern der Reporter schrieb\dots}
\item \textbf{Article défini :} « le journal » = \all{die Zeitung} (féminin) ; « les informations » = \all{die Nachrichten} (pluriel obligatoire).
\item \textbf{Préposition :} « à la télévision » se dit \all{im Fernsehen} ; « à Douala » (ville) = \all{nach Douala}, sans article.
\item \textbf{Cas :} « un article » après le verbe = accusatif masculin : \all{einen Artikel} ; « avec les habitants » = \all{mit} + datif pluriel : \all{mit den Einwohnern}.
\end{itemize}
\end{corrige}

\begin{corrige}[Exercice 6 — Ordre des mots : phrases mêlées]
\begin{enumerate}
\item \all{Der Reporter las gestern den Artikel.} $\rightarrow$ \all{Gestern las der Reporter den Artikel.}
\item \all{Wir sahen am Abend die Nachrichten.} $\rightarrow$ \all{Am Abend sahen wir die Nachrichten.}
\item \all{Die Journalistin kam letzte Woche nach Yaoundé.} $\rightarrow$ \all{Letzte Woche kam die Journalistin nach Yaoundé.}
\item \all{Mein Bruder sah jeden Tag das Fernsehen.} $\rightarrow$ \all{Jeden Tag sah mein Bruder das Fernsehen.}
\end{enumerate}
\textbf{Règle appliquée :} dans la phrase déclarative allemande, le verbe conjugué occupe \textbf{toujours la deuxième place}. Si un complément (ici de temps) ouvre la phrase, il compte pour la première place : le sujet suit immédiatement le verbe (inversion sujet--verbe). \textbf{Point de vigilance :} \all{den Artikel} est à l'accusatif masculin (\all{der} $\rightarrow$ \all{den}) ; ne pas oublier cette marque dans la phrase remise en ordre. Remarque : pour « regarder la télévision », l'allemand emploie souvent le verbe à particule \all{fernsehen} : \all{Mein Bruder sah jeden Tag fern.} — la particule \all{fern} se détache et va en fin de phrase.
\end{corrige}

\begin{corrige}[Exercice 7 — Texte à trous : les médias]
Mein Vater kauft jeden Morgen die \all{Zeitung} am Kiosk. Er liest zuerst die \all{Schlagzeilen} und dann die Sportseite. Mein Bruder ist \all{Reporter} bei einer Zeitung in Douala : er schreibt jeden Tag einen \all{Artikel}. Am Abend sehen wir zusammen die \all{Nachrichten} im Fernsehen. Danach sehen wir eine \all{Sendung} über Fußball.

\textbf{Justification :} \all{die Zeitung} suit \all{kauft} (accusatif féminin, article inchangé) ; \all{die Schlagzeilen} est au pluriel car on lit plusieurs titres ; après le verbe \all{sein}, la profession s'emploie \textbf{sans article} : \all{Mein Bruder ist Reporter} ; \all{einen Artikel} : accusatif masculin après \all{schreibt} ; \all{eine Sendung} : accusatif féminin après \all{sehen}. \textbf{Point de vigilance :} \all{bei einer Zeitung} : \all{bei} + datif pour le lieu de travail (« dans un journal »), à ne pas confondre avec \all{nach Douala} (ville, sans article).
\end{corrige}

\begin{corrige}[Exercice 8 — Appariements : mots et définitions]
1 -- c (\all{die Zeitung} = le journal) ; 2 -- e (\all{der Reporter} = le journaliste qui enquête) ; 3 -- d (\all{die Schlagzeile} = le titre en gros caractères) ; 4 -- b (\all{die Sendung} = l'émission) ; 5 -- a (\all{der Zuschauer} = celui qui regarde la télévision, le téléspectateur).

\textbf{Point de vigilance :} \all{der Zuschauer} vient de \all{zusehen} (regarder) ; ne pas le confondre avec \all{der Leser} (le lecteur) ni avec \all{der Zuhörer} (l'auditeur).
\end{corrige}

\begin{corrige}[Exercice 9 — Compréhension écrite : un journaliste camerounais en Allemagne]
\textbf{Barème indicatif (sur 10) :} A : 5 $\times$ 1 pt = 5 pts ; B : 3 $\times$ 1 pt = 3 pts (0,5 pour \all{richtig}/\all{falsch} + 0,5 pour la justification) ; C : 2 pts (1 pt pour les cinq verbes, 1 pt pour les deux mots traduits).

\textbf{A. Fragen zum Text}
\begin{enumerate}
\item \all{Jean Mbarga kommt aus Yaoundé (aus Kamerun).}
\item \all{Vor zwei Jahren kam er nach Hamburg.}
\item \all{Jeden Morgen las er die Schlagzeilen.}
\item \all{Er sprach oft mit deutschen Studenten.}
\item \all{Heute arbeitet er bei einer Zeitung in Douala.}
\end{enumerate}

\textbf{B. Richtig oder falsch ?}
\begin{enumerate}
\item \all{Falsch.} Justification : \all{„Dort arbeitete er bei einer großen Zeitung.“} (Il travaillait dans un journal, pas à la télévision.)
\item \all{Richtig.} Justification : \all{„Er sprach oft mit deutschen Studenten und gab ihnen Informationen über sein Land.“}
\item \all{Falsch.} Justification : \all{„Die Medien sind sehr wichtig“, sagte er.}
\end{enumerate}

\textbf{C. Fragen de langue}
\begin{enumerate}
\item Cinq verbes au prétérit (au choix) : \all{kam} (kommen), \all{arbeitete} (arbeiten), \all{las} (lesen), \all{schrieb} (schreiben), \all{sprach} (sprechen) ; aussi possibles : \all{gab} (geben), \all{sah} (sehen), \all{sagte} (sagen), \all{fuhr} (fahren).
\item \all{die Schlagzeile} = le titre en gros caractères ; \all{der Artikel} = l'article (de journal).
\end{enumerate}
\textbf{Points de vigilance :} répondre par des \textbf{phrases complètes} avec le verbe en deuxième position ; reprendre le temps de la question (prétérit) ; \all{nach Hamburg} avec \all{nach} devant une ville, mais \all{bei einer Zeitung} avec \all{bei} + datif pour le lieu de travail. Remarque : \all{arbeitete} et \all{sagte} sont des verbes \textbf{faibles} (prétérit régulier en \all{-te}), sans changement de voyelle.
\end{corrige}

\begin{corrige}[Exercice 10 — Thème et version]
\textbf{Barème indicatif (sur 10) :} thème : 5 $\times$ 1,5 pt = 7,5 pts ; version : 2,5 pts (fidélité 1,5 pt, qualité du français 1 pt).

\textbf{A. Thème — corrigé.}
\begin{enumerate}
\item \all{Gestern schrieb der Reporter einen Artikel über Fußball.}
\item \all{Wir sahen gestern Abend die Nachrichten im Fernsehen.}
\item \all{Am Morgen las mein Vater die Zeitung und trank seinen Kaffee.} (\all{trinken -- trank -- hat getrunken})
\item \all{Die Journalistin kam nach Kamerun und sah die Märkte von Yaoundé.}
\item \all{Die Zeitungen gaben jede Woche Nachrichten aus der Welt.}
\end{enumerate}
\textbf{Points de vigilance (thème) :} « hier » (adverbe de temps) = \all{gestern} ; attention au faux ami : \all{hier} en allemand signifie « ici » ! « hier soir » = \all{gestern Abend} ; « au Cameroun » = \all{nach Kamerun} (pays neutre sans article) ; « les marchés » = \all{die Märkte} (pluriel avec Umlaut : \all{der Markt, die Märkte}) ; inversion après \all{gestern} et \all{am Morgen} ; « chaque semaine » = \all{jede Woche} (accusatif féminin après l'adjectif indéfini \all{jede}).

\textbf{B. Version — traduction de référence.}
La famille Berger habite à Berne. Autrefois, le père lisait le journal chaque matin à la table du petit-déjeuner. La mère écoutait les informations à la radio dans la voiture. Les enfants regardaient une émission à la télévision le soir. Le week-end, toute la famille se réunissait et parlait des grands titres de la semaine. Aujourd'hui, tous lisent les informations sur leur téléphone portable, mais le dimanche, ils achètent encore toujours un journal.

\textbf{Points de vigilance (version) :} \all{früher} = « autrefois » (marque de l'imparfait en français) ; \all{die ganze Familie} = « toute la famille » ; \all{immer noch} = « encore toujours » ; \all{das Handy} = « le téléphone portable » (mot allemand, faux ami de l'anglais) ; rendre les prétérits allemands exprimant une habitude passée par des imparfaits français et non par des passés composés.
\end{corrige}

\begin{corrige}[Exercice 11 — Production écrite guidée]
\textbf{Barème indicatif (sur 10) :} A : 4 pts (cinq phrases au prétérit, fidélité au texte) ; B : 6 pts (contenu 2 pts, prétérit des verbes forts 2 pts, mots-charnières dans l'ordre 1 pt, ordre des mots et orthographe 1 pt).

\textbf{A. Proposition de résumé (Zusammenfassung).}
\all{Jean Mbarga kam vor zwei Jahren nach Hamburg. Dort arbeitete er bei einer großen Zeitung. Jeden Morgen las er die Schlagzeilen und schrieb einen Artikel über Kamerun. Er sprach mit deutschen Studenten und gab ihnen Informationen über sein Land. Nach einem Jahr fuhr er nach Kamerun zurück.}

\textbf{B. Production modèle (Meine Medien).}
\all{Gestern Abend sah ich mit meiner Familie die Nachrichten im Fernsehen. Zuerst las mein Vater die Zeitung und die Schlagzeilen. Dann sahen wir zusammen eine Sendung über Fußball. Danach sprachen wir über die Nachrichten aus Kamerun. Schließlich schrieb ich einen kurzen Text für die Schule. Die Medien sind wichtig für mich : sie geben mir jeden Tag Informationen über mein Land und die Welt.}

\textbf{Points de vigilance :}
\begin{itemize}
\item Employer au moins trois verbes forts au prétérit correctement conjugués : \all{sah, las, sprach, schrieb, gab, kam}.
\item Respecter l'ordre des mots-charnières : \all{zuerst} $\rightarrow$ \all{dann} $\rightarrow$ \all{danach} $\rightarrow$ \all{schließlich} ; chacun occupe la première place et provoque l'inversion : \all{Zuerst las mein Vater\dots}
\item Majuscule à tous les noms : \all{die Zeitung, die Nachrichten, die Sendung, die Schule}.
\item \all{zurückfahren} est un verbe à particule : \all{er fuhr zurück} (particule en fin de phrase).
\item Vérifier la longueur : 60 à 80 mots environ, et relire pour corriger les accords, les cas et les articles.
\end{itemize}
\end{corrige}

\newpage

\coursec{Fiche bilan}

\begin{popfigure}
\begin{center}
\begin{tikzpicture}[mindmap, grow cyclic,
  every node/.style={concept, font=\small},
  root concept/.append style={concept color=popblue, font=\bfseries\small, minimum size=3.2cm},
  level 1/.append style={level distance=3.4cm, sibling angle=60},
  level 1 concept/.append style={font=\scriptsize, minimum size=2.1cm}]
\node[root concept, align=center]{Die Medien\\ Presse und Fernsehen}
  child[concept color=poporange]{ node[align=center]{Lexique\\ Zeitung, Reporter,\\ Schlagzeile, Sendung} }
  child[concept color=popgreen]{ node[align=center]{Präteritum\\ der starken Verben\\ kam, las, sah, ging} }
  child[concept color=poppurple]{ node[align=center]{Wortstellung\\ Verb an 2. Stelle\\ Gestern las er...} }
  child[concept color=popgold]{ node[align=center]{Textverständnis\\ global, dann detailliert} }
  child[concept color=poppink]{ node[align=center]{Produktion\\ Zusammenfassung\\ kurzer Text} }
  child[concept color=popblue!60]{ node[align=center]{Kultur\\ Medien in Deutschland\\ und Kamerun} };
\end{tikzpicture}
\end{center}
\legende{Carte mentale de la leçon : la presse écrite et le journal télévisé}
\end{popfigure}

\begin{aretenir}[L'essentiel de la leçon : le lexique]
\textbf{1. Lexique clé (à connaître avec article et pluriel)}
\begin{itemize}
  \item \all{die Zeitung, -en} : le journal ; \all{die Zeitschrift, -en} : la revue, le magazine ; \all{die Schlagzeile, -n} : le gros titre
  \item \all{der Artikel, -} : l'article ; \all{der Reporter, - / die Reporterin, -nen} : le/la journaliste reporter
  \item \all{die Nachricht, -en} : une nouvelle, un message ; \all{die Nachrichten (Pl.)} : les informations, le journal télévisé
  \item \all{das Fernsehen} (sans pluriel) : la télévision ; \all{die Sendung, -en} : l'émission ; \all{die Tagesschau} : le journal télévisé de la première chaîne allemande
  \item \all{erscheinen (erscheint, erschien, ist erschienen)} : paraître ; \all{berichten über + Akkusativ} : rendre compte de, informer sur
  \item \all{täglich / wöchentlich} : quotidien / hebdomadaire ; \all{die Auflage, -n} : le tirage
\end{itemize}
\textbf{2. Familles de mots}
\begin{itemize}
  \item \all{die Zeit} (le temps) $\rightarrow$ \all{die Zeitung} (le journal) $\rightarrow$ \all{die Zeitschrift} (la revue)
  \item \all{sehen} (voir) $\rightarrow$ \all{das Fernsehen} (la télévision) $\rightarrow$ \all{der Fernseher, -} (le poste de télévision)
  \item \all{senden} (émettre, envoyer) $\rightarrow$ \all{die Sendung} (l'émission) $\rightarrow$ \all{der Sender, -} (la chaîne, la station)
  \item \all{berichten} (informer) $\rightarrow$ \all{der Bericht, -e} (le reportage, le compte rendu)
\end{itemize}
\end{aretenir}

\begin{aretenir}[L'essentiel de la leçon : la grammaire]
\begin{center}
\begin{tabularx}{\linewidth}{|l|X|X|}
\hline
\rowcolor{popblueL} Point & Règle & Exemple \\
\hline
Prétérit des verbes forts & Radical modifié ; aucune terminaison à la 1\textsuperscript{re} et 3\textsuperscript{e} pers. sg. ; \emph{-st, -en, -t, -en} aux autres personnes & \all{kommen : ich kam, du kamst, er kam, wir kamen, ihr kamt, sie kamen} \\
\hline
Verbes forts fréquents du texte & \all{gehen -- ging ; lesen -- las ; sehen -- sah ; schreiben -- schrieb ; kommen -- kam ; bleiben -- blieb} & \all{Der Reporter sah das Spiel und schrieb einen Artikel.} \\
\hline
Auxiliaires au prétérit & \all{sein -- war ; haben -- hatte ; werden -- wurde} & \all{Die Sendung war interessant.} \\
\hline
Ordre des mots & Verbe conjugué toujours en 2\textsuperscript{e} position ; si un complément ouvre la phrase, le sujet passe après le verbe (inversion) & \all{Gestern las er die Zeitung.} (et non *\all{Gestern er las...}) \\
\hline
Phrase au Perfekt & Sujet -- auxiliaire en 2\textsuperscript{e} position -- compléments -- participe passé en fin de phrase & \all{Sie hat gestern die Nachrichten gesehen.} \\
\hline
\end{tabularx}
\end{center}
\end{aretenir}

\begin{attention}[Pièges fréquents des francophones]
\begin{itemize}
  \item Ne pas ajouter de terminaison à la 3\textsuperscript{e} personne du singulier au prétérit fort : \all{er kam}, \all{sie las} (et non *\all{er kamte}, *\all{sie laste}).
  \item \all{die Nachrichten} ne signifie pas « les nouvelles » au sens de messages personnels, mais « les informations » (journal télévisé) ; une seule nouvelle = \all{eine Nachricht}.
  \item Faux ami : \all{die Zeitung} = le journal (quotidien), alors que \all{die Zeitschrift} = la revue, le magazine.
  \item Après un complément en tête de phrase, inversion obligatoire : \all{Am Abend sah die Familie die Tagesschau.} « Le soir, la famille regardait le journal télévisé. »
  \item Majuscule à tous les noms : \all{das Fernsehen, die Sendung, der Reporter}.
\end{itemize}
\end{attention}

\begin{methode}[Méthodes express]
\begin{itemize}
  \item \textbf{Comprendre un texte de presse :} 1) lire le titre et la \all{Schlagzeile} ; 2) lecture globale (qui ? quoi ? où ? quand ?) ; 3) lecture détaillée avec le glossaire ; 4) repérer les verbes au prétérit et leur infinitif.
  \item \textbf{Résumer :} reprendre les informations essentielles en 3 à 5 phrases, au prétérit ou au présent, sans recopier le texte ; commencer par exemple par \all{Der Text berichtet über...} « Le texte rend compte de... ».
  \item \textbf{Produire :} phrases simples, verbe en 2\textsuperscript{e} position, majuscule à tous les noms, vocabulaire du thème ; relire les Umlaute (ä, ö, ü) et le ß.
\end{itemize}
\end{methode}

\begin{aretenir}[Checklist avant le Bac]
\begin{itemize}
  \item $\square$ Je connais le lexique des médias avec article et pluriel (\all{die Zeitung, -en ; die Sendung, -en ; der Reporter, -}).
  \item $\square$ Je sais conjuguer au prétérit les verbes forts fréquents : \all{gehen -- ging, lesen -- las, sehen -- sah, kommen -- kam, schreiben -- schrieb}.
  \item $\square$ Je connais le prétérit de \all{sein} (\all{war}) et de \all{haben} (\all{hatte}).
  \item $\square$ Je n'ajoute pas de terminaison à la 1\textsuperscript{re} et 3\textsuperscript{e} personne du singulier au prétérit fort : \all{er kam} (et non *\all{er kamte}).
  \item $\square$ Je place toujours le verbe conjugué en 2\textsuperscript{e} position, même après un complément de temps : \all{Gestern sah ich die Sendung.}
  \item $\square$ Je mets une majuscule à tous les noms : \all{die Zeitung, das Fernsehen, die Nachrichten}.
  \item $\square$ Je distingue \all{die Nachricht} (une nouvelle) et \all{die Nachrichten} (les informations, le journal télévisé).
  \item $\square$ J'écris les guillemets allemands „ … “ dans les citations.
  \item $\square$ Je sais lire une \all{Schlagzeile} et dégager l'idée générale d'un article en une phrase.
  \item $\square$ Je sais rédiger un résumé de 3 à 5 phrases avec le vocabulaire du thème et des verbes correctement conjugués.
  \item $\square$ Je relis ma production : place du verbe, accords, orthographe des Umlaute (ä, ö, ü) et du ß.
  \item $\square$ Je peux présenter oralement un fait divers ou une émission en 4 à 5 phrases simples.
\end{itemize}
\end{aretenir}

\findecours

\end{document}
